%% OVERFORE-Subdocumento2-Anexos.tex — generado desde OVERFORE-Subdocumento2-Anexos.docx
%% Formato: clase overfore, modo anexo. El formato lo pone la clase; aquí sólo va el contenido.
\documentclass[anexo]{../../comun/overfore}
\numsubdoc{2}
\informe{Informe Preparatorio N° 2}
\fecha{12 de octubre de 2026}
\externos{S2}
\begin{document}

\anexo{Listado de Requerimientos}
Este anexo presenta el detalle de los 248 requerimientos que la sección 2.5 del Subdocumento 2 resume por problema: 126 funcionales, 96 no funcionales, 15 de transición y preparación y 11 del proyecto. Cada requerimiento se entrega en una ficha con su identificador, el problema que lo origina y su origen en los documentos del caso, agrupado según su tipo. Dentro de cada tipo, el orden es el del catálogo consolidado. Los problemas que origina cada requerimiento son el Problema 1, trazabilidad comercial y conciliación; el Problema 2, productividad de la operación; el Problema 3, cumplimiento regulatorio y gestión de personas; y los requerimientos transversales, que responden a condiciones comunes a los tres.

Cada identificador se compone del tipo de requerimiento y del ámbito al que pertenece. Los tipos son RF (funcional), RNF (no funcional), TRA (transición y preparación) y RP (del proyecto). Los ámbitos de los requerimientos funcionales son TRZ (trazabilidad), CON (conciliación), PLA (cumplimiento del plan), TER (registro en terreno), EQU (equipos y mantenimiento), ACR (acreditación y acceso), CUM (cumplimiento regulatorio y ambiental), ADM (administración), NOT (notificaciones), DAT (calidad de datos), ANA (indicadores y analítica), CNX (conectividad y operación sin enlace). Los no funcionales usan su categoría de calidad: DES (desempeño), SEG (seguridad), USA (usabilidad), TRV (requisitos transversales), DIS (disponibilidad y continuidad), CAP (capacidad y escalabilidad), INT (interoperabilidad), OPE (operabilidad), POR (portabilidad), CUM (cumplimiento) e INF (infraestructura). En los de transición, MIG corresponde a migración de datos. El módulo y la etapa de cada requerimiento se indican para su trazabilidad con el Subdocumento 3 y el Formulario T-12.

\seccion{Requerimientos funcionales}
Los requerimientos funcionales describen lo que hace la solución ante un hecho de la operación. Cada ficha indica el módulo, la etapa de implementación en que se entrega (92 requerimientos en la Etapa 1 y 34 en la Etapa 2), el actor que interviene y el resultado esperado.

\begin{tablasintitulo}{I{0.200}I{0.800}}{\cab{Campo} & \cab{RF-TRZ-001}}
Problema & Problema 1 \\
Módulo & Operación Mina \\
Etapa & Etapa 1 \\
Requerimiento & El sistema debe registrar cada ciclo de carguío (cada palada: la pala llena el balde y lo vacía sobre el camión) con el polígono de extracción vigente, el viaje de camión, la pala que cargó, la persona operadora y la hora de inicio y de término. Los datos se toman del sistema de despacho de flota, sin pedirle nada a la persona operadora. \\
Actor & Sistema de despacho de flota \\
Resultado esperado & Cada ciclo de carguío (cada palada) queda registrado con polígono, viaje, pala, persona operadora e inicio y término, sin ninguna acción de la persona operadora. \\
Origen en las bases & Bases técnicas del caso Minería, cap. 9.1; Bases técnicas del caso Minería, cap. 4.3; Bases técnicas del caso Minería, restricción no negociable 5; Bases técnicas del caso Minería, decisión pendiente 1; supuesto SUP-01; supuesto SUP-16 \\
\end{tablasintitulo}

\begin{tablasintitulo}{I{0.200}I{0.800}}{\cab{Campo} & \cab{RF-TRZ-002}}
Problema & Problema 1 \\
Módulo & Trazabilidad y Embarque \\
Etapa & Etapa 1 \\
Requerimiento & El sistema debe mostrar, para cualquier turno, lote de chancado (mineral que entra a la planta) o lote comercial (concentrado que se embarca), los viajes de camión que lo forman y el polígono de origen de cada viaje, desde la extracción hasta el embarque. \\
Actor & Área comercial \\
Resultado esperado & Cualquier turno, lote de chancado o lote comercial se descompone en los viajes que lo forman, y cada viaje conserva su polígono de origen. \\
Origen en las bases & Bases técnicas del caso Minería, cap. 9.1; Bases técnicas del caso Minería, criterio de aceptación 1; Bases técnicas del caso Minería, decisión pendiente 1; supuesto SUP-01 \\
\end{tablasintitulo}

\begin{tablasintitulo}{I{0.200}I{0.800}}{\cab{Campo} & \cab{RF-TRZ-003}}
Problema & Problema 1 \\
Módulo & Trazabilidad y Embarque \\
Etapa & Etapa 1 \\
Requerimiento & El sistema debe permitir reconstruir, con evidencia, el origen de cualquier lote embarcado: de qué polígonos de extracción viene, cuándo se extrajo, por qué stocks pasó, cuándo se procesó y qué análisis de laboratorio lo respaldan. Cada dato debe enlazar al registro original que lo respalda o, en los tres casos en que no existe un registro medido, al supuesto declarado que lo sustituye: la carga en el borde entre dos polígonos, asignada al polígono donde estaba la pala (SUP-11); la composición del material retomado de un stock intermedio, calculada con la regla de mezcla (SUP-03); y el material acopiado antes del arranque del sistema, marcado como de origen anterior al sistema (SUP-21). En pantalla se indica cuál dato es registro original y cuál es supuesto. \\
Actor & Área comercial \\
Resultado esperado & Para el lote consultado se reconstruye el origen con polígonos, fechas de extracción, stocks recorridos, fecha de procesamiento y análisis, cada dato con enlace a su registro primario o, en los tres casos declarados (borde de polígono, mezcla en stock y material anterior al sistema), al supuesto que lo sustituye, distinguido en pantalla. \\
Origen en las bases & Bases técnicas del caso Minería, cap. 9.1; Bases técnicas del caso Minería, criterio de aceptación 1; supuesto SUP-01; supuesto SUP-03; supuesto SUP-11; supuesto SUP-21 \\
\end{tablasintitulo}

\begin{tablasintitulo}{I{0.200}I{0.800}}{\cab{Campo} & \cab{RF-TRZ-004}}
Problema & Problema 1 \\
Módulo & Trazabilidad y Embarque \\
Etapa & Etapa 1 \\
Requerimiento & El sistema debe permitir hacer el camino inverso: a partir de un polígono extraído en una fecha dada, ver en qué lotes comerciales y en qué embarques terminó su material y con qué proporción participó en cada uno, con al menos diez años de historial consultable. \\
Actor & Planificación minera \\
Resultado esperado & A partir del polígono se listan los lotes comerciales y los embarques en que terminó su material, con la proporción de participación, para cualquier polígono extraído dentro de los últimos diez años. \\
Origen en las bases & Bases técnicas del caso Minería, cap. 9.1; Bases técnicas del caso Minería, criterio de aceptación 1; Derivado del requisito de retención de 10 años (RT-05.10 de las Bases técnicas del caso Minería, equivalente al RT-16.10 de las Bases técnicas transversales) \\
\end{tablasintitulo}

\begin{tablasintitulo}{I{0.200}I{0.800}}{\cab{Campo} & \cab{RF-TRZ-005}}
Problema & Problema 1 \\
Módulo & Trazabilidad y Embarque \\
Etapa & Etapa 1 \\
Requerimiento & El sistema debe indicar, en cada etapa de la cadena de origen de un lote, si el dato fue medido o es un supuesto (una estimación calculada con una regla), y en ese caso mostrar la regla aplicada y su versión. \\
Actor & Área comercial \\
Resultado esperado & Cada etapa se muestra marcada como medido o como supuesto, y en este último caso con la regla y la versión aplicadas. \\
Origen en las bases & Bases técnicas del caso Minería, cap. 9.1; Bases técnicas del caso Minería, criterio de aceptación 1; supuesto SUP-01; supuesto SUP-21 \\
\end{tablasintitulo}

\begin{tablasintitulo}{I{0.200}I{0.800}}{\cab{Campo} & \cab{RF-TRZ-006}}
Problema & Problema 1 \\
Módulo & Trazabilidad y Embarque \\
Etapa & Etapa 1 \\
Requerimiento & El sistema debe llevar el registro de cada stock intermedio: lo que entra (polígono de origen, tonelaje, ley de cada elemento, fecha y hora, equipo y viaje), lo que se retoma con cargador frontal (tonelaje, destino, hora y qué proporción vino de cada polígono) y los traspasos entre stocks y el reproceso del stock de baja ley. La proporción de cada polígono se calcula con el promedio ponderado por tonelaje de lo que hay vigente en el stock, se recalcula tras cada aporte, retoma y traspaso, y la regla es parametrizable por el Cliente. Cada movimiento conserva la cadena de origen por polígono. Reemplaza la planilla del supervisor. \\
Actor & Supervisor de turno \\
Resultado esperado & Todo aporte, retoma, traspaso y reproceso del stock queda registrado con su polígono de origen, y la planilla del supervisor deja de usarse. \\
Origen en las bases & Bases técnicas del caso Minería, cap. 4.4; Bases técnicas del caso Minería, cap. 5; Bases técnicas del caso Minería, criterio de aceptación 1; Bases técnicas del caso Minería, decisión pendiente 3; supuesto SUP-03; supuesto SUP-23 \\
\end{tablasintitulo}

\begin{tablasintitulo}{I{0.200}I{0.800}}{\cab{Campo} & \cab{RF-TRZ-007}}
Problema & Problema 1 \\
Módulo & Trazabilidad y Embarque \\
Etapa & Etapa 1 \\
Requerimiento & El sistema debe mantener actualizado, para cada stock intermedio y para el stock de baja ley, su saldo, su composición por polígono de origen y su ley resultante según la regla de mezcla declarada, como máximo 15 minutos después de cada ingreso o retoma. \\
Actor & Supervisor de turno \\
Resultado esperado & El saldo, la composición por polígono y la ley del stock se muestran actualizados en un máximo de 15 minutos después de cada aporte y cada retoma. \\
Origen en las bases & Bases técnicas del caso Minería, cap. 4.4; Bases técnicas del caso Minería, cap. 5; Bases técnicas del caso Minería, criterio de aceptación 1; Bases técnicas del caso Minería, decisión pendiente 3; requisito técnico RT-05.29; supuesto SUP-03; supuesto SUP-21 \\
\end{tablasintitulo}

\begin{tablasintitulo}{I{0.200}I{0.800}}{\cab{Campo} & \cab{RF-TRZ-008}}
Problema & Problema 1 \\
Módulo & Trazabilidad y Embarque \\
Etapa & Etapa 1 \\
Requerimiento & El sistema debe calcular la ley del material que se retoma de un stock con la regla de mezcla declarada (hoy se asume el promedio ponderado por tonelaje de lo acopiado) y mostrar junto al resultado la regla usada y cuánto aportó cada polígono. \\
Actor & Metalurgia y planta \\
Resultado esperado & La retoma queda con su ley calculada y se muestran, junto al resultado, la regla de mezcla aplicada y la proporción de cada polígono de origen. \\
Origen en las bases & Bases técnicas del caso Minería, cap. 4.4; Bases técnicas del caso Minería, decisión pendiente 3; supuesto SUP-03 \\
\end{tablasintitulo}

\begin{tablasintitulo}{I{0.200}I{0.800}}{\cab{Campo} & \cab{RF-TRZ-009}}
Problema & Problema 1 \\
Módulo & Trazabilidad y Embarque \\
Etapa & Etapa 1 \\
Requerimiento & El sistema debe registrar, en cada campaña de verificación de un stock (muestreo del material retomado para comprobar su ley), la ley que se había asumido y la ley medida, y publicar la diferencia por stock y por período, con un responsable asignado. \\
Actor & Metalurgia y planta \\
Resultado esperado & La desviación entre ley asumida y ley medida queda publicada por stock y período, con responsable asignado. \\
Origen en las bases & Bases técnicas del caso Minería, cap. 4.4; Bases técnicas del caso Minería, decisión pendiente 3; supuesto SUP-03 \\
\end{tablasintitulo}

\begin{tablasintitulo}{I{0.200}I{0.800}}{\cab{Campo} & \cab{RF-TRZ-010}}
Problema & Problema 1 \\
Módulo & Operación Mina \\
Etapa & Etapa 1 \\
Requerimiento & El sistema debe asignar cada viaje de camión a un polígono de origen. Mientras el Cliente no decida otra cosa, todo el tonelaje del viaje se asigna al polígono donde está ubicada la pala, aunque la pala esté en el borde entre dos polígonos (es la práctica actual). \\
Actor & Sistema de despacho de flota \\
Resultado esperado & Cada viaje queda asignado a un polígono de origen según la regla vigente de SUP-11. \\
Origen en las bases & Bases técnicas del caso Minería, cap. 4.3; Bases técnicas del caso Minería, criterio de aceptación 1; Bases técnicas del caso Minería, decisión pendiente 11; supuesto SUP-11 \\
\end{tablasintitulo}

\begin{tablasintitulo}{I{0.200}I{0.800}}{\cab{Campo} & \cab{RF-TRZ-011}}
Problema & Problema 1 \\
Módulo & Operación Mina \\
Etapa & Etapa 1 \\
Requerimiento & El sistema debe registrar, en cada ciclo de carguío (palada), dónde está la pala respecto del contacto que declara el plan (el límite entre dos polígonos o dos tipos de material), tomando el dato del sistema de despacho, sin pedir nada a la persona operadora. \\
Actor & Sistema de despacho de flota \\
Resultado esperado & Cada ciclo de carguío queda con la posición de la pala respecto del contacto, sin acción de la persona operadora. \\
Origen en las bases & Bases técnicas del caso Minería, cap. 4.3; supuesto SUP-11 \\
\end{tablasintitulo}

\begin{tablasintitulo}{I{0.200}I{0.800}}{\cab{Campo} & \cab{RF-TRZ-012}}
Problema & Problema 1 \\
Módulo & Trazabilidad y Embarque \\
Etapa & Etapa 1 \\
Requerimiento & El sistema debe preparar y entregar a SAP S/4HANA, para su emisión electrónica, los datos de la guía de despacho del concentrado a puerto, con identificador del lote, tonelaje húmedo y seco, humedad, ley de cada elemento (cobre, oro, plata y arsénico) y vínculo al origen del lote, y conservar el folio y el código QR de la guía emitida, reemplazando la guía y la planilla en papel. \\
Actor & Despachador de concentrado \\
Resultado esperado & La guía de despacho se emite en forma electrónica con los datos del lote y su vínculo al origen, sin guía ni planilla en papel. \\
Origen en las bases & Bases técnicas del caso Minería, cap. 4.6; supuesto SUP-23 \\
\end{tablasintitulo}

\begin{tablasintitulo}{I{0.200}I{0.800}}{\cab{Campo} & \cab{RF-TRZ-013}}
Problema & Problema 1 \\
Módulo & Trazabilidad y Embarque \\
Etapa & Etapa 1 \\
Requerimiento & El sistema debe registrar la recepción y el acopio del concentrado en el recinto portuario por lote, con tonelaje, humedad, fecha y responsable, aun cuando el recinto sea administrado por un tercero. \\
Actor & Despachador de concentrado \\
Resultado esperado & La recepción y el acopio quedan registrados por lote con tonelaje, humedad, fecha y responsable. \\
Origen en las bases & Bases técnicas del caso Minería, cap. 4.6; supuesto SUP-24 \\
\end{tablasintitulo}

\begin{tablasintitulo}{I{0.200}I{0.800}}{\cab{Campo} & \cab{RF-TRZ-014}}
Problema & Problema 1 \\
Módulo & Trazabilidad y Embarque \\
Etapa & Etapa 1 \\
Requerimiento & El sistema debe registrar, con marca de tiempo y responsable, cada evento del embarque en el terminal portuario (programación del embarque, conformación del lote comercial, toma de muestras, recepción de la nave, inicio y término de la carga, emisión de los documentos de embarque y zarpe), además de la recepción y el acopio ya cubiertos en RF-TRZ-013. \\
Actor & Área comercial \\
Resultado esperado & Cada evento del embarque, desde su programación hasta el zarpe, queda registrado con hora y responsable. \\
Origen en las bases & Bases técnicas del caso Minería, cap. 4.6; Bases técnicas del caso Minería, cap. 11; requisito técnico RT-12.12; supuesto SUP-24 \\
\end{tablasintitulo}

\begin{tablasintitulo}{I{0.200}I{0.800}}{\cab{Campo} & \cab{RF-TRZ-015}}
Problema & Problema 1 \\
Módulo & Trazabilidad y Embarque \\
Etapa & Etapa 1 \\
Requerimiento & El sistema debe conformar el lote comercial a partir de los lotes de concentrado que lo componen, conservando la trazabilidad al origen de cada uno y la proporción con que participa en la mezcla embarcada. \\
Actor & Área comercial \\
Resultado esperado & El lote comercial queda conformado con la proporción de cada lote de concentrado y su trazabilidad al origen. \\
Origen en las bases & Bases técnicas del caso Minería, cap. 4.6; Bases técnicas del caso Minería, criterio de aceptación 1 \\
\end{tablasintitulo}

\begin{tablasintitulo}{I{0.200}I{0.800}}{\cab{Campo} & \cab{RF-TRZ-016}}
Problema & Problema 1 \\
Módulo & Trazabilidad y Embarque \\
Etapa & Etapa 2 \\
Requerimiento & El sistema debe poner a disposición del comprador, en un portal con acceso autenticado, el certificado de análisis y el expediente de trazabilidad (el resumen del origen del lote) de cada lote, registrando quién lo consultó y cuándo, en reemplazo del envío por correo. \\
Actor & Comprador \\
Resultado esperado & El comprador consulta y descarga el certificado y el expediente del lote, y la consulta queda registrada. \\
Origen en las bases & Bases técnicas del caso Minería, cap. 4.6; requisito técnico RT-16.30 de las Bases técnicas del caso Minería (portal de trazabilidad de lote para compradores, autenticado); requisito técnico RT-12.12; supuesto SUP-22 \\
\end{tablasintitulo}

\begin{tablasintitulo}{I{0.200}I{0.800}}{\cab{Campo} & \cab{RF-TRZ-017}}
Problema & Problema 1 \\
Módulo & Trazabilidad y Embarque \\
Etapa & Etapa 1 \\
Requerimiento & El sistema debe registrar el acta de aceptación del embarque con firma electrónica avanzada y vincularla al lote comercial y a su expediente de trazabilidad. \\
Actor & Área comercial \\
Resultado esperado & El acta de aceptación queda firmada con firma electrónica avanzada y vinculada al lote comercial y a su expediente. \\
Origen en las bases & Bases técnicas del caso Minería, cap. 4.6; requisito técnico RT-16.14 de las Bases técnicas del caso Minería, equivalente al RT-16.17 de las Bases técnicas transversales; supuesto SUP-25 \\
\end{tablasintitulo}

\begin{tablasintitulo}{I{0.200}I{0.800}}{\cab{Campo} & \cab{RF-TRZ-018}}
Problema & Problema 1 \\
Módulo & Trazabilidad y Embarque \\
Etapa & Etapa 1 \\
Requerimiento & Cuando el laboratorio rechace o invalide una muestra, el sistema debe bloquear el lote para su venta y pedir una nueva toma de muestra dentro de 24 horas desde el rechazo; ningún lote sin análisis válido puede liberarse a embarque. Además debe guardar cada cambio de estado del lote en un registro que nadie pueda modificar: quién lo hizo, cuándo (con zona horaria) y los valores antes y después. \\
Actor & Laboratorio químico \\
Resultado esperado & El lote queda bloqueado para embarque, se genera el remuestreo con plazo de 24 horas, ningún lote sin análisis válido puede liberarse y cada cambio de estado queda en la auditoría con responsable, fecha y hora y valores anterior y posterior, sin posibilidad de modificarlo. \\
Origen en las bases & Bases técnicas del caso Minería, cap. 4.6; Bases técnicas del caso Minería, decisión pendiente 6; requisito técnico RT-16.06; requisito técnico RT-16.07; supuesto SUP-06 \\
\end{tablasintitulo}

\begin{tablasintitulo}{I{0.200}I{0.800}}{\cab{Campo} & \cab{RF-CON-001}}
Problema & Problema 1 \\
Módulo & Conciliación Metalúrgica \\
Etapa & Etapa 1 \\
Requerimiento & El sistema debe registrar y conservar las tres mediciones de tonelaje tal como se generan (pesaje de a bordo del camión, báscula de control de camiones y balanza de alimentación de la planta), sin transformarlas, sin descartarlas y sin intentar explicar sus diferencias. \\
Actor & Metalurgia y planta \\
Resultado esperado & Las tres mediciones quedan almacenadas tal como se generaron, sin transformación ni descarte. \\
Origen en las bases & Bases técnicas del caso Minería, cap. 9.2; Bases técnicas del caso Minería, cap. 4.5; Bases técnicas del caso Minería, decisión pendiente 2; supuesto SUP-02; supuesto SUP-29 \\
\end{tablasintitulo}

\begin{tablasintitulo}{I{0.200}I{0.800}}{\cab{Campo} & \cab{RF-CON-002}}
Problema & Problema 1 \\
Módulo & Conciliación Metalúrgica \\
Etapa & Etapa 1 \\
Requerimiento & El sistema debe usar en la conciliación metalúrgica una única fuente de verdad del tonelaje (la declarada como oficial: balanza de alimentación para la planta y pesaje de a bordo, corregido con la báscula, para la mina), mostrar en cada resultado cuál fuente se usó y no emitir cifras contables fuera de SAP. \\
Actor & Metalurgia y planta; Finanzas y costos \\
Resultado esperado & Cada conciliación usa una sola fuente de verdad del tonelaje y muestra cuál fue; la cifra contable la emite SAP y el sistema no emite cifra contable propia. \\
Origen en las bases & Bases técnicas del caso Minería, cap. 9.2; Bases técnicas del caso Minería, restricción no negociable 4; Bases técnicas del caso Minería, criterio de aceptación 3; Bases técnicas del caso Minería, decisión pendiente 2; supuesto SUP-02 \\
\end{tablasintitulo}

\begin{tablasintitulo}{I{0.200}I{0.800}}{\cab{Campo} & \cab{RF-CON-003}}
Problema & Problema 1 \\
Módulo & Conciliación Metalúrgica \\
Etapa & Etapa 1 \\
Requerimiento & El sistema debe calcular la conciliación metalúrgica del mes dentro del sistema, con acceso al dato de origen de cada componente, a partir de los registros primarios de tonelaje, ley y recuperación, sustituyendo las planillas personales del metalurgista. \\
Actor & Metalurgia y planta \\
Resultado esperado & La conciliación del mes se obtiene desde el sistema, sin planillas personales, con acceso al dato de origen de cada componente. \\
Origen en las bases & Bases técnicas del caso Minería, cap. 9.2; Bases técnicas del caso Minería, criterio de aceptación 2 \\
\end{tablasintitulo}

\begin{tablasintitulo}{I{0.200}I{0.800}}{\cab{Campo} & \cab{RF-CON-004}}
Problema & Problema 1 \\
Módulo & Conciliación Metalúrgica \\
Etapa & Etapa 1 \\
Requerimiento & El sistema debe descomponer la diferencia de conciliación por punto de la cadena -mina, stocks intermedios, planta y embarque- y por periodo, en lugar de entregarla como una cifra única sin explicación. \\
Actor & Metalurgia y planta \\
Resultado esperado & La diferencia de conciliación queda descompuesta por mina, stocks, planta y embarque, y por período. \\
Origen en las bases & Bases técnicas del caso Minería, cap. 9.2; Bases técnicas del caso Minería, criterio de aceptación 3; supuesto SUP-20 \\
\end{tablasintitulo}

\begin{tablasintitulo}{I{0.200}I{0.800}}{\cab{Campo} & \cab{RF-CON-005}}
Problema & Problema 1 \\
Módulo & Conciliación Metalúrgica \\
Etapa & Etapa 1 \\
Requerimiento & El sistema debe alertar cuando la diferencia de conciliación del período supere el umbral comprometido, indicando el punto de la cadena que más aporta a la diferencia y su responsable. \\
Actor & Metalurgia y planta \\
Resultado esperado & Se emite alerta con el punto de la cadena que más aporta a la diferencia y su responsable. \\
Origen en las bases & Bases técnicas del caso Minería, cap. 9.2; Bases técnicas del caso Minería, criterio de aceptación 3; Derivado de la sección 9.2 y del criterio de aceptación 3 (explicar la diferencia por punto de muestreo): la alerta y el responsable permiten actuar sobre esa diferencia \\
\end{tablasintitulo}

\begin{tablasintitulo}{I{0.200}I{0.800}}{\cab{Campo} & \cab{RF-DAT-001}}
Problema & Problema 1 \\
Módulo & Administración y Servicios Comunes \\
Etapa & Etapa 1 \\
Requerimiento & El sistema debe cuidar la calidad de los datos según la norma ISO/IEC 25012: validarlos al capturarlos y publicar en un tablero, por dominio de datos, indicadores de completitud, exactitud y consistencia, con metas de al menos 95 \%, 98 \% y 98 \% respectivamente, medidos cada mes. \\
Actor & Administrador funcional del Cliente \\
Resultado esperado & El dato inválido se rechaza en la captura con mensaje, y el tablero muestra la calidad de cada dominio frente a su umbral objetivo. \\
Origen en las bases & requisito técnico RT-05.04; Derivado de RT-05.04 (indicadores de calidad, ISO/IEC 25012): los umbrales y la periodicidad mensual se fijan para poder verificarlos \\
\end{tablasintitulo}

\begin{tablasintitulo}{I{0.200}I{0.800}}{\cab{Campo} & \cab{RF-CON-006}}
Problema & Problema 1 \\
Módulo & Conciliación Metalúrgica \\
Etapa & Etapa 1 \\
Requerimiento & El sistema debe publicar la diferencia entre las tres mediciones de tonelaje (pesaje de a bordo, báscula y balanza de planta) por período, por equipo y por punto de la cadena, en toneladas y en porcentaje. \\
Actor & Metalurgia y planta \\
Resultado esperado & La diferencia entre las tres mediciones queda publicada por período, equipo y punto de la cadena, en toneladas y en porcentaje. \\
Origen en las bases & Bases técnicas del caso Minería, cap. 4.3; Bases técnicas del caso Minería, cap. 9.2; Bases técnicas del caso Minería, cap. 4.5; Bases técnicas del caso Minería, criterio de aceptación 3; supuesto SUP-02 \\
\end{tablasintitulo}

\begin{tablasintitulo}{I{0.200}I{0.800}}{\cab{Campo} & \cab{RF-CON-007}}
Problema & Problema 1 \\
Módulo & Conciliación Metalúrgica \\
Etapa & Etapa 1 \\
Requerimiento & El sistema debe guardar cada factor de ajuste de la conciliación (corrección que se aplica al cálculo) como registro con versiones: fórmula, vigencia, responsable, fundamento y fecha de incorporación, incluidos los acumulados desde 2013. Debe poder reproducir exactamente la conciliación de cualquier período pasado con los factores que regían entonces, sin depender del conocimiento de una persona determinada. \\
Actor & Metalurgia y planta \\
Resultado esperado & Cada factor queda como entidad versionada con sus atributos, incluidos los acumulados desde 2013, y la conciliación de cualquier período pasado se reproduce exactamente con los factores vigentes en ese período. \\
Origen en las bases & Bases técnicas del caso Minería, cap. 4.5; Bases técnicas del caso Minería, cap. 9.2; Bases técnicas del caso Minería, criterio de aceptación 3; Bases técnicas del caso Minería, cap. 8 (entrevista de Sanhueza); Derivado de esa entrevista: el versionado de los factores de ajuste \\
\end{tablasintitulo}

\begin{tablasintitulo}{I{0.200}I{0.800}}{\cab{Campo} & \cab{RF-CON-008}}
Problema & Problema 1 \\
Módulo & Conciliación Metalúrgica \\
Etapa & Etapa 1 \\
Requerimiento & El sistema debe calcular el balance metalúrgico de cada turno (cuánto cobre entró, se recuperó y salió en el turno) con los mismos registros y las mismas reglas que el balance del mes, eliminando las dos planillas distintas. \\
Actor & Metalurgia y planta \\
Resultado esperado & El balance de turno se obtiene con los mismos registros y reglas que el balance del mes. \\
Origen en las bases & Bases técnicas del caso Minería, cap. 4.5 \\
\end{tablasintitulo}

\begin{tablasintitulo}{I{0.200}I{0.800}}{\cab{Campo} & \cab{RF-CON-009}}
Problema & Problema 1 \\
Módulo & Conciliación Metalúrgica \\
Etapa & Etapa 1 \\
Requerimiento & El sistema debe recibir los resultados de análisis directamente del sistema de gestión de laboratorio, conectado en ambos sentidos, sin que nadie los transcriba a mano en ningún punto. \\
Actor & Sistema de gestión de laboratorio \\
Resultado esperado & El resultado queda disponible en la conciliación sin transcripción manual. \\
Origen en las bases & Bases técnicas del caso Minería, cap. 4.5; Bases técnicas del caso Minería, restricción no negociable 6; supuesto SUP-20; supuesto SUP-16 \\
\end{tablasintitulo}

\begin{tablasintitulo}{I{0.200}I{0.800}}{\cab{Campo} & \cab{RF-CON-010}}
Problema & Problema 1 \\
Módulo & Conciliación Metalúrgica \\
Etapa & Etapa 1 \\
Requerimiento & El sistema debe mostrar junto a cada resultado de conciliación su memoria de cálculo: el detalle paso a paso con los datos de entrada, los factores aplicados y el resultado intermedio de cada paso. \\
Actor & Metalurgia y planta \\
Resultado esperado & Junto al resultado se muestra la memoria de cálculo con datos de entrada, factores y resultados intermedios. \\
Origen en las bases & Bases técnicas del caso Minería, cap. 4.5; Bases técnicas del caso Minería, cap. 9.2; Bases técnicas del caso Minería, criterio de aceptación 3 \\
\end{tablasintitulo}

\begin{tablasintitulo}{I{0.200}I{0.800}}{\cab{Campo} & \cab{RF-CON-011}}
Problema & Problema 1 \\
Módulo & Conciliación Metalúrgica \\
Etapa & Etapa 1 \\
Requerimiento & El sistema debe capturar la muestra en los puntos de transferencia (donde el material pasa de una etapa a otra) sin agregar pasos al ciclo de la pala ni del camión. El Adjudicatario debe declarar en su propuesta, para cada punto de muestreo, cuántos segundos agrega al ciclo y cuántas toneladas menos al año se moverían por ello. \\
Actor & Metalurgia y planta; Proponente (declaración en la propuesta) \\
Resultado esperado & Cada muestra queda registrada en su punto de transferencia sin agregar pasos al ciclo de pala ni de camión, y la propuesta declara por punto de muestreo el costo en segundos de ciclo y el efecto en toneladas por año. \\
Origen en las bases & Bases técnicas del caso Minería, cap. 9.4; Bases técnicas del caso Minería, restricción no negociable 8; Bases técnicas del caso Minería, decisión pendiente 12; supuesto SUP-12 \\
\end{tablasintitulo}

\begin{tablasintitulo}{I{0.200}I{0.800}}{\cab{Campo} & \cab{RF-CON-012}}
Problema & Problema 1 \\
Módulo & Conciliación Metalúrgica \\
Etapa & Etapa 1 \\
Requerimiento & El sistema debe publicar por polígono la diferencia entre la ley estimada del plan y del modelo de bloques y la ley medida en la alimentación de planta, tratando el resultado del laboratorio como fuente de verdad de la ley medida. \\
Actor & Planificación minera \\
Resultado esperado & La diferencia entre ley estimada y ley medida queda publicada por polígono. \\
Origen en las bases & Bases técnicas del caso Minería, cap. 9.4; Bases técnicas del caso Minería, cap. 4.1; Bases técnicas del caso Minería, cap. 4.5; Bases técnicas del caso Minería, restricción no negociable 6; supuesto SUP-20 \\
\end{tablasintitulo}

\begin{tablasintitulo}{I{0.200}I{0.800}}{\cab{Campo} & \cab{RF-CON-013}}
Problema & Problema 1 \\
Módulo & Indicadores y Centro Integrado \\
Etapa & Etapa 2 \\
Requerimiento & El sistema debe calcular el costo por tonelada movida y por libra de cobre fino todos los días, a partir de los datos contables que entrega SAP, y dejarlo disponible en los tableros de indicadores. \\
Actor & Finanzas y costos \\
Resultado esperado & El costo por tonelada movida y por libra de cobre fino del día anterior queda disponible en los tableros de indicadores, calculado con los datos contables que entrega SAP. \\
Origen en las bases & Bases técnicas del caso Minería, restricción no negociable 4; Bases técnicas del caso Minería, cap. 8 (entrevista de Ortúzar) \\
\end{tablasintitulo}

\begin{tablasintitulo}{I{0.200}I{0.800}}{\cab{Campo} & \cab{RF-CON-014}}
Problema & Problema 1 \\
Módulo & Indicadores y Centro Integrado \\
Etapa & Etapa 2 \\
Requerimiento & El sistema debe permitir ir desde el costo por tonelada (indicador resumen) hasta cada transacción de origen que lo compone. \\
Actor & Finanzas y costos \\
Resultado esperado & Desde el indicador de costo por tonelada se llega a las transacciones de origen que lo componen. \\
Origen en las bases & requisito técnico RT-05.26 \\
\end{tablasintitulo}

\begin{tablasintitulo}{I{0.200}I{0.800}}{\cab{Campo} & \cab{RF-CON-015}}
Problema & Problema 1 \\
Módulo & Conciliación Metalúrgica \\
Etapa & Etapa 1 \\
Requerimiento & El sistema debe entregar a SAP los datos operacionales que este necesita para su registro contable, en formato estructurado abierto (JSON, a través de la interfaz de programación de S/4HANA), todos los días al cierre del día operacional y con un consolidado mensual a más tardar el tercer día hábil del mes siguiente. SAP sigue emitiendo la contabilidad y el sistema no la duplica ni la reemplaza. \\
Actor & ERP SAP S/4HANA \\
Resultado esperado & SAP recibe los datos en formato estructurado abierto, diariamente y con consolidado mensual al tercer día hábil, y la emisión contable sigue en SAP. \\
Origen en las bases & Bases técnicas del caso Minería, cap. 9.2; Bases técnicas del caso Minería, cap. 5; Bases técnicas del caso Minería, restricción no negociable 4; supuesto SUP-16 \\
\end{tablasintitulo}

\begin{tablasintitulo}{I{0.200}I{0.800}}{\cab{Campo} & \cab{RF-TER-001}}
Problema & Problema 2 \\
Módulo & Operación Mina \\
Etapa & Etapa 1 \\
Requerimiento & El sistema debe permitir a la persona operadora de la pala avisar, con una sola pulsación, que el material que carga cambió respecto de lo planificado, sin elegir causa, sin escribir y sin pantallas adicionales. Cada aviso queda asociado automáticamente al viaje, al ciclo de carguío (palada), al polígono, al equipo, a la persona y a la hora. La cabina debe confirmar visualmente que el aviso quedó registrado. \\
Actor & Persona operadora de equipo \\
Resultado esperado & El aviso queda registrado con una sola acción y asociado a viaje, ciclo, polígono, equipo, persona y hora, sin que se pida ningún dato, y la cabina muestra la confirmación visible del registro. \\
Origen en las bases & Bases técnicas del caso Minería, cap. 9.4; Bases técnicas del caso Minería, cap. 4.3; Bases técnicas del caso Minería, restricción no negociable 8; Bases técnicas del caso Minería, criterio de aceptación 14; Bases técnicas del caso Minería, decisión pendiente 1; supuesto SUP-11 \\
\end{tablasintitulo}

\begin{tablasintitulo}{I{0.200}I{0.800}}{\cab{Campo} & \cab{RF-TER-002}}
Problema & Problema 2 \\
Módulo & Operación Mina \\
Etapa & Etapa 1 \\
Requerimiento & El sistema debe indicar en la cabina de la pala, con una luz por cada destino posible (planta, stock de baja ley y botadero), cuál es el destino que el plan vigente asigna al polígono donde la pala está cargando, sin que la persona operadora haga nada. \\
Actor & Persona operadora de equipo \\
Resultado esperado & La cabina muestra con una señal luminosa el destino del plan vigente para el polígono de carguío, sin que la persona operadora haga nada, y la señal cambia cuando cambia el polígono o la versión del plan. \\
Origen en las bases & Bases técnicas del caso Minería, cap. 9.4; Bases técnicas del caso Minería, cap. 4.3; Bases técnicas del caso Minería, restricción no negociable 8; Bases técnicas del caso Minería, criterio de aceptación 14; Bases técnicas del caso Minería, decisión pendiente 1; Derivado de la sección 9.4 y del criterio de aceptación 14 (registro en terreno sin agregar pasos al operador): la señal por destino forma parte del diseño de la botonera \\
\end{tablasintitulo}

\begin{tablasintitulo}{I{0.200}I{0.800}}{\cab{Campo} & \cab{RF-TER-003}}
Problema & Problema 2 \\
Módulo & Operación Mina \\
Etapa & Etapa 1 \\
Requerimiento & El sistema debe avisar a la persona despachadora (quien asigna los camiones) del aviso de cambio de material mediante una notificación en la aplicación, en no más de 10 segundos en el 95 \% de los casos con cobertura de red, sin pasar por la radio. \\
Actor & Despachador de flota \\
Resultado esperado & El despachador recibe el aviso dentro de 10 segundos (en el 95 \% de los casos, con cobertura), sin radio. \\
Origen en las bases & Bases técnicas del caso Minería, cap. 9.4; Bases técnicas del caso Minería, cap. 4.3; requisito técnico RT-16.21 de las Bases técnicas del caso Minería, equivalente al RT-16.20 de las Bases técnicas transversales; Derivado de RT-16.21 (canales de notificación): el plazo de 10 segundos se fija para reemplazar la radio \\
\end{tablasintitulo}

\begin{tablasintitulo}{I{0.200}I{0.800}}{\cab{Campo} & \cab{RF-TER-004}}
Problema & Problema 2 \\
Módulo & Operación Mina \\
Etapa & Etapa 1 \\
Requerimiento & El sistema debe registrar el cambio de destino que la persona despachadora ejecuta en el sistema de despacho y vincularlo al aviso que lo originó, para que la razón del cambio quede asociada a la decisión. \\
Actor & Despachador de flota \\
Resultado esperado & El cambio de destino queda vinculado al aviso que lo originó. \\
Origen en las bases & Bases técnicas del caso Minería, cap. 4.3; Bases técnicas del caso Minería, cap. 9.4; Bases técnicas del caso Minería, restricción no negociable 5 \\
\end{tablasintitulo}

\begin{tablasintitulo}{I{0.200}I{0.800}}{\cab{Campo} & \cab{RF-TER-005}}
Problema & Problema 2 \\
Módulo & Operación Mina \\
Etapa & Etapa 1 \\
Requerimiento & El sistema debe conservar el aviso de cambio de material cuando el equipo opere hasta 48 horas continuas sin señal y transmitirlo al recuperar cobertura, con sincronización en no más de 4 horas y sin intervención manual, sin que la persona operadora deba repetirlo ni verificarlo. \\
Actor & Persona operadora de equipo \\
Resultado esperado & El aviso queda guardado en el equipo durante hasta 48 horas sin señal y se transmite al recuperar señal en no más de 4 horas, sin intervención manual ni que la persona lo repita. \\
Origen en las bases & Bases técnicas del caso Minería, cap. 4.3; Bases técnicas del caso Minería, cap. 9.4; Bases técnicas del caso Minería, cap. 17.4; Bases técnicas del caso Minería, criterio de aceptación 13; Bases técnicas del caso Minería, decisión pendiente 9; requisito técnico RT-03.10; requisito técnico RT-03.13; supuesto SUP-09 \\
\end{tablasintitulo}

\begin{tablasintitulo}{I{0.200}I{0.800}}{\cab{Campo} & \cab{RF-TER-006}}
Problema & Problema 2 \\
Módulo & Operación Mina \\
Etapa & Etapa 1 \\
Requerimiento & El sistema debe capturar los datos de cada pozo perforado (identificador, posición, profundidad, velocidad de penetración y hora): de forma automática en las perforadoras que ya los miden y, en las que no, mediante registro de la persona operadora desde un dispositivo de terreno que funciona sin señal, en el momento de la perforación, eliminando la planilla de fin de turno. \\
Actor & Persona operadora de equipo \\
Resultado esperado & Cada pozo queda registrado con sus parámetros, sin intervención manual en perforadoras instrumentadas y desde el dispositivo de terreno en las no instrumentadas, sin planilla de fin de turno. \\
Origen en las bases & Bases técnicas del caso Minería, cap. 4.2; Bases técnicas del caso Minería, cap. 9.4; requisito técnico RT-17.01 \\
\end{tablasintitulo}

\begin{tablasintitulo}{I{0.200}I{0.800}}{\cab{Campo} & \cab{RF-TER-007}}
Problema & Problema 2 \\
Módulo & Operación Mina \\
Etapa & Etapa 1 \\
Requerimiento & El sistema debe entregar a geología los datos de perforación en formato estructurado y exportable (archivo CSV con pozo, posición, profundidad y parámetros de perforación, y el estándar abierto de intercambio de datos de equipos mineros cuando el fabricante lo soporte), sin intervención del área de perforación y tronadura y sin trabajo de conversión previo. \\
Actor & Geología \\
Resultado esperado & Geología descarga los datos de perforación en formato estructurado sin conversión previa. \\
Origen en las bases & Bases técnicas del caso Minería, cap. 4.2; requisito técnico RT-05.23; requisito técnico RT-05.22 \\
\end{tablasintitulo}

\begin{tablasintitulo}{I{0.200}I{0.800}}{\cab{Campo} & \cab{RF-TER-008}}
Problema & Problema 2 \\
Módulo & Operación Mina \\
Etapa & Etapa 1 \\
Requerimiento & El sistema debe registrar cada tronadura con fecha, hora, sector, pozos y desplazamiento declarado del contacto (cuánto se movió el límite entre materiales), vinculándola a los polígonos afectados. También debe registrar la liberación del sector (aviso de que quedó disponible para carguío) con responsable, sector y hora, y notificarla a despacho de flota y a la operación mina, reemplazando el aviso por radio. \\
Actor & Perforación y tronadura \\
Resultado esperado & La tronadura queda registrada y vinculada a los polígonos afectados, y la liberación queda registrada y notificada a despacho, mina y geología. \\
Origen en las bases & Bases técnicas del caso Minería, cap. 4.2; requisito técnico RT-16.21 de las Bases técnicas del caso Minería, equivalente al RT-16.20 de las Bases técnicas transversales \\
\end{tablasintitulo}

\begin{tablasintitulo}{I{0.200}I{0.800}}{\cab{Campo} & \cab{RF-PLA-001}}
Problema & Problema 2 \\
Módulo & Operación Mina \\
Etapa & Etapa 1 \\
Requerimiento & El sistema debe consumir automáticamente las salidas del software de planificación minera -polígonos, leyes estimadas, destinos y tonelajes planificados- eliminando la exportación a carpeta compartida y la carga manual en el sistema de despacho. \\
Actor & Planificación minera \\
Resultado esperado & Polígonos, leyes, destinos y tonelajes quedan disponibles en el sistema y en el despacho sin exportación a carpeta ni carga manual. \\
Origen en las bases & Bases técnicas del caso Minería, cap. 4.1; Bases técnicas del caso Minería, restricción no negociable 7; supuesto SUP-16 \\
\end{tablasintitulo}

\begin{tablasintitulo}{I{0.200}I{0.800}}{\cab{Campo} & \cab{RF-PLA-002}}
Problema & Problema 2 \\
Módulo & Operación Mina \\
Etapa & Etapa 1 \\
Requerimiento & El sistema debe guardar el plan minero con versiones: para cada una, su fecha y hora de entrada en vigencia, su autor, el motivo del cambio y la versión que reemplaza. La Superintendencia de Planificación debe poder registrar en el sistema un cambio de plan durante la semana, dejando quién, cuándo, qué y por qué. Cada movimiento de material se asocia automáticamente a la versión vigente cuando ocurre, para conocer el cumplimiento del plan al inicio del turno siguiente contra esa versión. \\
Actor & Planificación minera \\
Resultado esperado & Cada versión del plan queda almacenada con sus atributos, el cambio de plan queda registrado con quién, cuándo, qué y por qué y da origen a una nueva versión vigente, y cada movimiento queda asociado a la versión vigente al ocurrir, con el cumplimiento del turno comparado contra esa versión al inicio del turno siguiente. \\
Origen en las bases & Bases técnicas del caso Minería, cap. 4.1; Bases técnicas del caso Minería, cap. 9.4; Bases técnicas del caso Minería, criterio de aceptación 4; Bases técnicas del caso Minería, decisión pendiente 14; Bases técnicas del caso Minería, entrevista a Ledermann; supuesto SUP-14 \\
\end{tablasintitulo}

\begin{tablasintitulo}{I{0.200}I{0.800}}{\cab{Campo} & \cab{RF-PLA-003}}
Problema & Problema 2 \\
Módulo & Operación Mina \\
Etapa & Etapa 1 \\
Requerimiento & El sistema debe calcular el cumplimiento del plan minero (tonelaje, ley, destino y secuencia, es decir el orden de extracción) comparando lo ejecutado contra la versión del plan efectivamente vigente, y dejarlo disponible al inicio del turno siguiente. \\
Actor & Planificación minera \\
Resultado esperado & El cumplimiento del turno contra la versión vigente queda disponible al inicio del turno siguiente. \\
Origen en las bases & Bases técnicas del caso Minería, cap. 4.1; Bases técnicas del caso Minería, criterio de aceptación 4; supuesto SUP-14 \\
\end{tablasintitulo}

\begin{tablasintitulo}{I{0.200}I{0.800}}{\cab{Campo} & \cab{RF-PLA-004}}
Problema & Problema 2 \\
Módulo & Operación Mina \\
Etapa & Etapa 1 \\
Requerimiento & El sistema debe registrar el motivo de cada cambio de plan eligiéndolo de una lista cerrada (motivos predefinidos: rampa caída, desplazamiento del contacto por tronadura, pedido de planta y los que el Cliente agregue), sin depender de radio ni de mensajería del turno, y notificar el cambio a despacho, mina y planta por los canales de notificación del proyecto dentro de 15 minutos desde que entra en vigencia. \\
Actor & Planificación minera \\
Resultado esperado & El cambio queda registrado con un motivo de la lista cerrada y se notifica a despacho, mina y planta por los canales del Caso dentro de 15 minutos tras su entrada en vigencia. \\
Origen en las bases & Bases técnicas del caso Minería, cap. 4.1; Bases técnicas del caso Minería, decisión pendiente 14; requisito técnico RT-16.21 de las Bases técnicas del caso Minería, equivalente al RT-16.20 de las Bases técnicas transversales; Derivado de RT-16.21 (canales de notificación): el plazo de 15 minutos se fija para que el cambio de plan llegue a tiempo \\
\end{tablasintitulo}

\begin{tablasintitulo}{I{0.200}I{0.800}}{\cab{Campo} & \cab{RF-PLA-005}}
Problema & Problema 2 \\
Módulo & Operación Mina \\
Etapa & Etapa 1 \\
Requerimiento & El sistema debe mantener un catálogo cerrado y con versiones de los tipos de detención, con la definición de cada tipo, su tratamiento en el cálculo de disponibilidad y el rol responsable de clasificarlo. \\
Actor & Operaciones mina \\
Resultado esperado & Cada tipo de detención tiene definición, tratamiento en disponibilidad y rol responsable, con versión. \\
Origen en las bases & Bases técnicas del caso Minería, cap. 7.2; Bases técnicas del caso Minería, criterio de aceptación 5; Bases técnicas del caso Minería, decisión pendiente 5; supuesto SUP-05 \\
\end{tablasintitulo}

\begin{tablasintitulo}{I{0.200}I{0.800}}{\cab{Campo} & \cab{RF-PLA-006}}
Problema & Problema 2 \\
Módulo & Equipos y Mantenimiento \\
Etapa & Etapa 2 \\
Requerimiento & El sistema debe clasificar automáticamente las detenciones de equipos y de planta usando la telemetría, el sistema de despacho y el historiador de proceso, y mostrar al supervisor de turno solo las que no pudo resolver, para que las clasifique desde una lista cerrada, sin texto libre. Todo debe quedar clasificado dentro del turno siguiente al de la detención; las que queden sin clasificar al cierre del turno se escalan al supervisor de la guardia siguiente, con constancia. \\
Actor & Supervisor de turno \\
Resultado esperado & Las detenciones resolubles por fuente automática quedan clasificadas con causa dentro del turno siguiente, el residuo queda clasificado por el supervisor con causa de lista cerrada, y toda detención pendiente al cierre del turno queda escalada al supervisor de la guardia siguiente con constancia. \\
Origen en las bases & Bases técnicas del caso Minería, cap. 7.2; Bases técnicas del caso Minería, restricción no negociable 3; Bases técnicas del caso Minería, criterio de aceptación 5; Bases técnicas del caso Minería, decisión pendiente 5; requisito técnico RT-05.23; supuesto SUP-05 \\
\end{tablasintitulo}

\begin{tablasintitulo}{I{0.200}I{0.800}}{\cab{Campo} & \cab{RF-PLA-007}}
Problema & Problema 2 \\
Módulo & Equipos y Mantenimiento \\
Etapa & Etapa 2 \\
Requerimiento & El sistema debe registrar la causa raíz (la causa de fondo, más profunda que el tipo de detención) de cada detención no programada de equipo y de planta, con responsable y fecha de determinación, y alertar cuando siga sin causa raíz pasadas 72 horas desde la detención. \\
Actor & Supervisor de mantenimiento \\
Resultado esperado & La causa raíz queda registrada con responsable y fecha, y se alerta la detención sin causa raíz pasadas 72 horas. \\
Origen en las bases & Bases técnicas del caso Minería, cap. 7.2; Bases técnicas del caso Minería, criterio de aceptación 5; supuesto SUP-05 \\
\end{tablasintitulo}

\begin{tablasintitulo}{I{0.200}I{0.800}}{\cab{Campo} & \cab{RF-PLA-008}}
Problema & Problema 2 \\
Módulo & Operación Mina \\
Etapa & Etapa 1 \\
Requerimiento & El sistema debe calcular y publicar por equipo, área y turno la disponibilidad física, la utilización efectiva, las horas de detención no programada y el porcentaje de tiempo registrado sin clasificar. \\
Actor & Operaciones mina \\
Resultado esperado & Disponibilidad física, utilización efectiva, horas de detención y porcentaje sin clasificar quedan publicados por equipo, área y turno. \\
Origen en las bases & Bases técnicas del caso Minería, cap. 7.2; Bases técnicas del caso Minería, criterio de aceptación 5; supuesto SUP-05 \\
\end{tablasintitulo}

\begin{tablasintitulo}{I{0.200}I{0.800}}{\cab{Campo} & \cab{RF-PLA-009}}
Problema & Problema 2 \\
Módulo & Operación Mina \\
Etapa & Etapa 1 \\
Requerimiento & El sistema debe registrar las horas de detención no programada de planta leyendo del historiador de proceso. \\
Actor & Metalurgia y planta \\
Resultado esperado & Las horas de detención no programada de planta quedan registradas sin digitación. \\
Origen en las bases & Bases técnicas del caso Minería, cap. 7.2; Bases técnicas del caso Minería, restricción no negociable 3; requisito técnico RT-05.23; requisito técnico RT-05.29; Traducción del proponente del indicador del problema de la sección 7.2 (214 horas de detención no programada de planta en el año, frente a la meta interna de bajo 120), necesaria como insumo para la clasificación con causa del criterio de aceptación 5 (ver RF-PLA-006) \\
\end{tablasintitulo}

\begin{tablasintitulo}{I{0.200}I{0.800}}{\cab{Campo} & \cab{RF-ACR-001}}
Problema & Problema 3 \\
Módulo & Acreditación y Acceso \\
Etapa & Etapa 1 \\
Requerimiento & El sistema debe mantener un expediente digital por persona con sus contratos, exámenes, cursos, protocolos de vigilancia de salud y los demás documentos que exige la habilitación, cada uno con su fecha de emisión, su fecha de vencimiento y su estado, reemplazando la planilla que hoy mantiene una sola persona. \\
Actor & Seguridad y salud ocupacional \\
Resultado esperado & Existe un expediente digital por persona con cada documento, su vencimiento y su estado, y la planilla deja de usarse. \\
Origen en las bases & Bases técnicas del caso Minería, cap. 9.5; Bases técnicas del caso Minería, cap. 4.8; Bases técnicas del caso Minería, criterio de aceptación 8 \\
\end{tablasintitulo}

\begin{tablasintitulo}{I{0.200}I{0.800}}{\cab{Campo} & \cab{RF-ACR-002}}
Problema & Problema 3 \\
Módulo & Acreditación y Acceso \\
Etapa & Etapa 1 \\
Requerimiento & El sistema debe ejecutar el flujo de acreditación de una persona contratista con etapas, responsables y plazos, registrando la hora de entrada y de salida de cada etapa. \\
Actor & Seguridad y salud ocupacional \\
Resultado esperado & La acreditación avanza por etapas con responsable y plazo, y queda registrada la hora de entrada y salida de cada etapa. \\
Origen en las bases & Bases técnicas del caso Minería, cap. 9.5; Bases técnicas del caso Minería, cap. 4.8; Bases técnicas del caso Minería, criterio de aceptación 8 \\
\end{tablasintitulo}

\begin{tablasintitulo}{I{0.200}I{0.800}}{\cab{Campo} & \cab{RF-ACR-003}}
Problema & Problema 3 \\
Módulo & Acreditación y Acceso \\
Etapa & Etapa 1 \\
Requerimiento & El sistema debe validar automáticamente la completitud y la vigencia de la documentación cargada e indicar a la empresa contratista qué documento falta o está vencido, sin intervención del área de seguridad. \\
Actor & Empresa contratista \\
Resultado esperado & La empresa ve qué documento falta o está vencido sin intervención del área de seguridad. \\
Origen en las bases & Bases técnicas del caso Minería, cap. 9.5; Bases técnicas del caso Minería, cap. 4.8; Bases técnicas del caso Minería, criterio de aceptación 8 \\
\end{tablasintitulo}

\begin{tablasintitulo}{I{0.200}I{0.800}}{\cab{Campo} & \cab{RF-ACR-004}}
Problema & Problema 3 \\
Módulo & Acreditación y Acceso \\
Etapa & Etapa 1 \\
Requerimiento & El sistema debe permitir consultar en cualquier momento el estado de habilitación de cada persona presente en faena, con el detalle del documento que lo sustenta y el histórico auditable de sus cambios. \\
Actor & Seguridad y salud ocupacional \\
Resultado esperado & El estado de habilitación y su historial se consultan en cualquier momento con el documento que lo sustenta. \\
Origen en las bases & Bases técnicas del caso Minería, cap. 4.8; Bases técnicas del caso Minería, cap. 9.5 \\
\end{tablasintitulo}

\begin{tablasintitulo}{I{0.200}I{0.800}}{\cab{Campo} & \cab{RF-ACR-005}}
Problema & Problema 3 \\
Módulo & Acreditación y Acceso \\
Etapa & Etapa 1 \\
Requerimiento & El sistema debe recibir la documentación de acreditación a través de un portal autenticado para empresas contratistas, eliminando la recepción por correo electrónico y el registro manual. \\
Actor & Empresa contratista \\
Resultado esperado & La documentación se recibe por el portal, sin correo ni registro manual. \\
Origen en las bases & Bases técnicas del caso Minería, cap. 9.5; Bases técnicas del caso Minería, cap. 4.8; requisito técnico RT-16.30 de las Bases técnicas del caso Minería (portal autenticado para empresas contratistas); requisito técnico RT-12.12; supuesto SUP-08 \\
\end{tablasintitulo}

\begin{tablasintitulo}{I{0.200}I{0.800}}{\cab{Campo} & \cab{RF-ACR-006}}
Problema & Problema 3 \\
Módulo & Acreditación y Acceso \\
Etapa & Etapa 1 \\
Requerimiento & El sistema debe avisar del vencimiento próximo de cada documento de habilitación con 30, 15 y 7 días de anticipación (valores por defecto, modificables), notificando a la vez al trabajador, a su empleador y al área de seguridad. Si a los 7 días el documento no se ha renovado, la alerta se marca como escalada, y al llegar la fecha de vencimiento se bloquea la habilitación. \\
Actor & Seguridad y salud ocupacional \\
Resultado esperado & Trabajador, empleador y área de seguridad reciben la alerta a los 30, 15 y 7 días antes del vencimiento, la alerta de 7 días queda marcada como escalada si el documento no se ha renovado y la habilitación se bloquea al vencer el documento. \\
Origen en las bases & Bases técnicas del caso Minería, cap. 4.8; Bases técnicas del caso Minería, cap. 9.5; Bases técnicas del caso Minería, criterio de aceptación 7; Bases técnicas del caso Minería, entrevista a Tapia; Derivado de la entrevista a Tapia y del criterio de aceptación 7 (avisos con semanas de anticipación): los hitos de 30, 15 y 7 días, la escalada y el bloqueo al vencimiento concretan esos avisos \\
\end{tablasintitulo}

\begin{tablasintitulo}{I{0.200}I{0.800}}{\cab{Campo} & \cab{RF-ACR-007}}
Problema & Problema 3 \\
Módulo & Acreditación y Acceso \\
Etapa & Etapa 1 \\
Requerimiento & El sistema debe calcular el vencimiento de cada protocolo de vigilancia de salud (altura geográfica y exposición a sílice) según su periodicidad, anual por defecto y configurable por protocolo conforme a la normativa, y tratarlo como documento de habilitación de la persona, con los avisos previos y el bloqueo al vencimiento de RF-ACR-006. \\
Actor & Seguridad y salud ocupacional \\
Resultado esperado & Cada protocolo de vigilancia de la persona queda con su vencimiento calculado según su periodicidad, y ese vencimiento genera los avisos previos y el bloqueo de la habilitación de RF-ACR-006. \\
Origen en las bases & Bases técnicas del caso Minería, cap. 4.8; Bases técnicas del caso Minería, cap. 9.5; Bases técnicas del caso Minería, cap. 12; Derivado de la sección 12 (protocolos de vigilancia): la periodicidad anual por defecto se fija a partir de ella \\
\end{tablasintitulo}

\begin{tablasintitulo}{I{0.200}I{0.800}}{\cab{Campo} & \cab{RF-ACR-008}}
Problema & Problema 3 \\
Módulo & Acreditación y Acceso \\
Etapa & Etapa 1 \\
Requerimiento & El sistema debe verificar si la persona que presenta su credencial en el torniquete está habilitada, y registrar el resultado del ingreso o del rechazo con su causa, su hora y el punto de control. \\
Actor & Guardia de portería \\
Resultado esperado & La habilitación en portería queda resuelta y el ingreso o rechazo queda registrado con causa, hora y punto de control. \\
Origen en las bases & Bases técnicas del caso Minería, cap. 9.5; Bases técnicas del caso Minería, cap. 4.8; supuesto SUP-13 \\
\end{tablasintitulo}

\begin{tablasintitulo}{I{0.200}I{0.800}}{\cab{Campo} & \cab{RF-ACR-009}}
Problema & Problema 3 \\
Módulo & Acreditación y Acceso \\
Etapa & Etapa 2 \\
Requerimiento & El sistema debe verificar la identidad de la persona que entra o sale por portería con reconocimiento facial uno a uno contra la persona asociada a la credencial presentada, guardando una plantilla cifrada que no permite reconstruir el rasgo (no la imagen) y usándola solo para verificar identidad en el ingreso y la salida. Solo se activa con acuerdo previo del sindicato. Si el rostro no coincide con la persona asociada a la credencial, el ingreso no se autoriza y queda registrado con esa causa. Si la persona no tiene plantilla enrolada o el lector facial está fuera de servicio, el ingreso se controla con la credencial y la validación de habilitación. \\
Actor & Guardia de portería \\
Resultado esperado & La identidad queda verificada con biometría facial uno a uno o, si la persona no tiene plantilla enrolada o el lector facial está fuera de servicio, con credencial y validación de habilitación, y se registra el resultado. \\
Origen en las bases & Bases técnicas del caso Minería, cap. 4.8; Bases técnicas del caso Minería, cap. 9.5; requisito técnico RT-05.08; Bases técnicas del caso Minería, decisión pendiente 13; Ley N° 21.719; Bases técnicas del caso Minería, entrevista a Tapia; supuesto SUP-13 \\
\end{tablasintitulo}

\begin{tablasintitulo}{I{0.200}I{0.800}}{\cab{Campo} & \cab{RF-ACR-010}}
Problema & Problema 3 \\
Módulo & Acreditación y Acceso \\
Etapa & Etapa 2 \\
Requerimiento & El sistema debe registrar el consentimiento informado de la persona cuando el tratamiento de un dato personal sensible se base en su consentimiento, como la plantilla biométrica de portería, conservar la versión del consentimiento y permitir su revocación, dejando registro de ella. \\
Actor & Persona contratista \\
Resultado esperado & El consentimiento queda registrado con versión y la revocación queda trazable. \\
Origen en las bases & Bases técnicas del caso Minería, cap. 4.8; Bases técnicas del caso Minería, cap. 9.5; Bases técnicas del caso Minería, cap. 12; requisito técnico RT-05.08; Bases administrativas, artículo 85; Ley N° 21.719; supuesto SUP-13; Derivado de la Ley 21.719 (datos personales sensibles, sección 12): el registro del consentimiento y su revocación permiten cumplirla \\
\end{tablasintitulo}

\begin{tablasintitulo}{I{0.200}I{0.800}}{\cab{Campo} & \cab{RF-ACR-011}}
Problema & Problema 3 \\
Módulo & Acreditación y Acceso \\
Etapa & Etapa 1 \\
Requerimiento & El sistema debe identificar a la persona operadora de un equipo compartido en el relevo de turno leyendo su credencial personal y contrastándola con el operador que el despacho asigna al equipo, sin usuario, sin contraseña y sin pantallas adicionales. Si nadie se identifica, el registro queda marcado como operador no identificado y se alerta al supervisor de turno, sin detener el equipo. \\
Actor & Persona operadora de equipo \\
Resultado esperado & La persona queda identificada en el equipo sin login ni contraseña o, si nadie se identifica, el registro queda como operador no identificado con alerta al supervisor, sin detener el equipo. \\
Origen en las bases & Bases técnicas del caso Minería, cap. 4.8; Bases técnicas del caso Minería, cap. 9.5; Bases técnicas del caso Minería, restricción no negociable 8; Bases técnicas del caso Minería, decisión pendiente 4; requisito técnico RT-12.11; supuesto SUP-04; supuesto SUP-17 \\
\end{tablasintitulo}

\begin{tablasintitulo}{I{0.200}I{0.800}}{\cab{Campo} & \cab{RF-ACR-012}}
Problema & Problema 3 \\
Módulo & Acreditación y Acceso \\
Etapa & Etapa 1 \\
Requerimiento & El sistema debe autenticar a las personas contratistas en terreno sin correo corporativo ni contraseña escrita, usando su credencial y un PIN numérico. \\
Actor & Persona contratista \\
Resultado esperado & La persona se autentica en terreno sin correo corporativo ni contraseña alfanumérica. \\
Origen en las bases & Bases técnicas del caso Minería, cap. 4.8; Bases técnicas del caso Minería, cap. 9.5; Bases técnicas del caso Minería, decisión pendiente 8; requisito técnico RT-12.11; requisito técnico RT-12.12; supuesto SUP-08 \\
\end{tablasintitulo}

\begin{tablasintitulo}{I{0.200}I{0.800}}{\cab{Campo} & \cab{RF-ACR-013}}
Problema & Problema 3 \\
Módulo & Acreditación y Acceso \\
Etapa & Etapa 1 \\
Requerimiento & El sistema debe registrar el vínculo entre persona, empresa contratista y faena y, al término del contrato o de la asignación, dar de baja automáticamente la habilitación y los accesos de la persona, sin intervención manual. \\
Actor & Empresa contratista \\
Resultado esperado & El vínculo queda registrado mientras dure el contrato o la asignación y, a su término, la habilitación y los accesos quedan dados de baja automáticamente. \\
Origen en las bases & Bases técnicas del caso Minería, cap. 9.5; Bases técnicas del caso Minería, cap. 4.8; requisito técnico RT-12.10; requisito técnico RT-12.12; supuesto SUP-08 \\
\end{tablasintitulo}

\begin{tablasintitulo}{I{0.200}I{0.800}}{\cab{Campo} & \cab{RF-ACR-014}}
Problema & Problema 3 \\
Módulo & Acreditación y Acceso \\
Etapa & Etapa 1 \\
Requerimiento & El sistema debe conservar en cada registro operacional y de seguridad la identidad de la persona que lo generó, de modo que ningún registro quede sin autor identificable salvo el caso declarado de operador no identificado. \\
Actor & Seguridad y salud ocupacional \\
Resultado esperado & Todo registro conserva la identidad de quien lo generó, salvo el caso declarado de operador no identificado. \\
Origen en las bases & Bases técnicas del caso Minería, cap. 4.8; Bases técnicas del caso Minería, cap. 9.5; Bases técnicas del caso Minería, criterio de aceptación 9; Bases técnicas del caso Minería, decisión pendiente 4; supuesto SUP-04 \\
\end{tablasintitulo}

\begin{tablasintitulo}{I{0.200}I{0.800}}{\cab{Campo} & \cab{RF-ACR-015}}
Problema & Problema 3 \\
Módulo & Acreditación y Acceso \\
Etapa & Etapa 1 \\
Requerimiento & El sistema debe permitir a las personas usuarias externas (empresas contratistas y compradores) resolver por su cuenta en el portal las consultas más frecuentes: estado de acreditación, documentos vencidos o por vencer y recuperación de acceso. Los compradores consultan el estado de sus lotes en el Portal para Compradores de la Etapa 2 (RF-TRZ-016). \\
Actor & Empresa contratista \\
Resultado esperado & La consulta se resuelve en el portal sin contacto telefónico. \\
Origen en las bases & requisito técnico RT-16.32; requisito técnico RT-12.12 \\
\end{tablasintitulo}

\begin{tablasintitulo}{I{0.200}I{0.800}}{\cab{Campo} & \cab{RF-ACR-016}}
Problema & Problema 3 \\
Módulo & Acreditación y Acceso \\
Etapa & Etapa 1 \\
Requerimiento & El sistema debe controlar el acceso por roles, complementado con atributos -sitio, empresa contratista y área- donde el proceso lo exija, y aplicar una matriz de segregación de funciones documentada y verificable. \\
Actor & TI del Cliente \\
Resultado esperado & La persona accede sólo a lo que su rol y sus atributos permiten, y una combinación de roles prohibida por la matriz queda rechazada. \\
Origen en las bases & Bases técnicas del caso Minería, restricción no negociable 9; requisito técnico RT-12.05 \\
\end{tablasintitulo}

\begin{tablasintitulo}{I{0.200}I{0.800}}{\cab{Campo} & \cab{RF-ACR-017}}
Problema & Transversal \\
Módulo & Acreditación y Acceso \\
Etapa & Etapa 1 \\
Requerimiento & El sistema debe verificar la identidad de cada persona que ingresa a la sala técnica, en el acceso principal y nuevamente en la esclusa al término del pasillo, con biometría facial uno a uno como mecanismo principal y, como respaldo, con huella dactilar (AFIS) cuando el lector facial esté fuera de servicio o no logre verificar a la persona. La esclusa se libera solo tras una verificación positiva y para una sola persona a la vez; si ni la verificación facial ni la huella resultan positivas, el acceso no se libera y el intento queda registrado. \\
Actor & Personal técnico / operador de sala técnica \\
Resultado esperado & Solo ingresan a la sala técnica personas con identidad verificada por biometría facial o, como respaldo, por AFIS, de a una por vez a través de la esclusa, y cada intento sin verificación positiva queda registrado con fecha, hora y punto de acceso. \\
Origen en las bases & requisito técnico RT-06.20; requisito técnico RT-06.22; requisito técnico RT-06.23 \\
\end{tablasintitulo}

\begin{tablasintitulo}{I{0.200}I{0.800}}{\cab{Campo} & \cab{RF-ACR-018}}
Problema & Transversal \\
Módulo & Acreditación y Acceso \\
Etapa & Etapa 1 \\
Requerimiento & El sistema debe registrar cada ingreso y cada egreso de la sala técnica en una bitácora que nadie puede modificar ni eliminar, con la identidad de la persona, la fecha, la hora y el motivo que la persona declara al ingresar. \\
Actor & Personal técnico / operador de sala técnica \\
Resultado esperado & Cada ingreso y egreso de la sala técnica queda en la bitácora con identidad, fecha, hora y motivo, sin posibilidad de modificarlo ni eliminarlo. \\
Origen en las bases & requisito técnico RT-06.21 \\
\end{tablasintitulo}

\begin{tablasintitulo}{I{0.200}I{0.800}}{\cab{Campo} & \cab{RF-CUM-001}}
Problema & Problema 3 \\
Módulo & Seguridad y Salud Ocupacional \\
Etapa & Etapa 2 \\
Requerimiento & El sistema debe permitir registrar incidentes, cuasi accidentes, observaciones e inspecciones de seguridad en el lugar del hecho, desde un dispositivo móvil que funciona sin señal, en el momento en que ocurren. \\
Actor & Supervisor de turno \\
Resultado esperado & El incidente, cuasi accidente, observación o inspección queda registrado en el lugar y en el momento. \\
Origen en las bases & Bases técnicas del caso Minería, cap. 4.8; Bases técnicas del caso Minería, criterio de aceptación 9; requisito técnico RT-17.01 \\
\end{tablasintitulo}

\begin{tablasintitulo}{I{0.200}I{0.800}}{\cab{Campo} & \cab{RF-CUM-002}}
Problema & Problema 3 \\
Módulo & Seguridad y Salud Ocupacional \\
Etapa & Etapa 2 \\
Requerimiento & El sistema debe registrar, junto a cada registro de seguridad, la fecha y hora del hecho y la fecha y hora del registro. \\
Actor & Seguridad y salud ocupacional \\
Resultado esperado & El registro conserva la hora del hecho y la hora del registro. \\
Origen en las bases & Bases técnicas del caso Minería, cap. 4.8; Bases técnicas del caso Minería, criterio de aceptación 9; Derivado de la sección 4.8 y del criterio de aceptación 9 (registro digital en el lugar del hecho): se separan la hora del hecho y la del registro para que la evidencia sea auditable \\
\end{tablasintitulo}

\begin{tablasintitulo}{I{0.200}I{0.800}}{\cab{Campo} & \cab{RF-CUM-003}}
Problema & Problema 3 \\
Módulo & Seguridad y Salud Ocupacional \\
Etapa & Etapa 2 \\
Requerimiento & El sistema debe generar el expediente de fiscalización (los registros del período que pide la autoridad, con su evidencia) en formato exportable, sin armar carpetas a mano. \\
Actor & Seguridad y salud ocupacional \\
Resultado esperado & El expediente del período se genera desde el sistema con su evidencia y en formato exportable. \\
Origen en las bases & Bases técnicas del caso Minería, cap. 4.8; Bases técnicas del caso Minería, criterio de aceptación 9 \\
\end{tablasintitulo}

\begin{tablasintitulo}{I{0.200}I{0.800}}{\cab{Campo} & \cab{RF-ADM-005}}
Problema & Problema 3 \\
Módulo & Administración y Servicios Comunes \\
Etapa & Etapa 1 \\
Requerimiento & El sistema debe registrar toda operación que cree, modifique o elimine información: quién (persona o sistema), cuándo (fecha y hora con zona horaria), desde dónde y con qué dispositivo, y los valores antes y después. Ese registro de auditoría no puede ser modificado ni eliminado por nadie, ni siquiera por el administrador. \\
Actor & Cumplimiento y auditoría interna \\
Resultado esperado & Para cualquier registro se reconstruye quién, qué, cuándo, desde dónde y con qué valores antes y después durante todo su plazo de retención, y todo intento de alterar la auditoría queda rechazado y registrado. \\
Origen en las bases & Bases técnicas del caso Minería, restricción no negociable 9; requisito técnico RT-05.03; requisito técnico RT-16.06; requisito técnico RT-16.07 \\
\end{tablasintitulo}

\begin{tablasintitulo}{I{0.200}I{0.800}}{\cab{Campo} & \cab{RF-ADM-006}}
Problema & Problema 3 \\
Módulo & Administración y Servicios Comunes \\
Etapa & Etapa 1 \\
Requerimiento & El sistema debe permitir consultar y exportar la auditoría desde la interfaz, con filtros por persona, período, entidad y tipo de operación, sin acceso a la base de datos, y registrar en la auditoría toda exportación de información sensible. \\
Actor & Cumplimiento y auditoría interna \\
Resultado esperado & La auditoría filtrada se exporta en formato abierto, y la exportación queda registrada con quién, qué y cuándo. \\
Origen en las bases & requisito técnico RT-16.08; requisito técnico RT-16.30; Consulta N.º 18 \\
\end{tablasintitulo}

\begin{tablasintitulo}{I{0.200}I{0.800}}{\cab{Campo} & \cab{RF-ADM-011}}
Problema & Problema 3 \\
Módulo & Administración y Servicios Comunes \\
Etapa & Etapa 1 \\
Requerimiento & El sistema debe generar documentos -actas, certificados, expedientes y reportes- a partir de plantillas administrables por el cliente, con los datos de la transacción y salida en formato abierto. \\
Actor & Administrador funcional del Cliente \\
Resultado esperado & El documento se genera con los datos de la transacción, sin edición manual, en formato abierto. \\
Origen en las bases & requisito técnico RT-16.19 \\
\end{tablasintitulo}

\begin{tablasintitulo}{I{0.200}I{0.800}}{\cab{Campo} & \cab{RF-NOT-001}}
Problema & Problema 3 \\
Módulo & Administración y Servicios Comunes \\
Etapa & Etapa 1 \\
Requerimiento & El sistema debe enviar notificaciones por correo electrónico, notificación en aplicación, mensajería instantánea para personas contratistas, mensaje de texto para alertas críticas y de seguridad, e integración con el sistema de radio y megafonía de faena para las alertas de evacuación. \\
Actor & Seguridad y salud ocupacional \\
Resultado esperado & El mensaje llega por el canal definido para el evento, y la alerta de evacuación sale por radio y megafonía. \\
Origen en las bases & requisito técnico RT-16.21 de las Bases técnicas del caso Minería, equivalente al RT-16.20 de las Bases técnicas transversales; Consulta N.º 17 \\
\end{tablasintitulo}

\begin{tablasintitulo}{I{0.200}I{0.800}}{\cab{Campo} & \cab{RF-DAT-002}}
Problema & Problema 3 \\
Módulo & Administración y Servicios Comunes \\
Etapa & Etapa 1 \\
Requerimiento & El sistema debe eliminar o anonimizar de forma segura y verificable los datos cuyo plazo de retención haya vencido, conforme a la política de retención por dominio, dejando constancia auditable de cada eliminación. \\
Actor & Cumplimiento y auditoría interna \\
Resultado esperado & Los registros quedan eliminados o anonimizados y existe constancia auditable de la operación. \\
Origen en las bases & requisito técnico RT-05.07; Ley N° 21.719; Consulta N.º 10 \\
\end{tablasintitulo}

\begin{tablasintitulo}{I{0.200}I{0.800}}{\cab{Campo} & \cab{RF-CUM-004}}
Problema & Problema 3 \\
Módulo & Seguridad y Salud Ocupacional \\
Etapa & Etapa 2 \\
Requerimiento & El sistema debe soportar el programa de inspecciones planificadas de seguridad, cargado mensualmente, con su calendario diario, su responsable y su resultado, dimensionado para unas 8.200 inspecciones al año (unas 680 al mes y 22 al día; es la proyección a 3 años del Caso). \\
Actor & Inspector de seguridad \\
Resultado esperado & Cada inspección planificada del mes queda con su fecha programada, su responsable y su resultado. \\
Origen en las bases & Bases técnicas del caso Minería, cap. 4.8; Bases técnicas del caso Minería, criterio de aceptación 9; requisito técnico RT-17.01; Bases técnicas del caso Minería, cap. 14.1; Derivado de la sección 14.1 (volumetría de inspecciones): el calendario mensual y diario se calcula a partir de las 8.200 inspecciones al año \\
\end{tablasintitulo}

\begin{tablasintitulo}{I{0.200}I{0.800}}{\cab{Campo} & \cab{RF-ADM-010}}
Problema & Problema 3 \\
Módulo & Administración y Servicios Comunes \\
Etapa & Etapa 1 \\
Requerimiento & El sistema debe gestionar documentos con versiones, metadatos, control de acceso, búsqueda por contenido y por metadato y vista previa sin descargar. Debe guardarlos cifrados, verificar que no se alteren y conservarlos 10 años (seguridad y salud ocupacional, ambientales y trazabilidad de mineral y embarques), 7 años (tributarios) o 5 años (operacionales). \\
Actor & Seguridad y salud ocupacional \\
Resultado esperado & El documento se recupera por versión y por metadato, cualquier alteración de su contenido queda detectada y se conserva durante el plazo de retención de su categoría. \\
Origen en las bases & requisito técnico RT-16.15; requisito técnico RT-16.16; requisito técnico RT-05.10; requisito técnico RT-16.10; requisito técnico RT-05.07 \\
\end{tablasintitulo}

\begin{tablasintitulo}{I{0.200}I{0.800}}{\cab{Campo} & \cab{RF-CUM-005}}
Problema & Problema 3 \\
Módulo & Seguridad y Salud Ocupacional \\
Etapa & Etapa 2 \\
Requerimiento & El sistema debe recibir las alertas de fatiga y somnolencia desde la plataforma de las cámaras de control instaladas en los equipos, registrando cada una con equipo, persona operadora, marca de tiempo y severidad, y exigir una acción documentada para cada alerta, con responsable, tipo de acción y marca de tiempo, escalándola cuando no se registre dentro de 30 minutos desde la alerta. \\
Actor & Supervisor de turno \\
Resultado esperado & La alerta queda registrada con equipo, persona, hora y severidad, y su acción queda documentada con responsable, tipo y hora, o escalada a los 30 minutos. \\
Origen en las bases & Bases técnicas del caso Minería, cap. 4.8; Bases técnicas del caso Minería, criterio de aceptación 10; requisito técnico RT-17.06; Derivado del criterio de aceptación 10 y de RT-17.06 (toda alerta de fatiga con acción documentada): el plazo de escalamiento de 30 minutos concreta ese seguimiento \\
\end{tablasintitulo}

\begin{tablasintitulo}{I{0.200}I{0.800}}{\cab{Campo} & \cab{RF-CUM-006}}
Problema & Problema 3 \\
Módulo & Seguridad y Salud Ocupacional \\
Etapa & Etapa 2 \\
Requerimiento & El sistema debe conservar durante 10 años el histórico completo de alertas de fatiga y de las acciones tomadas sobre ellas, de modo que ante un accidente por somnolencia la compañía pueda acreditar qué hizo con las alertas previas. \\
Actor & Seguridad y salud ocupacional \\
Resultado esperado & El historial de alertas y acciones se consulta completo durante 10 años. \\
Origen en las bases & Bases técnicas del caso Minería, cap. 4.8; Bases técnicas del caso Minería, criterio de aceptación 10; requisito técnico RT-05.10; requisito técnico RT-16.10; requisito técnico RT-05.07; Bases técnicas del caso Minería, entrevistas \\
\end{tablasintitulo}

\begin{tablasintitulo}{I{0.200}I{0.800}}{\cab{Campo} & \cab{RF-CUM-007}}
Problema & Problema 3 \\
Módulo & Seguridad y Salud Ocupacional \\
Etapa & Etapa 2 \\
Requerimiento & El sistema debe restringir el acceso al dato individual de las alertas de fatiga a perfiles autorizados y registrar las consultas a ese dato, además de las modificaciones. \\
Actor & Seguridad y salud ocupacional \\
Resultado esperado & El dato individual de fatiga solo es accesible a perfiles autorizados y cada consulta y modificación queda registrada. \\
Origen en las bases & Bases técnicas del caso Minería, cap. 4.8; Bases técnicas del caso Minería, cap. 12; requisito técnico RT-16.09; Ley N° 21.719 \\
\end{tablasintitulo}

\begin{tablasintitulo}{I{0.200}I{0.800}}{\cab{Campo} & \cab{RF-CUM-008}}
Problema & Problema 3 \\
Módulo & Cumplimiento Ambiental \\
Etapa & Etapa 2 \\
Requerimiento & El sistema debe capturar automáticamente los datos de la instrumentación de monitoreo ambiental y geotécnico, con punto de medición, marca de tiempo y valor, sin transcripción manual. \\
Actor & Sustentabilidad y asuntos regulatorios \\
Resultado esperado & La medición queda almacenada con punto, hora y valor, sin transcripción. \\
Origen en las bases & Bases técnicas del caso Minería, cap. 9.6; Bases técnicas del caso Minería, cap. 4.9; requisito técnico RT-17.06; supuesto SUP-29 \\
\end{tablasintitulo}

\begin{tablasintitulo}{I{0.200}I{0.800}}{\cab{Campo} & \cab{RF-CUM-009}}
Problema & Problema 3 \\
Módulo & Cumplimiento Ambiental \\
Etapa & Etapa 2 \\
Requerimiento & El sistema debe recibir por integración los resultados de los laboratorios externos de calidad del aire y de la empresa geotécnica de relaves. \\
Actor & Sustentabilidad y asuntos regulatorios \\
Resultado esperado & El resultado queda registrado como dato estructurado vinculado a su punto de medición. \\
Origen en las bases & Bases técnicas del caso Minería, cap. 9.6; Bases técnicas del caso Minería, cap. 4.9 \\
\end{tablasintitulo}

\begin{tablasintitulo}{I{0.200}I{0.800}}{\cab{Campo} & \cab{RF-CUM-010}}
Problema & Problema 3 \\
Módulo & Cumplimiento Ambiental \\
Etapa & Etapa 2 \\
Requerimiento & El sistema debe registrar el consumo de agua por punto de medición y permitir declarar la regla de reparto del consumo entre áreas (planta, mina, campamento y supresión de polvo) como registro con versiones, con vigencia y responsable, en reemplazo de la regla fija de 2019. \\
Actor & Sustentabilidad y asuntos regulatorios \\
Resultado esperado & El consumo de agua queda registrado por punto de medición y repartido entre áreas con la regla vigente y su versión. \\
Origen en las bases & Bases técnicas del caso Minería, cap. 9.6; Bases técnicas del caso Minería, cap. 4.9; Bases técnicas del caso Minería, criterio de aceptación 11; Bases técnicas del caso Minería, entrevista a Guerrero; supuesto SUP-28 \\
\end{tablasintitulo}

\begin{tablasintitulo}{I{0.200}I{0.800}}{\cab{Campo} & \cab{RF-CUM-011}}
Problema & Problema 3 \\
Módulo & Cumplimiento Ambiental \\
Etapa & Etapa 2 \\
Requerimiento & El sistema debe registrar el combustible efectivamente consumido por equipo y por área, obtenido de la medición de consumo y no de los litros comprados. \\
Actor & Sustentabilidad y asuntos regulatorios \\
Resultado esperado & El combustible consumido queda registrado por equipo y por área. \\
Origen en las bases & Bases técnicas del caso Minería, cap. 9.6; Bases técnicas del caso Minería, cap. 4.9; Bases técnicas del caso Minería, criterio de aceptación 12; supuesto SUP-28 \\
\end{tablasintitulo}

\begin{tablasintitulo}{I{0.200}I{0.800}}{\cab{Campo} & \cab{RF-CUM-012}}
Problema & Problema 3 \\
Módulo & Cumplimiento Ambiental \\
Etapa & Etapa 2 \\
Requerimiento & El sistema debe generar los reportes a la autoridad ambiental desde los datos registrados, con trazabilidad hasta la medición de origen de cada cifra, y aplicarles firma electrónica avanzada al remitirlos. \\
Actor & Sustentabilidad y asuntos regulatorios \\
Resultado esperado & El reporte se genera desde el sistema, con trazabilidad al origen de cada cifra y con firma electrónica avanzada. \\
Origen en las bases & Bases técnicas del caso Minería, cap. 9.6; Bases técnicas del caso Minería, cap. 4.9; Bases técnicas del caso Minería, criterio de aceptación 11; requisito técnico RT-16.14 de las Bases técnicas del caso Minería, equivalente al RT-16.17 de las Bases técnicas transversales \\
\end{tablasintitulo}

\begin{tablasintitulo}{I{0.200}I{0.800}}{\cab{Campo} & \cab{RF-CUM-013}}
Problema & Problema 3 \\
Módulo & Cumplimiento Ambiental \\
Etapa & Etapa 2 \\
Requerimiento & El sistema debe registrar los compromisos de la Resolución de Calificación Ambiental con su frecuencia, su punto de medición y su destinatario, y alertar el vencimiento de cada obligación de reporte. \\
Actor & Sustentabilidad y asuntos regulatorios \\
Resultado esperado & Cada compromiso queda con frecuencia, punto y destinatario, y se alerta el vencimiento de cada reporte. \\
Origen en las bases & Bases técnicas del caso Minería, cap. 4.9; Bases técnicas del caso Minería, cap. 9.6; Bases técnicas del caso Minería, cap. 12 \\
\end{tablasintitulo}

\begin{tablasintitulo}{I{0.200}I{0.800}}{\cab{Campo} & \cab{RF-CUM-014}}
Problema & Problema 3 \\
Módulo & Cumplimiento Ambiental \\
Etapa & Etapa 2 \\
Requerimiento & El sistema debe construir los indicadores de empleo local y de compras locales del convenio con la comunidad desde el dato de origen, con la periodicidad comprometida y sin planillas solicitadas a otras áreas. \\
Actor & Sustentabilidad y asuntos regulatorios \\
Resultado esperado & Los indicadores de empleo local y compras locales se obtienen desde el sistema sin pedir planillas. \\
Origen en las bases & Bases técnicas del caso Minería, cap. 4.9; Bases técnicas del caso Minería, cap. 9.6; Bases técnicas del caso Minería, cap. 7.3; Bases técnicas del caso Minería, cap. 13.2 \\
\end{tablasintitulo}

\begin{tablasintitulo}{I{0.200}I{0.800}}{\cab{Campo} & \cab{RF-CUM-015}}
Problema & Problema 3 \\
Módulo & Cumplimiento Ambiental \\
Etapa & Etapa 2 \\
Requerimiento & El sistema debe calcular cada mes la huella de carbono por tonelada de cobre fino en sus tres alcances: 1 (emisiones directas de la operación), 2 (energía comprada) y 3 (emisiones de terceros: transporte del concentrado a puerto y por mar, insumos comprados según SAP y traslado del personal a la faena). El cálculo parte de los consumos de combustible y energía registrados por área. \\
Actor & Sustentabilidad y asuntos regulatorios \\
Resultado esperado & La huella de carbono por tonelada de cobre fino del mes queda calculada por alcance. \\
Origen en las bases & Bases técnicas del caso Minería, cap. 4.9; Bases técnicas del caso Minería, cap. 9.6; Bases técnicas del caso Minería, cap. 12; Bases técnicas del caso Minería, criterio de aceptación 12; requisito técnico RT-05.23; supuesto SUP-28 \\
\end{tablasintitulo}

\begin{tablasintitulo}{I{0.200}I{0.800}}{\cab{Campo} & \cab{RF-CUM-016}}
Problema & Problema 3 \\
Módulo & Cumplimiento Ambiental \\
Etapa & Etapa 2 \\
Requerimiento & El sistema debe guardar con versiones los factores de emisión (cuánto se emite por cada unidad de combustible o energía) y la evidencia de cada dato de actividad (consumo o cantidad medida), de modo que el cálculo de un período pasado pueda reproducirse exactamente. \\
Actor & Sustentabilidad y asuntos regulatorios \\
Resultado esperado & El cálculo de cualquier período pasado se reproduce exactamente con los factores y datos de ese período. \\
Origen en las bases & Bases técnicas del caso Minería, cap. 4.9; Bases técnicas del caso Minería, cap. 9.6; Bases técnicas del caso Minería, cap. 12; Bases técnicas del caso Minería, criterio de aceptación 12 \\
\end{tablasintitulo}

\begin{tablasintitulo}{I{0.200}I{0.800}}{\cab{Campo} & \cab{RF-CUM-017}}
Problema & Problema 3 \\
Módulo & Cumplimiento Ambiental \\
Etapa & Etapa 2 \\
Requerimiento & El sistema debe generar el expediente de verificación por tercero independiente con los datos de actividad, los factores aplicados, la memoria de cálculo por alcance y la trazabilidad de cada cifra hasta su medición de origen, en formato abierto y documentado. \\
Actor & Verificador externo de la huella de carbono \\
Resultado esperado & El expediente de verificación se entrega en formato abierto con datos, factores, memoria de cálculo y trazabilidad. \\
Origen en las bases & Bases técnicas del caso Minería, cap. 4.9; Bases técnicas del caso Minería, cap. 9.6; Bases técnicas del caso Minería, cap. 12; Bases técnicas del caso Minería, criterio de aceptación 12 \\
\end{tablasintitulo}

\begin{tablasintitulo}{I{0.200}I{0.800}}{\cab{Campo} & \cab{RF-EQU-001}}
Problema & Problema 2 \\
Módulo & Equipos y Mantenimiento \\
Etapa & Etapa 2 \\
Requerimiento & El sistema debe mostrar en una sola vista, para cada equipo, su estado actual, su historial de los últimos 30 días, sus alarmas de telemetría y sus pautas de mantenimiento pendientes, junto con su desempeño comparado con el de los demás equipos del mismo tipo y con su propio historial de 30 días. La vista debe poder consultarse desde terreno y desde el Centro Integrado de Operaciones con la misma información, sin abrir otros sistemas. \\
Actor & Supervisor de mantenimiento \\
Resultado esperado & El supervisor ve estado, historial de los últimos 30 días, alarmas de telemetría y pautas pendientes del equipo en una sola vista, desde terreno o desde el Centro Integrado de Operaciones, con la misma información. Además ve el desempeño del equipo comparado con el de los equipos del mismo tipo y con su propio historial de los últimos 30 días. \\
Origen en las bases & Bases técnicas del caso Minería, cap. 9.3; Bases técnicas del caso Minería, cap. 9.7; Bases técnicas del caso Minería, criterio de aceptación 6; Derivado de la sección 9.3 y del criterio de aceptación 6 (historial reciente y desempeño comparado): la ventana de 30 días y la comparación por tipo de equipo concretan esa vista \\
\end{tablasintitulo}

\begin{tablasintitulo}{I{0.200}I{0.800}}{\cab{Campo} & \cab{RF-EQU-002}}
Problema & Problema 2 \\
Módulo & Equipos y Mantenimiento \\
Etapa & Etapa 2 \\
Requerimiento & El sistema debe integrar la telemetría (alarmas, horas de operación, consumos y códigos de falla) de los portales de los tres fabricantes de equipos y llevarla a un formato común, usando el estándar abierto de intercambio de datos de equipos mineros (ISO 15143-3) en cada fabricante que lo soporte y declarando cuáles no lo soportan. \\
Actor & Portales de telemetría de fabricantes \\
Resultado esperado & Las señales de los tres fabricantes quedan en un modelo de datos común, usando el estándar abierto en los fabricantes que lo soportan, y se declara cuáles fabricantes no lo soportan. \\
Origen en las bases & Bases técnicas del caso Minería, cap. 4.7; Bases técnicas del caso Minería, criterio de aceptación 6; Bases técnicas del caso Minería, decisión pendiente 10; requisito técnico RT-05.23; Bases técnicas del caso Minería, entrevista al Jefe de Mantenimiento \\
\end{tablasintitulo}

\begin{tablasintitulo}{I{0.200}I{0.800}}{\cab{Campo} & \cab{RF-EQU-003}}
Problema & Problema 2 \\
Módulo & Equipos y Mantenimiento \\
Etapa & Etapa 2 \\
Requerimiento & El sistema debe permitir cambiar la fuente de telemetría de un fabricante sin afectar la vista única de equipo, aunque cambie o deba reemplazarse el acceso a sus datos (hoy no está definido quién es dueño de esos datos ni con qué condiciones pueden extraerse). \\
Actor & Portales de telemetría de fabricantes \\
Resultado esperado & El reemplazo de la fuente de telemetría de un fabricante no afecta la vista única de equipo. \\
Origen en las bases & Bases técnicas del caso Minería, cap. 4.7; Bases técnicas del caso Minería, decisión pendiente 10; requisito técnico RT-05.20; supuesto SUP-10 \\
\end{tablasintitulo}

\begin{tablasintitulo}{I{0.200}I{0.800}}{\cab{Campo} & \cab{RF-EQU-004}}
Problema & Problema 2 \\
Módulo & Equipos y Mantenimiento \\
Etapa & Etapa 2 \\
Requerimiento & El sistema debe exponer las pautas de mantenimiento y las órdenes de trabajo pendientes de cada equipo tomándolas de SAP, sin replicar la gestión que este realiza ni modificar SAP, con una actualización desde SAP no mayor a 15 minutos. \\
Actor & Supervisor de mantenimiento \\
Resultado esperado & Las pautas y órdenes pendientes se muestran en la vista del equipo, actualizadas desde SAP con una antigüedad no mayor a 15 minutos y sin duplicar la gestión de SAP. \\
Origen en las bases & Bases técnicas del caso Minería, cap. 4.7; Bases técnicas del caso Minería, restricción no negociable 4; Bases técnicas del caso Minería, criterio de aceptación 6; Derivado de la sección 4.7 y del criterio de aceptación 6: la frescura de 15 minutos concreta la lectura de SAP \\
\end{tablasintitulo}

\begin{tablasintitulo}{I{0.200}I{0.800}}{\cab{Campo} & \cab{RF-EQU-005}}
Problema & Problema 2 \\
Módulo & Equipos y Mantenimiento \\
Etapa & Etapa 2 \\
Requerimiento & El sistema debe permitir levantar un aviso de falla desde la tablet, en no más de 3 acciones en pantalla y 30 segundos, sin que un planificador lo transcriba desde la radio, y entregarlo a SAP con equipo, componente declarado, hora y persona que lo registra. SAP genera la orden de trabajo; el sistema no. \\
Actor & Supervisor de mantenimiento \\
Resultado esperado & El aviso de falla queda creado en no más de 3 acciones y 30 segundos (validado en la marcha blanca) y llega a SAP sin transcripción. \\
Origen en las bases & Bases técnicas del caso Minería, cap. 4.7; Bases técnicas del caso Minería, cap. 9.4; Bases técnicas del caso Minería, restricción no negociable 8; Bases técnicas del caso Minería, restricción no negociable 4; Derivado de la sección 9.4 y de la restricción no negociable 8 (justificar en segundos todo registro adicional al operador): el máximo de 3 acciones y 30 segundos concreta ese límite \\
\end{tablasintitulo}

\begin{tablasintitulo}{I{0.200}I{0.800}}{\cab{Campo} & \cab{RF-EQU-006}}
Problema & Problema 2 \\
Módulo & Equipos y Mantenimiento \\
Etapa & Etapa 2 \\
Requerimiento & El sistema debe conservar el aviso de falla levantado en zona sin cobertura, con integridad y sin pérdida de datos, durante una autonomía mínima de 48 horas continuas sin enlace, y transmitirlo al recuperar señal, sin que la persona que lo registró deba repetirlo. \\
Actor & Supervisor de mantenimiento \\
Resultado esperado & El aviso de falla queda guardado localmente hasta 48 horas sin enlace y se transmite al recuperar señal, sin repetirlo y sin pérdida. \\
Origen en las bases & Bases técnicas del caso Minería, cap. 4.7; Bases técnicas del caso Minería, restricción no negociable 2; Bases técnicas del caso Minería, criterio de aceptación 13; Bases técnicas del caso Minería, decisión pendiente 9; requisito técnico RT-03.10; requisito técnico RT-03.11; supuesto SUP-09 \\
\end{tablasintitulo}

\begin{tablasintitulo}{I{0.200}I{0.800}}{\cab{Campo} & \cab{RF-CNX-001}}
Problema & Transversal \\
Módulo & Integrador de Sistemas y Nodo de Faena \\
Etapa & Etapa 1 \\
Requerimiento & El sistema debe sincronizar automáticamente, sin intervención manual, los registros hechos en faena durante un corte de enlace, cuando este se restablece. Si un mismo dato quedó distinto en faena y en la nube, debe resolverlo con la regla declarada para su tipo de registro, dejar en una bitácora cada conflicto con la regla aplicada y el resultado, y escalar al equipo de TI del Cliente los conflictos de un tipo de registro que no tenga regla declarada, sin detener la sincronización de los demás registros. \\
Actor & TI del Cliente \\
Resultado esperado & Los registros acumulados durante el corte quedan sincronizados sin intervención manual, con cada conflicto resuelto por regla documentada o escalado al equipo de TI del Cliente, y su traza en la bitácora auditable. \\
Origen en las bases & Bases técnicas del caso Minería, cap. 6; Bases técnicas del caso Minería, cap. 9.8; Bases técnicas del caso Minería, criterio de aceptación 13; requisito técnico RT-03.12; requisito técnico RT-03.11 \\
\end{tablasintitulo}

\begin{tablasintitulo}{I{0.200}I{0.800}}{\cab{Campo} & \cab{RF-CNX-002}}
Problema & Transversal \\
Módulo & Integrador de Sistemas y Nodo de Faena \\
Etapa & Etapa 1 \\
Requerimiento & Sin enlace, el sistema debe validar la credencial y las autorizaciones de una persona con la última lista de habilitaciones vigentes guardada en faena, aunque no se haya recuperado el enlace. \\
Actor & Guardia de portería \\
Resultado esperado & La credencial se valida contra la lista guardada en faena mientras no hay enlace. \\
Origen en las bases & Bases técnicas del caso Minería, cap. 6; Bases técnicas del caso Minería, cap. 9.8; Bases técnicas del caso Minería, restricción no negociable 2; requisito técnico RT-17.01; Bases técnicas del caso Minería, decisión pendiente 9 \\
\end{tablasintitulo}

\begin{tablasintitulo}{I{0.200}I{0.800}}{\cab{Campo} & \cab{RF-CNX-003}}
Problema & Transversal \\
Módulo & Integrador de Sistemas / Nodo de Faena \\
Etapa & Etapa 1 \\
Requerimiento & El dispositivo de terreno debe operar sin enlace desde el origen: guardar cada registro cifrado y enviarlo sin duplicarlo al recuperar cobertura, aplicando la regla de conflictos declarada para cada tipo de dato. Cada hecho debe llevar una hora que no dependa del enlace ni pueda cambiar la persona usuaria, y el relevo de turno debe poder cerrar y abrir sesión localmente. \\
Actor & Persona usuaria \\
Resultado esperado & Cada registro queda guardado cifrado en el dispositivo con una hora independiente del enlace y no modificable, y se transmite al recuperar cobertura sin duplicados; el relevo cierra y abre sesión localmente y los registros posteriores quedan a nombre de quien abrió sesión. \\
Origen en las bases & Bases técnicas del caso Minería, cap. 9.8; Bases técnicas del caso Minería, cap. 6; Bases técnicas del caso Minería, restricción no negociable 2; Bases técnicas del caso Minería, restricción no negociable 9; Bases técnicas del caso Minería, criterio de aceptación 13; Bases técnicas del caso Minería, decisión pendiente 9; requisito técnico RT-03.11; requisito técnico RT-17.05; requisito técnico RT-03.12; requisito técnico RT-17.01; requisito técnico RT-05.03; supuesto SUP-09; supuesto SUP-25 \\
\end{tablasintitulo}

\begin{tablasintitulo}{I{0.200}I{0.800}}{\cab{Campo} & \cab{RF-CNX-004}}
Problema & Transversal \\
Módulo & Integrador de Sistemas / Nodo de Faena \\
Etapa & Etapa 1 \\
Requerimiento & El sistema debe administrar de forma remota y centralizada los dispositivos de terreno y los gabinetes de borde (equipos locales instalados cerca de cada sitio): inventario, configuración, actualización de la aplicación y del firmware y bloqueo y borrado remoto ante pérdida o desvinculación de la persona usuaria, con la información almacenada cifrada. El bloqueo debe aplicarse en no más de 15 minutos desde la orden, o desde que el dispositivo recupere el enlace si estaba sin conexión. \\
Actor & TI del Cliente \\
Resultado esperado & El inventario muestra el estado de cada dispositivo, y una actualización, cambio de configuración, bloqueo o borrado remoto queda aplicado desde la consola, el bloqueo dentro de 15 minutos desde la orden (o desde que el dispositivo recupera conexión). \\
Origen en las bases & requisito técnico RT-03.18; requisito técnico RT-17.05; supuesto SUP-33 \\
\end{tablasintitulo}

\begin{tablasintitulo}{I{0.200}I{0.800}}{\cab{Campo} & \cab{RF-ADM-001}}
Problema & Transversal \\
Módulo & Administración y Servicios Comunes \\
Etapa & Etapa 1 \\
Requerimiento & El sistema debe mantener accesibles, en un repositorio de consulta de solo lectura, los datos históricos que no se migren, durante el plazo de retención que fija el Caso para cada tipo de registro (RT-05.10). \\
Actor & Cumplimiento y auditoría interna \\
Resultado esperado & El registro histórico no migrado se consulta en solo lectura, con el mismo control de acceso y la misma auditoría que el registro vivo, durante el plazo de retención de RT-05.10 para su tipo de registro. \\
Origen en las bases & requisito técnico RT-05.15; requisito técnico RT-05.10; requisito técnico RT-16.10; supuesto SUP-30 \\
\end{tablasintitulo}

\begin{tablasintitulo}{I{0.200}I{0.800}}{\cab{Campo} & \cab{RF-ADM-002}}
Problema & Transversal \\
Módulo & Administración y Servicios Comunes \\
Etapa & Etapa 1 \\
Requerimiento & El sistema debe permitir al Cliente, sin intervención del Adjudicatario, administrar personas usuarias, roles, permisos, unidades organizacionales y sus jerarquías. \\
Actor & TI del Cliente \\
Resultado esperado & Altas, modificaciones y bajas de personas usuarias, roles y unidades quedan ejecutadas por el Cliente y registradas en la auditoría. \\
Origen en las bases & Bases técnicas del caso Minería, restricción no negociable 10; requisito técnico RT-16.01 \\
\end{tablasintitulo}

\begin{tablasintitulo}{I{0.200}I{0.800}}{\cab{Campo} & \cab{RF-ADM-003}}
Problema & Transversal \\
Módulo & Administración y Servicios Comunes \\
Etapa & Etapa 1 \\
Requerimiento & El sistema debe permitir configurar desde la interfaz de administración las reglas de negocio parametrizables -umbrales, plazos, tolerancias, montos, catálogos, listas de valores y textos de notificación-, con control de versiones y registro de quién cambió qué y cuándo, y el Proponente debe declarar en un catálogo de parámetros qué elementos son parametrizables y cuáles requieren desarrollo. \\
Actor & Administrador funcional del Cliente \\
Resultado esperado & El nuevo valor rige desde su fecha de vigencia sin despliegue de código, la versión anterior queda consultable con autor y fecha, y el catálogo declara para cada elemento si es parametrizable o requiere desarrollo, coincidiendo con la demostración. \\
Origen en las bases & requisito técnico RT-16.02; requisito técnico RT-16.04 \\
\end{tablasintitulo}

\begin{tablasintitulo}{I{0.200}I{0.800}}{\cab{Campo} & \cab{RF-ADM-004}}
Problema & Transversal \\
Módulo & Administración y Servicios Comunes \\
Etapa & Etapa 1 \\
Requerimiento & El sistema debe exigir la aprobación de una segunda persona (otro perfil) para todo cambio de parámetro con impacto operacional y registrar la justificación del cambio. \\
Actor & Administrador funcional del Cliente \\
Resultado esperado & El cambio no entra en vigencia sin la aprobación del segundo perfil y queda registrado con su justificación. \\
Origen en las bases & requisito técnico RT-16.03 \\
\end{tablasintitulo}

\begin{tablasintitulo}{I{0.200}I{0.800}}{\cab{Campo} & \cab{RF-ADM-007}}
Problema & Transversal \\
Módulo & Administración y Servicios Comunes \\
Etapa & Etapa 1 \\
Requerimiento & El sistema debe registrar el ciclo de vida completo de cada cuenta de usuario (creación, modificación, elevación de privilegios, bloqueo y baja) de modo que nadie pueda negar ni alterar lo registrado, y conservarlo 10 años. \\
Actor & TI del Cliente \\
Resultado esperado & El historial de cada cuenta se consulta completo, con autor y fecha de cada evento, con no repudio y conservado por el período de retención de auditoría. \\
Origen en las bases & Bases técnicas del caso Minería, restricción no negociable 9; requisito técnico RT-12.09; requisito técnico RT-16.10; requisito técnico RT-05.10 \\
\end{tablasintitulo}

\begin{tablasintitulo}{I{0.200}I{0.800}}{\cab{Campo} & \cab{RF-ADM-008}}
Problema & Transversal \\
Módulo & Administración y Servicios Comunes \\
Etapa & Etapa 1 \\
Requerimiento & El sistema debe soportar flujos de trabajo con estados, transiciones, responsables, plazos, escalamiento automático por vencimiento y delegación por ausencia, configurables por el Cliente sin desarrollo, al menos en lo relativo a responsables, plazos y niveles de aprobación. \\
Actor & Administrador funcional del Cliente \\
Resultado esperado & Un cambio de responsable, plazo o nivel de aprobación queda aplicado sin desarrollo, las tareas vencidas escalan en forma automática y la delegación por ausencia reasigna la tarea al delegado. \\
Origen en las bases & requisito técnico RT-16.11; requisito técnico RT-16.12 \\
\end{tablasintitulo}

\begin{tablasintitulo}{I{0.200}I{0.800}}{\cab{Campo} & \cab{RF-DAT-003}}
Problema & Transversal \\
Módulo & Administración y Servicios Comunes \\
Etapa & Etapa 1 \\
Requerimiento & El sistema debe permitir la carga y descarga masiva de información en formatos abiertos, con validación previa, informe de errores por registro y procesamiento parcial. \\
Actor & Administrador funcional del Cliente \\
Resultado esperado & Los registros válidos quedan cargados y los inválidos quedan informados uno a uno. \\
Origen en las bases & requisito técnico RT-05.22 \\
\end{tablasintitulo}

\begin{tablasintitulo}{I{0.200}I{0.800}}{\cab{Campo} & \cab{RF-ADM-009}}
Problema & Transversal \\
Módulo & Administración y Servicios Comunes \\
Etapa & Etapa 1 \\
Requerimiento & El sistema debe mostrar a cada persona usuaria, en una bandeja de tareas unificada, todas las solicitudes pendientes que tiene asignadas, con priorización y alerta de vencimiento. \\
Actor & Persona usuaria \\
Resultado esperado & La persona ve en una sola bandeja todas sus tareas pendientes ordenadas por prioridad y vencimiento. \\
Origen en las bases & requisito técnico RT-16.13 \\
\end{tablasintitulo}

\begin{tablasintitulo}{I{0.200}I{0.800}}{\cab{Campo} & \cab{RF-ADM-012}}
Problema & Transversal \\
Módulo & Administración y Servicios Comunes \\
Etapa & Etapa 1 \\
Requerimiento & El sistema debe ofrecer una búsqueda global por el contenido de los registros, que encuentre resultados aunque haya errores de tipeo y permita filtrar por categorías, mostrando sólo los registros que el rol de la persona puede ver. \\
Actor & Persona usuaria \\
Resultado esperado & La búsqueda devuelve resultados que incluyen sólo registros permitidos para el rol de la persona. \\
Origen en las bases & requisito técnico RT-16.27 \\
\end{tablasintitulo}

\begin{tablasintitulo}{I{0.200}I{0.800}}{\cab{Campo} & \cab{RF-ADM-013}}
Problema & Transversal \\
Módulo & Administración y Servicios Comunes \\
Etapa & Etapa 1 \\
Requerimiento & El sistema debe entregar listados ordenables, filtrables y paginados, exportables en formatos abiertos con el filtro aplicado, y procesar en segundo plano las exportaciones de gran volumen (más de 10.000 registros o más de 50 MB), avisando cuando terminen y sin bloquear la sesión. \\
Actor & Persona usuaria \\
Resultado esperado & El archivo exportado refleja el filtro aplicado, y una exportación sobre el umbral no bloquea la sesión y notifica a la persona usuaria al completarse. \\
Origen en las bases & requisito técnico RT-16.28; requisito técnico RT-16.29; Derivado de RT-16.28 y RT-16.29 (exportación asíncrona de gran volumen): el umbral de 10.000 registros o 50 MB define qué es gran volumen \\
\end{tablasintitulo}

\begin{tablasintitulo}{I{0.200}I{0.800}}{\cab{Campo} & \cab{RF-NOT-002}}
Problema & Transversal \\
Módulo & Administración y Servicios Comunes \\
Etapa & Etapa 1 \\
Requerimiento & El sistema debe permitir a cada persona elegir el canal y la frecuencia con que recibe notificaciones, respetando las que el Cliente defina como obligatorias (operacionales, de seguridad y de emergencia), y permitir darse de baja de las informativas, promocionales y de encuestas o boletines, como exige la normativa (artículo 28 B de la Ley 19.496 y Ley 21.719 de protección de datos personales). \\
Actor & Persona usuaria \\
Resultado esperado & Las notificaciones llegan según las preferencias de la persona, salvo las que el Cliente definió como obligatorias, y la persona puede darse de baja de las que la normativa lo exige, respetándose la baja. \\
Origen en las bases & requisito técnico RT-16.22; requisito técnico RT-16.25; Derivado de RT-16.22 y RT-16.25 (preferencias y baja de notificaciones): la normativa citada y los tipos de notificación sujetos a baja concretan esa exigencia \\
\end{tablasintitulo}

\begin{tablasintitulo}{I{0.200}I{0.800}}{\cab{Campo} & \cab{RF-NOT-003}}
Problema & Transversal \\
Módulo & Administración y Servicios Comunes \\
Etapa & Etapa 1 \\
Requerimiento & El sistema debe permitir al Cliente administrar las plantillas de notificación, con variables de la transacción y control de versiones. \\
Actor & Administrador funcional del Cliente \\
Resultado esperado & El texto de la notificación se modifica sin desarrollo y la versión anterior queda consultable. \\
Origen en las bases & requisito técnico RT-16.21; Consulta N.º 17 \\
\end{tablasintitulo}

\begin{tablasintitulo}{I{0.200}I{0.800}}{\cab{Campo} & \cab{RF-ANA-001}}
Problema & Transversal \\
Módulo & Indicadores y Centro Integrado \\
Etapa & Etapa 1 \\
Requerimiento & El sistema debe ofrecer tableros operacionales y de gestión con los indicadores del Caso, filtrables por período, unidad organizacional, equipo, turno, polígono y área, con acceso propio y permanente del Cliente y con exportación. Desde cada indicador debe poder llegarse hasta la transacción de origen. \\
Actor & Gerencia del Cliente \\
Resultado esperado & Desde cada indicador se navega hasta los registros de origen que lo componen con los filtros indicados, y el Cliente accede por cuenta propia y de forma permanente a los tableros, con exportación. \\
Origen en las bases & requisito técnico RT-05.25; requisito técnico RT-05.26; requisito técnico RT-14.02; requisito técnico RT-05.29; Derivado de RT-14.02 (datos en tiempo real): el retraso máximo de 60 segundos define qué se entiende por tiempo real \\
\end{tablasintitulo}

\begin{tablasintitulo}{I{0.200}I{0.800}}{\cab{Campo} & \cab{RF-ANA-002}}
Problema & Transversal \\
Módulo & Indicadores y Centro Integrado \\
Etapa & Etapa 2 \\
Requerimiento & El sistema debe permitir al Cliente construir sus propios informes sin intervención del Adjudicatario, mediante una herramienta de autoservicio con una guía documentada de los datos disponibles. \\
Actor & Administrador funcional del Cliente \\
Resultado esperado & Un informe nuevo queda construido y publicado por el Cliente sin desarrollo. \\
Origen en las bases & requisito técnico RT-05.27 \\
\end{tablasintitulo}

\begin{tablasintitulo}{I{0.200}I{0.800}}{\cab{Campo} & \cab{RF-ANA-003}}
Problema & Transversal \\
Módulo & Indicadores y Centro Integrado \\
Etapa & Etapa 2 \\
Requerimiento & El sistema debe permitir exportar todo informe en formatos abiertos y programar su envío automático por calendario. \\
Actor & Persona usuaria \\
Resultado esperado & El informe llega a sus destinatarios en la fecha programada, en formato abierto. \\
Origen en las bases & requisito técnico RT-05.28 \\
\end{tablasintitulo}

\begin{tablasintitulo}{I{0.200}I{0.800}}{\cab{Campo} & \cab{RF-TER-009}}
Problema & Problema 2 \\
Módulo & Operación Mina \\
Etapa & Etapa 1 \\
Requerimiento & El sistema debe calcular, a partir de los parámetros de perforación de cada pozo, un índice de dureza del macizo rocoso y ponerlo a disposición de geología y de planificación minera, asociado al pozo y al sector, en formato estructurado y exportable, sin trabajo de conversión previo. \\
Actor & Geología \\
Resultado esperado & Geología y planificación consultan y descargan el índice de dureza por pozo y por sector sin conversión previa. \\
Origen en las bases & Bases técnicas del caso Minería, cap. 4.2; Bases técnicas del caso Minería, cap. 4.2 (los datos de perforación contienen información sobre la dureza del macizo que hoy no se aprovecha; geología los pidió tres veces) \\
\end{tablasintitulo}

\begin{tablasintitulo}{I{0.200}I{0.800}}{\cab{Campo} & \cab{RF-EQU-007}}
Problema & Problema 2 \\
Módulo & Equipos y Mantenimiento \\
Etapa & Etapa 2 \\
Requerimiento & El sistema debe registrar el tiempo transcurrido entre la detención de un equipo de mina y la determinación de su causa, y publicar el tiempo medio entre falla y diagnóstico por equipo, por tipo de falla y por turno, para contrastarlo con la meta del Cliente de 1 hora (situación actual: 4,2 horas). \\
Actor & Supervisor de mantenimiento \\
Resultado esperado & El tiempo entre la detención y el diagnóstico queda registrado por evento y se publica el tiempo medio por equipo, tipo de falla y turno, contrastado con la meta de 1 hora. \\
Origen en las bases & Bases técnicas del caso Minería, cap. 9.3; Bases técnicas del caso Minería, cap. 7.2 (indicador "Tiempo medio entre falla y diagnóstico en equipo mina": 4,2 horas, meta interna 1 hora) \\
\end{tablasintitulo}

\begin{tablasintitulo}{I{0.200}I{0.800}}{\cab{Campo} & \cab{RF-PLA-010}}
Problema & Problema 2 \\
Módulo & Operación Mina \\
Etapa & Etapa 1 \\
Requerimiento & El sistema debe medir la pérdida efectiva de tiempo en el cambio de turno, con el desglose de sus etapas, a partir de eventos que el sistema ya registra y sin agregar pasos al ciclo productivo, y publicarla por turno y por área para contrastarla con la meta del Cliente de 25 minutos por turno (situación actual: 47 minutos). \\
Actor & Supervisor de turno \\
Resultado esperado & La pérdida por cambio de turno se publica por turno y por área con su desglose, contrastada con la meta de 25 minutos, sin que el operador registre nada adicional. \\
Origen en las bases & Bases técnicas del caso Minería, cap. 7.2; Bases técnicas del caso Minería, cap. 7.2 (indicador "Pérdida efectiva por cambio de turno": 47 minutos por turno, meta interna 25); Bases técnicas del caso Minería, restricción no negociable 8. Desglose por etapas: criterio del proponente, el Caso no lo exige \\
\end{tablasintitulo}

\begin{tablasintitulo}{I{0.200}I{0.800}}{\cab{Campo} & \cab{RF-CON-016}}
Problema & Problema 1 \\
Módulo & Conciliación Metalúrgica \\
Etapa & Etapa 1 \\
Requerimiento & El sistema debe registrar el tiempo transcurrido entre la toma de cada muestra y la disponibilidad de su resultado, tomando las marcas de tiempo del sistema de gestión de laboratorio, y alertar cuando exceda el plazo definido para el tipo de análisis (hoy entre 4 y 36 horas según el análisis). \\
Actor & Laboratorio químico \\
Resultado esperado & Cada muestra queda con su tiempo hasta resultado y se alerta la que excede el plazo definido para su tipo de análisis. \\
Origen en las bases & Bases técnicas del caso Minería, cap. 4.5; Bases técnicas del caso Minería, cap. 4.5 (tiempo entre la toma de la muestra y la disponibilidad del resultado de entre 4 y 36 horas); Bases técnicas del caso Minería, restricción no negociable 6. Plazo por tipo de análisis: parametrizable, a definir con el laboratorio \\
\end{tablasintitulo}

\begin{tablasintitulo}{I{0.200}I{0.800}}{\cab{Campo} & \cab{RF-CON-017}}
Problema & Problema 1 \\
Módulo & Conciliación Metalúrgica \\
Etapa & Etapa 1 \\
Requerimiento & El sistema debe enviar al sistema de gestión de laboratorio la identificación de cada muestra tomada en los puntos de transferencia de la planta (identificador de la muestra, punto de toma, fecha y hora y lote de turno al que corresponde), para que el laboratorio asocie el resultado del análisis a la muestra correcta. \\
Actor & Metalurgia y planta \\
Resultado esperado & El sistema de gestión de laboratorio recibe la identificación de cada muestra antes de que se emita su resultado, sin transcripción manual. \\
Origen en las bases & Bases técnicas del caso Minería, cap. 4.5; Bases técnicas del caso Minería, restricción no negociable 6; supuesto SUP-12; supuesto SUP-16 \\
\end{tablasintitulo}

\begin{tablasintitulo}{I{0.200}I{0.800}}{\cab{Campo} & \cab{RF-CNX-005}}
Problema & Transversal \\
Módulo & Integrador de Sistemas / Nodo de Faena \\
Etapa & Etapa 1 \\
Requerimiento & El sistema debe registrar cada corte del enlace con Antofagasta, con su inicio, su término, su duración y el volumen de datos acumulado localmente durante el corte, y mantener ese historial disponible para el equipo de TI del Cliente. \\
Actor & TI del Cliente \\
Resultado esperado & Cada corte queda registrado con inicio, término, duración y volumen acumulado, y el historial es consultable por el equipo de TI. \\
Origen en las bases & Bases técnicas del caso Minería, cap. 6 (entre tres y seis cortes de fibra al año, de cuatro a veinte horas); Bases técnicas del caso Minería, cap. 9.8; Bases técnicas del caso Minería, criterio de aceptación 13; Bases técnicas del caso Minería, restricción no negociable 2; Derivado de estos antecedentes: el registro de cada corte permite al equipo de TI demostrar el cumplimiento de las 48 horas de operación sin enlace \\
\end{tablasintitulo}

\begin{tablasintitulo}{I{0.200}I{0.800}}{\cab{Campo} & \cab{RF-TRZ-019}}
Problema & Problema 1 \\
Módulo & Trazabilidad y Embarque \\
Etapa & Etapa 1 \\
Requerimiento & El sistema debe registrar la liquidación que informa el comprador de cada embarque (el pago final, que se calcula meses después) con sus componentes: ley pagable, deducciones por humedad, cargos de tratamiento y refinación y penalidades por elementos deletéreos. Cada componente se vincula al lote comercial y a sus polígonos de origen, y los datos se entregan a SAP. El sistema no calcula la liquidación ni emite cifras contables. \\
Actor & Área comercial \\
Resultado esperado & La liquidación informada por el comprador queda registrada por componente, vinculada al lote y a sus polígonos de origen, y sus datos se entregan a SAP, sin que el sistema la calcule. \\
Origen en las bases & Bases técnicas del caso Minería, cap. 4.6; Bases técnicas del caso Minería, cap. 4.6 (la liquidación se realiza meses después del embarque); Bases técnicas del caso Minería, restricción no negociable 4 (SAP emite)". \\
\end{tablasintitulo}

\begin{tablasintitulo}{I{0.200}I{0.800}}{\cab{Campo} & \cab{RF-TRZ-020}}
Problema & Problema 1 \\
Módulo & Trazabilidad y Embarque \\
Etapa & Etapa 1 \\
Requerimiento & El sistema debe calcular, para cada lote comercial en conformación, el contenido de arsénico y de los demás elementos deletéreos a partir de la composición de origen por polígono y de los análisis disponibles, y alertar antes del embarque cuando el lote proyecte superar el límite del contrato del comprador (límite parametrizable por contrato). \\
Actor & Área comercial \\
Resultado esperado & El lote muestra su contenido proyectado de elementos deletéreos y se alerta antes del embarque si supera el límite del contrato del comprador. \\
Origen en las bases & Bases técnicas del caso Minería, cap. 9.1; Bases técnicas del caso Minería, cap. 1 (4 de 31 embarques con penalización por arsénico, USD 1,9 millones); Bases técnicas del caso Minería, cap. 7.1. La alerta preventiva antes del embarque es propuesta del proponente, el Caso no la exige expresamente \\
\end{tablasintitulo}

\begin{tablasintitulo}{I{0.200}I{0.800}}{\cab{Campo} & \cab{RF-EQU-008}}
Problema & Problema 2 \\
Módulo & Equipos y Mantenimiento \\
Etapa & Etapa 2 \\
Requerimiento & El sistema debe generar alertas de mantenimiento basadas en reglas definidas por el propio equipo de mantenimiento, con un límite semanal de alertas por equipo, y registrar cuántas de ellas terminaron en una acción de mantenimiento. \\
Actor & Jefe de mantenimiento / Supervisor de mantenimiento \\
Resultado esperado & Las alertas del equipo quedan generadas por regla, respetando el límite semanal, con registro de cuántas derivaron en una acción de mantenimiento. \\
Origen en las bases & Bases técnicas del caso Minería, cap. 8, entrevista a Bustos; Bases técnicas del caso Minería, cap. 9.3. Límite semanal y reglas propias: decisión del proponente para evitar la sobrecarga de alertas que causó el fracaso de las iniciativas anteriores. \\
\end{tablasintitulo}

\seccion{Requerimientos no funcionales}
Los requerimientos no funcionales fijan cómo debe operar la solución, con un umbral que se puede medir. Cada ficha indica el método con que se verifica su cumplimiento.

\begin{tablasintitulo}{I{0.200}I{0.800}}{\cab{Campo} & \cab{RNF-CAP-001}}
Problema & Problema 1 \\
Requerimiento & El sistema debe poder crecer hasta 3 veces la actividad actual de movimiento de material sin rediseñar su arquitectura: 1.026.000 viajes de camión al año, 4,11 millones de ciclos de carguío (paladas) al año y 54 polígonos activos a la vez. La proyección del Caso a 3 años (367.000 viajes y 22 polígonos) queda dentro de ese margen. \\
Método de verificación & Prueba de carga a 1,5 veces el peak declarado (momento de mayor actividad, RT-09.06) y demostración en Preproducción del escalamiento hasta 3 veces la volumetría actual sin cambio de arquitectura. Prueba de carga con el perfil de concurrencia derivado del numeral 14.1 y declarado conforme al numeral 14.2. \\
Origen en las bases & Bases técnicas del caso Minería, cap. 9.1; Bases técnicas del caso Minería, cap. 14.1; requisito técnico RT-09.03 \\
\end{tablasintitulo}

\begin{tablasintitulo}{I{0.200}I{0.800}}{\cab{Campo} & \cab{RNF-CAP-002}}
Problema & Problema 1 \\
Requerimiento & El sistema debe absorber hasta 3 veces el volumen actual de muestras y resultados de laboratorio, es decir 213.000 muestras analizadas al año, sin degradar la disponibilidad de la conciliación. \\
Método de verificación & Prueba de carga a 1,5 veces el peak declarado (RT-09.06) y demostración en Preproducción del escalamiento hasta 3 veces la volumetría actual sin cambio de arquitectura. Prueba de carga sobre la integración con el sistema de gestión de laboratorio en el peak de recepción de resultados. \\
Origen en las bases & Bases técnicas del caso Minería, cap. 4.5; Bases técnicas del caso Minería, cap. 14.1; requisito técnico RT-09.03 \\
\end{tablasintitulo}

\begin{tablasintitulo}{I{0.200}I{0.800}}{\cab{Campo} & \cab{RNF-CAP-003}}
Problema & Problema 3 \\
Requerimiento & El sistema debe absorber hasta 3 veces el volumen actual de control de acceso, es decir 8.700 ingresos y salidas diarios, 25.500 personas con acceso durante una detención mayor de planta y 270 empresas contratistas activas. \\
Método de verificación & Prueba de carga a 1,5 veces el peak declarado (RT-09.06) y demostración en Preproducción del escalamiento hasta 3 veces la volumetría actual sin cambio de arquitectura. Prueba de carga en el peak de cambio de turno y simulación del escenario de detención mayor de planta. \\
Origen en las bases & Bases técnicas del caso Minería, cap. 4.8; Bases técnicas del caso Minería, cap. 9.5; Bases técnicas del caso Minería, cap. 14.1; requisito técnico RT-09.03 \\
\end{tablasintitulo}

\begin{tablasintitulo}{I{0.200}I{0.800}}{\cab{Campo} & \cab{RNF-CAP-004}}
Problema & Problema 2 \\
Requerimiento & El sistema debe poder crecer hasta 3 veces la telemetría y las señales de proceso actuales: 156 equipos con telemetría (3 × 52), 15.000 señales de proceso de ingesta continua (3 × las 5.000 que se ingestan, no las 14.500 declaradas por el historiador) y 558 instrumentos de monitoreo (3 × 186), sin superar los 5 segundos de la vista única de equipo. \\
Método de verificación & Prueba de carga a 1,5 veces el peak declarado (RT-09.06) y demostración en Preproducción del escalamiento hasta 3 veces la volumetría actual sin cambio de arquitectura. Prueba de carga con el volumen proyectado y medición de la latencia de la vista de equipo. \\
Origen en las bases & Bases técnicas del caso Minería, cap. 4.7; Bases técnicas del caso Minería, cap. 14.1; requisito técnico RT-09.03; supuesto SUP-29 \\
\end{tablasintitulo}

\begin{tablasintitulo}{I{0.200}I{0.800}}{\cab{Campo} & \cab{RNF-CAP-005}}
Problema & Problema 3 \\
Requerimiento & El sistema debe absorber hasta 3 veces el volumen actual de registros de seguridad, es decir 34.500 incidentes, cuasi accidentes y observaciones al año, 23.400 inspecciones planificadas al año y 13.800 alertas de fatiga al año. \\
Método de verificación & Prueba de carga a 1,5 veces el peak declarado (RT-09.06) y demostración en Preproducción del escalamiento hasta 3 veces la volumetría actual sin cambio de arquitectura. Prueba de carga con el volumen proyectado y verificación del tiempo de generación del expediente de fiscalización. \\
Origen en las bases & Bases técnicas del caso Minería, cap. 4.8; Bases técnicas del caso Minería, cap. 14.1; requisito técnico RT-09.03 \\
\end{tablasintitulo}

\begin{tablasintitulo}{I{0.200}I{0.800}}{\cab{Campo} & \cab{RNF-CAP-006}}
Problema & Transversal \\
Requerimiento & El sistema debe almacenar localmente, en faena y en los dispositivos de terreno, hasta 150.000 eventos acumulados durante un corte de enlace de 48 horas (en la cola persistente del Nodo de Faena; cada gabinete de borde conserva su cola local), sin rechazar registros. \\
Método de verificación & Corte controlado de 48 horas con medición del volumen acumulado por nodo y del margen restante. \\
Origen en las bases & Bases técnicas del caso Minería, cap. 6; Bases técnicas del caso Minería, cap. 9.8; requisito técnico RT-03.10; Derivado de las 48 horas sin enlace (RT-03.10) y de la cantidad de nodos de borde: los 150.000 eventos cubren los 146.016 estimados en ese corte con tres veces la volumetría (RT-09.03), \\
\end{tablasintitulo}

\begin{tablasintitulo}{I{0.200}I{0.800}}{\cab{Campo} & \cab{RNF-CAP-007}}
Problema & Transversal \\
Requerimiento & El sistema debe sostener 40 transacciones por segundo en régimen normal y 200 en el momento de mayor carga, y en ese momento el tiempo de respuesta no debe aumentar más de 20 \% respecto del régimen normal. \\
Método de verificación & Prueba de carga con el perfil derivado del numeral 14.1 y declarado conforme al numeral 14.2. \\
Origen en las bases & Bases técnicas del caso Minería, cap. 14.2; requisito técnico RT-09.02; Derivado de la sección 14.2 (volumetría que el proponente debe estimar): los valores de 40 y 200 transacciones por segundo son esa estimación \\
\end{tablasintitulo}

\begin{tablasintitulo}{I{0.200}I{0.800}}{\cab{Campo} & \cab{RNF-CAP-008}}
Problema & Transversal \\
Requerimiento & El sistema debe escalar horizontal y automáticamente sus capas de aplicación e integración, sin intervención manual y con 0 transacciones perdidas durante el escalamiento. \\
Método de verificación & Prueba de carga creciente en Preproducción hasta gatillar el escalamiento, verificando las transacciones durante el evento. \\
Origen en las bases & requisito técnico RT-02.10; requisito técnico RT-09.04 \\
\end{tablasintitulo}

\begin{tablasintitulo}{I{0.200}I{0.800}}{\cab{Campo} & \cab{RNF-CUM-001}}
Problema & Problema 3 \\
Requerimiento & El tratamiento de los datos personales sensibles (salud ocupacional, control de fatiga, control de alcohol y drogas, y datos biométricos) debe cumplir la Ley N° 21.719 y el artículo 85 de las Bases administrativas: el 100 \% de los registros de esas categorías con la base legal que autoriza su tratamiento documentada y trazable, y 0 registros con plazo de retención vencido sin eliminar o anonimizar al momento de la auditoría. \\
Método de verificación & Revisión con el área legal del Cliente, sobre cada categoría de dato, de que el 100 \% de los registros tenga su base legal de tratamiento documentada y trazable, más prueba de eliminación o anonimización verificable sobre una muestra de registros con plazo de retención vencido. \\
Origen en las bases & Bases técnicas del caso Minería, cap. 4.8; Bases técnicas del caso Minería, cap. 9.5; Bases técnicas del caso Minería, cap. 12; requisito técnico RT-05.10 de las Bases técnicas del caso Minería, equivalente al RT-16.10 de las Bases técnicas transversales; requisito técnico RT-05.07; requisito técnico RT-05.08; Bases administrativas, artículo 85; Ley N° 21.719; El registro del consentimiento deriva de la Ley 21.719 (sección 12); la eliminación segura se apoya en RT-05.07 y en el artículo 85 de las Bases administrativas \\
\end{tablasintitulo}

\begin{tablasintitulo}{I{0.200}I{0.800}}{\cab{Campo} & \cab{RNF-CUM-002}}
Problema & Problema 3 \\
Requerimiento & El sistema debe firmar con firma electrónica avanzada (Ley 19.799) los registros de seguridad exigidos por la autoridad, las actas de aceptación de embarque y los informes ambientales que se envían a la autoridad, comprobando que el certificado sea válido al firmar, agregando un sello de tiempo y guardando la evidencia, para que cada documento pueda verificarse durante todo su plazo de retención, incluso con el certificado ya vencido. \\
Método de verificación & Verificación de la firma avanzada y de la validez del certificado al emitirse, y de la firma y del sello de tiempo sobre una muestra de documentos de cada tipo, incluidos documentos cuyo certificado de firma ya esté vencido (real o de prueba con vencimiento simulado), dentro del plazo de retención. \\
Origen en las bases & Bases técnicas del caso Minería, cap. 4.8; requisito técnico RT-05.10 de las Bases técnicas del caso Minería, equivalente al RT-16.10 de las Bases técnicas transversales; requisito técnico RT-16.14; requisito técnico RT-16.17; requisito técnico RT-16.18; Ley N° 19.799; Consulta N.º 11; supuesto SUP-25 \\
\end{tablasintitulo}

\begin{tablasintitulo}{I{0.200}I{0.800}}{\cab{Campo} & \cab{RNF-DES-001}}
Problema & Problema 2 \\
Requerimiento & El sistema debe registrar un movimiento de material desde el equipo en terreno, incluido el aviso de cambio de material en cabina, en no más de 2 segundos en el 95 \% de los casos, medidos desde la acción de la persona operadora hasta la confirmación en pantalla. \\
Método de verificación & Medición instrumentada sobre al menos 1.000 registros reales en operación normal, con el dispositivo definitivo de terreno. La medición incluye el aviso de cambio de material y se realiza durante la marcha blanca (período de prueba en paralelo con el método actual), en turno de día y de noche. \\
Origen en las bases & Bases técnicas del caso Minería, cap. 9.4; Bases técnicas del caso Minería, restricción no negociable 8; requisito técnico RT-09.01 \\
\end{tablasintitulo}

\begin{tablasintitulo}{I{0.200}I{0.800}}{\cab{Campo} & \cab{RNF-DES-002}}
Problema & Problema 1 \\
Requerimiento & El sistema debe completar el registro en la báscula de camiones en un máximo de 3 segundos por pesaje, sin detener la fila de pesaje. \\
Método de verificación & Medición sobre 200 pesajes consecutivos en operación normal. \\
Origen en las bases & Bases técnicas del caso Minería, cap. 4.3; requisito técnico RT-09.01 \\
\end{tablasintitulo}

\begin{tablasintitulo}{I{0.200}I{0.800}}{\cab{Campo} & \cab{RNF-DES-003}}
Problema & Problema 3 \\
Requerimiento & El sistema debe verificar la habilitación de una persona en portería en no más de 1,5 segundos, también cuando no hay enlace (usando la lista de habilitaciones guardada en faena). \\
Método de verificación & Medición sobre 500 ingresos en el peak de cambio de turno, con y sin enlace disponible. \\
Origen en las bases & Bases técnicas del caso Minería, cap. 4.8; Bases técnicas del caso Minería, cap. 9.5; requisito técnico RT-09.01 \\
\end{tablasintitulo}

\begin{tablasintitulo}{I{0.200}I{0.800}}{\cab{Campo} & \cab{RNF-DES-004}}
Problema & Problema 3 \\
Requerimiento & El sistema debe completar, desde la lectura de la credencial de proximidad en cabina hasta el registro de la identificación de la persona operadora, en no más de 2 segundos en el 95 \% de los casos. \\
Método de verificación & Prueba en un cambio de turno de noche en invierno, con guantes de invierno puestos y sin retirarlos, cronometrando el tiempo entre la lectura de la credencial y el registro de la identificación y verificando que no supere 2 segundos en el 95 \% de los intentos. \\
Origen en las bases & Bases técnicas del caso Minería, cap. 4.8; Bases técnicas del caso Minería, cap. 9.5; Bases técnicas del caso Minería, restricción no negociable 8; requisito técnico RT-12.11; supuesto SUP-04 \\
\end{tablasintitulo}

\begin{tablasintitulo}{I{0.200}I{0.800}}{\cab{Campo} & \cab{RNF-DES-005}}
Problema & Problema 1 \\
Requerimiento & El sistema debe completar la reconstrucción del origen de un lote ya embarcado en menos de 2 horas, con cada etapa de la cadena respaldada por evidencia o, en los tres casos en que no existe un registro medido (carga en el borde entre dos polígonos, mezcla en un stock intermedio y material anterior al arranque del sistema), por el supuesto declarado que la sustituye, indicando en pantalla cuál es cuál. \\
Método de verificación & Ejercicio cronometrado sobre tres lotes ya embarcados durante la marcha blanca, midiendo la reconstrucción en menos de 2 horas, y revisión de que cada etapa muestra su evidencia o su supuesto declarado y se distingue una de otro, con verificación por el área comercial. \\
Origen en las bases & Bases técnicas del caso Minería, cap. 9.1; Bases técnicas del caso Minería, criterio de aceptación 1; supuesto SUP-03; supuesto SUP-11; supuesto SUP-21 \\
\end{tablasintitulo}

\begin{tablasintitulo}{I{0.200}I{0.800}}{\cab{Campo} & \cab{RNF-DES-006}}
Problema & Problema 1 \\
Requerimiento & El sistema debe cerrar la conciliación metalúrgica del mes en no más de 3 días hábiles desde el cierre contable del mes. \\
Método de verificación & Medición durante seis cierres consecutivos, contra la línea base de 18 días hábiles. \\
Origen en las bases & Bases técnicas del caso Minería, cap. 9.2; Bases técnicas del caso Minería, criterio de aceptación 2; requisito técnico RT-05.29 \\
\end{tablasintitulo}

\begin{tablasintitulo}{I{0.200}I{0.800}}{\cab{Campo} & \cab{RNF-DES-007}}
Problema & Problema 1 \\
Requerimiento & El sistema debe dejar disponible, todos los días, el costo por tonelada movida y por libra de cobre fino del día operacional anterior (cerrado a las 08:00 con el término del turno de noche), en no más de 4 horas desde que SAP entrega los datos contables del día. \\
Método de verificación & Medición durante 30 días consecutivos, contra la situación actual en que el costo no existe con esa frecuencia y llega 18 días hábiles después del cierre, verificando que cada día esté disponible dentro de las 4 horas desde la entrega de SAP. \\
Origen en las bases & Bases técnicas del caso Minería, cap. 9.2; requisito técnico RT-05.29; Bases técnicas del caso Minería, cap. 8 (entrevista de Ortúzar); Derivado de esa entrevista: las horas de cierre y de disponibilidad se fijan para cumplir el costo diario \\
\end{tablasintitulo}

\begin{tablasintitulo}{I{0.200}I{0.800}}{\cab{Campo} & \cab{RNF-DES-008}}
Problema & Problema 2 \\
Requerimiento & El sistema debe dejar disponible el cumplimiento del plan minero, comparado contra la versión vigente del plan, al inicio del turno siguiente. \\
Método de verificación & Medición durante 30 turnos consecutivos, contra la línea base de 10 a 15 días de desfase. \\
Origen en las bases & Bases técnicas del caso Minería, cap. 4.1; Bases técnicas del caso Minería, criterio de aceptación 4; requisito técnico RT-05.29; supuesto SUP-14 \\
\end{tablasintitulo}

\begin{tablasintitulo}{I{0.200}I{0.800}}{\cab{Campo} & \cab{RNF-DES-009}}
Problema & Problema 2 \\
Requerimiento & El sistema debe clasificar automáticamente las detenciones a medida que ocurren y dejar el resto disponible para el supervisor, de modo que el 100 \% de las detenciones quede clasificado con causa dentro del turno siguiente al que ocurrieron. Además, el tiempo registrado sin clasificar se mantiene bajo 5 \% en el mes. \\
Método de verificación & Medición del retardo entre fin de detención y disponibilidad de la clasificación, y del porcentaje mensual sin clasificar, contra la línea base de 23 \%. Prueba previa a producción (SUP-05): al reclasificar tres meses de historia, un residuo automático superior al 5 \% indica que las reglas de clasificación están mal calibradas y deben ajustarse antes de la puesta en producción. \\
Origen en las bases & Bases técnicas del caso Minería, cap. 7.2; Bases técnicas del caso Minería, criterio de aceptación 5; supuesto SUP-05 \\
\end{tablasintitulo}

\begin{tablasintitulo}{I{0.200}I{0.800}}{\cab{Campo} & \cab{RNF-DES-010}}
Problema & Problema 2 \\
Requerimiento & El sistema debe mostrar los indicadores con un retraso máximo, respecto del hecho registrado, de 15 minutos para los indicadores operacionales de turno y de 4 horas para los de gestión y de cumplimiento. El cumplimiento del plan minero se publica al inicio del turno siguiente. El estado de equipos y señales se muestra con un retraso máximo de 60 segundos. \\
Método de verificación & Comparación entre marca de tiempo del evento y su disponibilidad en la capa analítica, sobre una muestra de 30 días. \\
Origen en las bases & Bases técnicas del caso Minería, cap. 4.1; requisito técnico RT-05.29 \\
\end{tablasintitulo}

\begin{tablasintitulo}{I{0.200}I{0.800}}{\cab{Campo} & \cab{RNF-DES-011}}
Problema & Problema 2 \\
Requerimiento & El sistema debe cargar la página de la vista única de equipo en no más de 2 segundos y mostrarla con todos sus datos en no más de 5 segundos cuando todas las fuentes están disponibles, o en no más de 8 segundos en modo degradado (cuando una o más fuentes de datos no responden). Los tres tiempos se miden desde que el supervisor solicita la vista y se cumplen en el 95 \% de las solicitudes. \\
Método de verificación & Medición sobre 200 consultas en operación normal y 50 con una integración de telemetría caída. \\
Origen en las bases & Bases técnicas del caso Minería, cap. 9.3; Bases técnicas del caso Minería, criterio de aceptación 6; Bases técnicas transversales, sección 9.1; requisito técnico RT-02.09; Derivado de la sección 9.3 y del criterio de aceptación 6 (vista única): los tiempos de 5 y 8 segundos concretan el requisito \\
\end{tablasintitulo}

\begin{tablasintitulo}{I{0.200}I{0.800}}{\cab{Campo} & \cab{RNF-DES-012}}
Problema & Transversal \\
Requerimiento & El sistema debe completar en un máximo de 4 horas la sincronización posterior a 48 horas de operación desconectada, sin intervención manual y sin degradar la operación en curso. \\
Método de verificación & Corte controlado de 48 horas con cambio de turno dentro de la ventana, midiendo el tiempo hasta la consistencia completa. \\
Origen en las bases & Bases técnicas del caso Minería, cap. 6; Bases técnicas del caso Minería, cap. 9.8; Bases técnicas del caso Minería, criterio de aceptación 13; requisito técnico RT-03.13 de las Bases técnicas del caso Minería, equivalente al RT-03.12 de las Bases técnicas transversales; supuesto SUP-09 \\
\end{tablasintitulo}

\begin{tablasintitulo}{I{0.200}I{0.800}}{\cab{Campo} & \cab{RNF-DES-013}}
Problema & Problema 3 \\
Requerimiento & El sistema debe dejar disponible el cálculo mensual de la huella de carbono por tonelada de cobre fino, auditable por un tercero, dentro de 5 días hábiles desde el cierre del mes. \\
Método de verificación & Medición durante tres cierres consecutivos, contra la línea base de un cálculo anual de dos meses de duración. \\
Origen en las bases & Bases técnicas del caso Minería, cap. 4.9; Bases técnicas del caso Minería, cap. 9.6; Bases técnicas del caso Minería, criterio de aceptación 12; requisito técnico RT-05.23; Derivado del criterio de aceptación 12 y de la sección 9.6 (huella de carbono mensual y auditable): el plazo de 5 días hábiles concreta el requisito \\
\end{tablasintitulo}

\begin{tablasintitulo}{I{0.200}I{0.800}}{\cab{Campo} & \cab{RNF-DES-014}}
Problema & Problema 1 \\
Requerimiento & El sistema debe sostener la diferencia de la conciliación metalúrgica mensual en no más de 3 \% al mes 22 y en no más de 2 \% desde el mes 28, y dejar explicada por punto de la cadena toda diferencia que persista. \\
Método de verificación & Medición durante doce cierres consecutivos, contra la línea base de 3,1 \% a 6,8 \% declarada en el numeral 7.1. \\
Origen en las bases & Bases técnicas del caso Minería, cap. 9.2; Bases técnicas del caso Minería, cap. 7.1; Bases técnicas del caso Minería, criterio de aceptación 3 \\
\end{tablasintitulo}

\begin{tablasintitulo}{I{0.200}I{0.800}}{\cab{Campo} & \cab{RNF-DES-015}}
Problema & Problema 3 \\
Requerimiento & El flujo de acreditación soportado por el sistema debe permitir habilitar a una persona contratista nueva en menos de 48 horas, desde que su empresa carga la documentación completa hasta su habilitación vigente. \\
Método de verificación & Medición sobre 50 acreditaciones consecutivas de empresas contratistas de distinto tamaño, contra la línea base de 9 días. \\
Origen en las bases & Bases técnicas del caso Minería, cap. 4.8; Bases técnicas del caso Minería, cap. 9.5; Bases técnicas del caso Minería, criterio de aceptación 8 \\
\end{tablasintitulo}

\begin{tablasintitulo}{I{0.200}I{0.800}}{\cab{Campo} & \cab{RNF-DES-016}}
Problema & Problema 2 \\
Requerimiento & El sistema debe registrar el aviso de cambio de material sin detener la pala: 0 segundos de pala detenida atribuibles al aviso. Todo paso adicional al ciclo productivo debe justificarse en segundos y demostrarse en la marcha blanca. \\
Método de verificación & Medición del tiempo de ciclo de carguío (el tiempo de una palada) con y sin el aviso habilitado, durante la marcha blanca, en turno de día y de noche, sobre al menos 500 ciclos. \\
Origen en las bases & Bases técnicas del caso Minería, cap. 4.3; Bases técnicas del caso Minería, cap. 9.4; Bases técnicas del caso Minería, restricción no negociable 8; Bases técnicas del caso Minería, entrevista a Painemal \\
\end{tablasintitulo}

\begin{tablasintitulo}{I{0.200}I{0.800}}{\cab{Campo} & \cab{RNF-DES-017}}
Problema & Problema 1 \\
Requerimiento & El sistema debe mantener los tiempos de registro en terreno (2 segundos), de báscula (3 segundos) y de consulta en portería (1,5 segundos) mientras se ejecutan las consultas analíticas de mayor costo. \\
Método de verificación & Prueba de carga concurrente de operación y analítica en Preproducción. \\
Origen en las bases & requisito técnico RT-05.05 \\
\end{tablasintitulo}

\begin{tablasintitulo}{I{0.200}I{0.800}}{\cab{Campo} & \cab{RNF-DES-018}}
Problema & Problema 3 \\
Requerimiento & El sistema debe enviar las notificaciones en segundo plano, con reintento si falla, sin duplicados y con registro de entrega, apertura y error, de modo que el 100 \% de los mensajes tenga su estado registrado. \\
Método de verificación & Prueba de envío con caída simulada del proveedor de un canal. \\
Origen en las bases & requisito técnico RT-16.23 \\
\end{tablasintitulo}

\begin{tablasintitulo}{I{0.200}I{0.800}}{\cab{Campo} & \cab{RNF-DIS-001}}
Problema & Transversal \\
Requerimiento & El sistema debe operar y registrar en faena durante un corte completo del enlace con Antofagasta de al menos 48 horas continuas, manteniendo disponibles las funciones que sostienen la extracción, el transporte, el procesamiento y el registro en terreno, sin que ninguna de ellas dependa de conectividad permanente hacia la nube. \\
Método de verificación & Corte controlado de 48 horas ejecutado en faena, verificando que ninguna de esas funciones se degrada ni se interrumpe y que ningún registro se pierde ni se rechaza. \\
Origen en las bases & Bases técnicas del caso Minería, cap. 6; Bases técnicas del caso Minería, cap. 9.8; Bases técnicas del caso Minería, restricción no negociable 2; requisito técnico RT-03.10; supuesto SUP-09; supuesto SUP-24 \\
\end{tablasintitulo}

\begin{tablasintitulo}{I{0.200}I{0.800}}{\cab{Campo} & \cab{RNF-DIS-002}}
Problema & Transversal \\
Requerimiento & El sistema debe sostener la disponibilidad que corresponde según la clasificación de servicios del Artículo 78.2 de las Bases administrativas, medida sobre la transacción de negocio de extremo a extremo: 99,9 \% mensual para las funciones que sostienen la operación en terreno (servicio crítico, unos 43 minutos de caída al mes como máximo) y 99,5 \% mensual para las de gestión, analítica y portales (servicio alto, unas 3,6 horas al mes como máximo). \\
Método de verificación & Medición mensual con el registro de indisponibilidad del sistema, excluyendo las ventanas de intervención aprobadas. \\
Origen en las bases & Bases técnicas del caso Minería, restricción no negociable 2; requisito técnico RT-10.05; requisito técnico RT-10.01; requisito técnico RT-10.02; Bases administrativas, artículo 20; Bases administrativas, artículo 78 \\
\end{tablasintitulo}

\begin{tablasintitulo}{I{0.200}I{0.800}}{\cab{Campo} & \cab{RNF-DIS-003}}
Problema & Transversal \\
Requerimiento & El sistema no debe perder ningún registro generado en faena durante un corte de enlace: 0 registros perdidos en cortes de hasta 48 horas. \\
Método de verificación & Conciliación entre los registros generados localmente y los registros presentes tras la sincronización, en tres cortes controlados. \\
Origen en las bases & Bases técnicas del caso Minería, cap. 6; Bases técnicas del caso Minería, cap. 9.8; Bases técnicas del caso Minería, criterio de aceptación 13; requisito técnico RT-03.10; requisito técnico RT-03.11; supuesto SUP-09 \\
\end{tablasintitulo}

\begin{tablasintitulo}{I{0.200}I{0.800}}{\cab{Campo} & \cab{RNF-DIS-004}}
Problema & Transversal \\
Requerimiento & El Adjudicatario debe intervenir el sistema solo dentro de la ventana operacional protegida de la faena: intervenciones mayores sólo en las detenciones programadas de marzo y septiembre, intervenciones menores sólo entre 03:00 y 05:00 y previa aprobación, y ninguna intervención durante el cierre contable ni en las 48 horas previas a un embarque, con aviso previo mínimo de 10 días hábiles. \\
Método de verificación & Revisión del registro de intervenciones del periodo contra el calendario operacional de la faena, verificando la ventana y el aviso previo de 10 días hábiles. \\
Origen en las bases & Bases técnicas del caso Minería, cap. 13.2; Bases técnicas del caso Minería, restricción no negociable 1; requisito técnico RT-10.05 \\
\end{tablasintitulo}

\begin{tablasintitulo}{I{0.200}I{0.800}}{\cab{Campo} & \cab{RNF-DIS-005}}
Problema & Problema 2 \\
Requerimiento & El sistema debe mantener operativa la vista única de equipo cuando una o más integraciones de telemetría no estén disponibles: en el 100 \% de esos casos la vista indica qué fuente falta y ninguna otra función queda inaccesible. \\
Método de verificación & Prueba de caída controlada de cada una de las tres integraciones de fabricante, por separado y en conjunto, midiendo el tiempo de presentación de la vista completa contra el umbral en modo degradado de RNF-DES-011 y verificando que la vista indica la fuente que falta y que las demás funciones siguen accesibles. \\
Origen en las bases & Bases técnicas del caso Minería, cap. 4.7; Bases técnicas del caso Minería, criterio de aceptación 6; requisito técnico RT-02.09; supuesto SUP-10 \\
\end{tablasintitulo}

\begin{tablasintitulo}{I{0.200}I{0.800}}{\cab{Campo} & \cab{RNF-DIS-006}}
Problema & Transversal \\
Requerimiento & El sistema debe cumplir para los componentes de terreno un RTO no superior a 1 hora y un RPO de 0 eventos perdidos para el registro operacional (ver también RNF-DIS-003). \\
Método de verificación & Prueba de recuperación de un nodo de borde con corte súbito, cronometrando el tiempo hasta la recuperación del servicio (máximo 1 hora) y verificando que ningún evento del registro operacional se pierde. \\
Origen en las bases & Bases técnicas del caso Minería, cap. 6; Bases técnicas del caso Minería, cap. 9.8; requisito técnico RT-03.11; requisito técnico RT-07.04; Derivado de las 48 horas sin enlace (sección 9.8) y de RT-03.11: el objetivo de recuperación de 1 hora en terreno es más exigente que el de RT-07.04 \\
\end{tablasintitulo}

\begin{tablasintitulo}{I{0.200}I{0.800}}{\cab{Campo} & \cab{RNF-DIS-007}}
Problema & Transversal \\
Requerimiento & El sistema debe degradar de forma controlada ante la indisponibilidad de un componente no crítico o al superar su capacidad: continuar en modo reducido, encolar, limitar la tasa e informar la degradación en pantalla en el 100 \% de los casos, con 0 transacciones perdidas y sin error genérico. \\
Método de verificación & Inyección controlada de fallas y prueba de estrés hasta el punto de quiebre en Preproducción, verificando cada transacción enviada contra las registradas. \\
Origen en las bases & requisito técnico RT-02.09; requisito técnico RT-09.08 \\
\end{tablasintitulo}

\begin{tablasintitulo}{I{0.200}I{0.800}}{\cab{Campo} & \cab{RNF-DIS-008}}
Problema & Transversal \\
Requerimiento & El sistema debe permitir desplegar cambios sin interrupción del servicio: 0 minutos de indisponibilidad de los servicios críticos atribuibles a un despliegue. \\
Método de verificación & Despliegue demostrado en Preproducción antes de cada paso a producción, con medición de disponibilidad durante el despliegue. \\
Origen en las bases & requisito técnico RT-10.06; requisito técnico RT-04.07 \\
\end{tablasintitulo}

\begin{tablasintitulo}{I{0.200}I{0.800}}{\cab{Campo} & \cab{RNF-DIS-009}}
Problema & Problema 1 \\
Requerimiento & El sistema debe mantener los tiempos de registro en terreno (2 segundos), báscula (3 segundos) y portería (1,5 segundos) cuando los portales de empresas contratistas y de compradores operan a 1,5 veces su momento de mayor carga. \\
Método de verificación & Prueba de carga sobre los portales en Preproducción con medición simultánea de los tiempos de terreno. \\
Origen en las bases & requisito técnico RT-16.34; requisito técnico RT-09.06 \\
\end{tablasintitulo}

\begin{tablasintitulo}{I{0.200}I{0.800}}{\cab{Campo} & \cab{RNF-INF-001}}
Problema & Transversal \\
Requerimiento & El Proponente debe especificar, y el sistema debe operar sobre, 1 sala técnica principal en faena de uso exclusivo de la solución, aislada de otras dependencias y con acceso independiente, que habilite o reemplace la sala actual de 40 m² (no cumple el Capítulo 6 de las Bases técnicas transversales; la obra civil es del Cliente y la especificación del Proponente), y 4 gabinetes de borde (equipos locales cerca de cada sitio) en portería, planta, depósito de relaves y acopio de puerto. \\
Método de verificación & Revisión de la arquitectura física contra RT-06.01 y verificación en terreno de cada emplazamiento. \\
Origen en las bases & Bases técnicas del caso Minería, cap. 6; requisito técnico RT-06.01; supuesto SUP-24; supuesto SUP-34 \\
\end{tablasintitulo}

\begin{tablasintitulo}{I{0.200}I{0.800}}{\cab{Campo} & \cab{RNF-INT-001}}
Problema & Problema 1 \\
Requerimiento & El sistema debe leer el historiador de proceso de la planta mediante OPC UA. \\
Método de verificación & Revisión de la integración y verificación del protocolo efectivamente empleado. \\
Origen en las bases & Bases técnicas del caso Minería, cap. 5; Bases técnicas del caso Minería, restricción no negociable 3; requisito técnico RT-05.23; supuesto SUP-16; supuesto SUP-19 \\
\end{tablasintitulo}

\begin{tablasintitulo}{I{0.200}I{0.800}}{\cab{Campo} & \cab{RNF-INT-002}}
Problema & Problema 1 \\
Requerimiento & La integración con el sistema de gestión de laboratorio debe usar una interfaz documentada del propio sistema, sin modificarlo, y entregar cada registro una sola vez en cada sentido (la identificación de cada muestra hacia el laboratorio y cada resultado de vuelta): 0 registros perdidos y 0 duplicados. \\
Método de verificación & Prueba de ida y vuelta sobre una muestra de 200 muestras: se envía la identificación de cada una al sistema de gestión de laboratorio y se recibe su resultado, contando en cada sentido los registros perdidos y los duplicados (se exige 0 y 0). \\
Origen en las bases & Bases técnicas del caso Minería, cap. 5; Bases técnicas del caso Minería, restricción no negociable 6; supuesto SUP-16 \\
\end{tablasintitulo}

\begin{tablasintitulo}{I{0.200}I{0.800}}{\cab{Campo} & \cab{RNF-INT-003}}
Problema & Problema 2 \\
Requerimiento & El sistema debe integrarse en ambos sentidos con el sistema de despacho de flota, cuyo contrato está vigente hasta 2031, sin sustituirlo y sin replicar ninguna de sus funciones. \\
Método de verificación & Revisión funcional conjunta con el proveedor del sistema de despacho. \\
Origen en las bases & Bases técnicas del caso Minería, cap. 5; Bases técnicas del caso Minería, restricción no negociable 5; supuesto SUP-16 \\
\end{tablasintitulo}

\begin{tablasintitulo}{I{0.200}I{0.800}}{\cab{Campo} & \cab{RNF-INT-004}}
Problema & Problema 1 \\
Requerimiento & El sistema debe entregar datos a SAP sin reemplazarlo ni modificarlo, con 0 emisiones contables fuera de SAP, de modo que no existan dos verdades contables. \\
Método de verificación & Revisión de las interfaces con SAP y conciliación entre el dato entregado y el dato contabilizado. \\
Origen en las bases & Bases técnicas del caso Minería, cap. 5; Bases técnicas del caso Minería, restricción no negociable 4; supuesto SUP-16 \\
\end{tablasintitulo}

\begin{tablasintitulo}{I{0.200}I{0.800}}{\cab{Campo} & \cab{RNF-INT-005}}
Problema & Problema 2 \\
Requerimiento & El sistema debe integrar operación y gestión con un modelo de referencia estándar de la industria (ISA-95) y recibir la telemetría de equipos en el estándar abierto de la minería cuando el fabricante lo soporte. \\
Método de verificación & Revisión del modelo de integración y del mapeo de señales por fabricante. \\
Origen en las bases & requisito técnico RT-05.23 \\
\end{tablasintitulo}

\begin{tablasintitulo}{I{0.200}I{0.800}}{\cab{Campo} & \cab{RNF-INT-006}}
Problema & Transversal \\
Requerimiento & El sistema debe documentar el 100 \% de sus interfaces, con servicios síncronos en OpenAPI 3.1 y flujos por eventos en AsyncAPI 2.6 o superior, mantener esa documentación al día con el código, versionar semánticamente sus contratos de interfaz, no romper a quienes ya usan una interfaz cuando esta cambia (compatibilidad hacia atrás) y avisar con al menos 6 meses de anticipación antes de retirar una versión. \\
Método de verificación & Revisión del catálogo de interfaces y de los contratos contra el código desplegado en cada entrega: servicios en OpenAPI 3.1, flujos por eventos en AsyncAPI 2.6 o superior, versionado semántico, pruebas de compatibilidad hacia atrás y registro de avisos de obsolescencia con 6 meses de preaviso. \\
Origen en las bases & requisito técnico RT-05.16; requisito técnico RT-05.17 \\
\end{tablasintitulo}

\begin{tablasintitulo}{I{0.200}I{0.800}}{\cab{Campo} & \cab{RNF-POR-001}}
Problema & Transversal \\
Requerimiento & El sistema debe poder replicarse a una nueva unidad operacional (otra faena, como la que la compañía evalúa adquirir hacia 2030) mediante parametrización, con catálogos, unidades organizacionales y reglas de negocio independientes, en no más de 20 días hábiles y sin modificar el código fuente. \\
Método de verificación & Ejercicio cronometrado de configuración de una unidad operacional adicional en ambiente de prueba, verificando que se completa en el plazo declarado y sin modificar código fuente. \\
Origen en las bases & Bases técnicas del caso Minería, cap. 13.2; requisito técnico RT-02.12; Derivado de la sección 13.2 y de RT-02.12 (replicar la solución a otras unidades): el plazo de 20 días hábiles concreta el requisito \\
\end{tablasintitulo}

\begin{tablasintitulo}{I{0.200}I{0.800}}{\cab{Campo} & \cab{RNF-POR-002}}
Problema & Problema 3 \\
Requerimiento & El sistema debe permitir exportar el 100 \% de la información del Cliente, incluidos los registros y sus evidencias, en formatos abiertos y documentados, en cualquier momento del Contrato y sin intervención del Adjudicatario, de modo que la información pueda ser auditada por un tercero sin herramientas propietarias. \\
Método de verificación & Prueba de exportación y lectura por un tercero sin acceso a la solución. \\
Origen en las bases & Bases técnicas del caso Minería, cap. 12; Bases técnicas del caso Minería, criterio de aceptación 12; requisito técnico RT-05.06 \\
\end{tablasintitulo}

\begin{tablasintitulo}{I{0.200}I{0.800}}{\cab{Campo} & \cab{RNF-POR-003}}
Problema & Problema 2 \\
Requerimiento & El sistema debe distribuir su aplicación móvil por las tiendas oficiales o por gestión de flota corporativa, con firma de la aplicación y verificación de integridad, con el 100 \% de las instalaciones con firma verificada. \\
Método de verificación & Revisión del proceso de distribución e intento de instalar una versión no firmada. \\
Origen en las bases & requisito técnico RT-17.04 \\
\end{tablasintitulo}

\begin{tablasintitulo}{I{0.200}I{0.800}}{\cab{Campo} & \cab{RNF-POR-004}}
Problema & Problema 2 \\
Requerimiento & El sistema debe soportar, en su aplicación móvil, la versión vigente del sistema operativo móvil y las 2 anteriores. \\
Método de verificación & Prueba de la aplicación en cada versión soportada. \\
Origen en las bases & requisito técnico RT-17.03 \\
\end{tablasintitulo}

\begin{tablasintitulo}{I{0.200}I{0.800}}{\cab{Campo} & \cab{RNF-SEG-001}}
Problema & Problema 3 \\
Requerimiento & El sistema debe cifrar campo por campo los datos de salud ocupacional, exámenes preocupacionales, resultados de control de fatiga y resultados de control de alcohol y drogas, de modo que quien acceda a la base de datos no pueda leer su contenido. El cifrado en tránsito, en reposo y en los dispositivos se exige aparte. \\
Método de verificación & Revisión de arquitectura y prueba de lectura directa de la base de datos, verificando que los campos de las 4 categorías sensibles no revelan su contenido. \\
Origen en las bases & Bases técnicas del caso Minería, cap. 4.8; requisito técnico RT-11.10; requisito técnico RT-05.08; Ley N° 21.719 \\
\end{tablasintitulo}

\begin{tablasintitulo}{I{0.200}I{0.800}}{\cab{Campo} & \cab{RNF-SEG-002}}
Problema & Problema 3 \\
Requerimiento & El sistema debe registrar el 100 \% de las consultas a información sensible, además de toda modificación, conservando quién, cuándo y sobre qué persona o registro, con una retención de 10 años. \\
Método de verificación & Auditoría sobre una muestra de accesos, contrastando el registro con los accesos efectivamente realizados. \\
Origen en las bases & Bases técnicas del caso Minería, cap. 4.8; requisito técnico RT-16.09; requisito técnico RT-05.10; requisito técnico RT-05.07; requisito técnico RT-16.10 \\
\end{tablasintitulo}

\begin{tablasintitulo}{I{0.200}I{0.800}}{\cab{Campo} & \cab{RNF-SEG-003}}
Problema & Problema 3 \\
Requerimiento & El sistema debe usar las cuentas del directorio corporativo del grupo para el personal propio (identidad federada), con doble factor de autenticación obligatorio, sin credenciales compartidas y con registro completo de accesos. Debe exigir doble factor en el 100 \% de los accesos administrativos, privilegiados y desde fuera de la red corporativa, y permitir que cada empresa contratista administre la identidad de su personal desde el portal de empresas contratistas. \\
Método de verificación & Revisión del registro de accesos y de la configuración de autenticación; prueba de intento de uso de credencial compartida; prueba de acceso administrativo, privilegiado y externo sin segundo factor; y prueba de alta y baja de una persona contratista desde el portal de empresas. \\
Origen en las bases & Bases técnicas del caso Minería, cap. 4.8; Bases técnicas del caso Minería, cap. 9.5; Bases técnicas del caso Minería, restricción no negociable 9; Bases técnicas del caso Minería, criterio de aceptación 8; Bases técnicas del caso Minería, decisión pendiente 8; requisito técnico RT-12.01; requisito técnico RT-12.03; requisito técnico RT-12.12; requisito técnico RT-16.30; supuesto SUP-08 \\
\end{tablasintitulo}

\begin{tablasintitulo}{I{0.200}I{0.800}}{\cab{Campo} & \cab{RNF-SEG-004}}
Problema & Problema 1 \\
Requerimiento & El sistema debe conectarse al servidor espejo del historiador de proceso de la planta como cliente OPC UA sin capacidad de escritura, dentro de una red segregada entre tecnología de información y tecnología de operación conforme a IEC 62443 y a la aprobación de la compañía, de modo que ningún intento de escritura pueda llegar a la red de control. \\
Método de verificación & Revisión de arquitectura de red y prueba de intento de escritura sobre el control de proceso. \\
Origen en las bases & Bases técnicas del caso Minería, cap. 5; Bases técnicas del caso Minería, cap. 12; Bases técnicas del caso Minería, restricción no negociable 3 \\
\end{tablasintitulo}

\begin{tablasintitulo}{I{0.200}I{0.800}}{\cab{Campo} & \cab{RNF-SEG-005}}
Problema & Problema 3 \\
Requerimiento & El sistema debe conservar los registros, con integridad verificable y disponibilidad para fiscalización, por los siguientes plazos: seguridad y salud ocupacional 10 años; ambientales y de monitoreo 10 años; trazabilidad de mineral y embarques 10 años; operacionales 5 años; tributarios 7 años; trazas técnicas 90 días. \\
Método de verificación & Prueba de recuperación de registros de cada categoría en el extremo de su plazo de retención. \\
Origen en las bases & Bases técnicas del caso Minería, cap. 4.8; requisito técnico RT-05.10; requisito técnico RT-05.07; requisito técnico RT-16.10; Derivado de RT-05.10 (retenciones definidas): la retención de 90 días de las trazas técnicas complementa esas retenciones \\
\end{tablasintitulo}

\begin{tablasintitulo}{I{0.200}I{0.800}}{\cab{Campo} & \cab{RNF-SEG-006}}
Problema & Problema 3 \\
Requerimiento & El sistema debe revocar el acceso de una persona en un máximo de 24 horas desde el término de su relación con la faena. \\
Método de verificación & Prueba de revocación cronometrada sobre 20 casos, incluida la propagación a los nodos de borde. \\
Origen en las bases & Bases técnicas del caso Minería, cap. 4.8; Bases técnicas del caso Minería, cap. 9.5; requisito técnico RT-12.10 \\
\end{tablasintitulo}

\begin{tablasintitulo}{I{0.200}I{0.800}}{\cab{Campo} & \cab{RNF-SEG-007}}
Problema & Transversal \\
Requerimiento & El sistema debe autenticar cada integración entre sistemas con credenciales seguras y no con claves fijas: 0 integraciones con una clave fija incluida en la dirección web (URL). \\
Método de verificación & Revisión de la configuración de cada integración (OAuth 2.1 con credenciales de cliente o autenticación mutua TLS, según RT-05.18) y prueba de acceso sin credencial válida. \\
Origen en las bases & requisito técnico RT-05.18 \\
\end{tablasintitulo}

\begin{tablasintitulo}{I{0.200}I{0.800}}{\cab{Campo} & \cab{RNF-SEG-008}}
Problema & Transversal \\
Requerimiento & El sistema debe cifrar el 100 \% de las comunicaciones en tránsito con la versión vigente del protocolo de cifrado (TLS 1.3), con 0 servicios que acepten versiones antiguas, y renovar los certificados automáticamente con alerta anticipada de vencimiento. \\
Método de verificación & Escaneo de configuración de todos los servicios expuestos e internos: TLS 1.3, ningún servicio con TLS 1.0 o 1.1, HSTS con precarga y rotación de certificados (RT-11.08). \\
Origen en las bases & requisito técnico RT-11.08 \\
\end{tablasintitulo}

\begin{tablasintitulo}{I{0.200}I{0.800}}{\cab{Campo} & \cab{RNF-SEG-009}}
Problema & Transversal \\
Requerimiento & El sistema debe cifrar el 100 \% de los datos en reposo, en nube, en faena y en respaldos, con claves en un servicio de gestión de claves o módulo de seguridad de hardware, rotación declarada y separación de funciones en la custodia de claves. \\
Método de verificación & Revisión de configuración de almacenamiento y de la política de claves. \\
Origen en las bases & requisito técnico RT-11.09 \\
\end{tablasintitulo}

\begin{tablasintitulo}{I{0.200}I{0.800}}{\cab{Campo} & \cab{RNF-SEG-010}}
Problema & Problema 1 \\
Requerimiento & El sistema debe proteger el 100 \% de los puntos de entrada públicos de los portales de empresas contratistas y de compradores contra bots y abuso automatizado, con reto progresivo conforme a WCAG 2.2 AA que no degrade la accesibilidad ni bloquee a personas usuarias legítimas. \\
Método de verificación & Prueba de abuso automatizado y revisión de accesibilidad del reto. \\
Origen en las bases & requisito técnico RT-16.30 de las Bases técnicas del caso Minería, equivalente al RT-16.31 de las Bases técnicas transversales; requisito técnico RT-11.12 \\
\end{tablasintitulo}

\begin{tablasintitulo}{I{0.200}I{0.800}}{\cab{Campo} & \cab{RNF-SEG-011}}
Problema & Problema 3 \\
Requerimiento & El sistema debe registrar los eventos de seguridad, incluidos los accesos a la sala técnica, de forma centralizada e inalterable, con una retención de 12 meses en línea y 24 meses adicionales en archivo recuperable. \\
Método de verificación & Revisión de la retención configurada y recuperación de un evento archivado. \\
Origen en las bases & requisito técnico RT-11.14 \\
\end{tablasintitulo}

\begin{tablasintitulo}{I{0.200}I{0.800}}{\cab{Campo} & \cab{RNF-SEG-012}}
Problema & Problema 3 \\
Requerimiento & El sistema debe ofrecer inicio de sesión único para todos sus módulos, con 1 autenticación por sesión, y propagar el cierre de sesión al 100 \% de ellos. \\
Método de verificación & Prueba de navegación entre módulos y de cierre de sesión. \\
Origen en las bases & requisito técnico RT-12.02 \\
\end{tablasintitulo}

\begin{tablasintitulo}{I{0.200}I{0.800}}{\cab{Campo} & \cab{RNF-SEG-013}}
Problema & Transversal \\
Requerimiento & El sistema debe dar permisos de administrador solo por tiempo limitado y con aprobación previa, grabar el 100 \% de las sesiones de administración de las operaciones de mayor riesgo y no tener cuentas de administrador permanente (0), salvo una cuenta de acceso de emergencia cuyo procedimiento, custodia, control y auditoría debe declarar el Proponente. \\
Método de verificación & Revisión de cuentas privilegiadas y de grabaciones sobre una muestra de sesiones; revisión de la declaración del procedimiento de acceso de emergencia y verificación del control y la auditoría de esa cuenta. \\
Origen en las bases & requisito técnico RT-12.06; requisito técnico RT-12.13 \\
\end{tablasintitulo}

\begin{tablasintitulo}{I{0.200}I{0.800}}{\cab{Campo} & \cab{RNF-SEG-014}}
Problema & Problema 3 \\
Requerimiento & El sistema debe cerrar la sesión de una persona por duración máxima o por inactividad, con tiempos que el Cliente pueda modificar, permitir revocarla de inmediato, controlar las sesiones simultáneas de una misma persona y no mostrar identificadores de sesión en la dirección web (0). \\
Método de verificación & Pruebas de caducidad, revocación y manipulación de credenciales de sesión (credenciales firmadas, de vida breve y con refresco rotatorio, según RT-12.07 y RT-12.08). \\
Origen en las bases & requisito técnico RT-12.07; requisito técnico RT-12.08 \\
\end{tablasintitulo}

\begin{tablasintitulo}{I{0.200}I{0.800}}{\cab{Campo} & \cab{RNF-SEG-015}}
Problema & Problema 3 \\
Requerimiento & El sistema debe excluir de sus registros técnicos los datos personales sensibles y las credenciales (0 casos), con acceso a los registros controlado y auditado. \\
Método de verificación & Revisión automatizada de muestras de registros técnicos. \\
Origen en las bases & requisito técnico RT-14.07; Ley N° 21.719 \\
\end{tablasintitulo}

\begin{tablasintitulo}{I{0.200}I{0.800}}{\cab{Campo} & \cab{RNF-SEG-016}}
Problema & Transversal \\
Requerimiento & El sistema debe cifrar el 100 \% de la información almacenada en los dispositivos de terreno y permitir su bloqueo y borrado remoto ante pérdida o desvinculación de la persona usuaria, con borrado efectivo en la siguiente conexión del dispositivo. \\
Método de verificación & Prueba de extracción del almacenamiento de un dispositivo y de borrado remoto. \\
Origen en las bases & requisito técnico RT-17.05 \\
\end{tablasintitulo}

\begin{tablasintitulo}{I{0.200}I{0.800}}{\cab{Campo} & \cab{RNF-USA-001}}
Problema & Problema 2 \\
Requerimiento & El sistema debe permitir operar el 100 \% de las funciones de registro en terreno con guantes de invierno puestos, sin que la persona usuaria deba retirarlos en ningún paso. \\
Método de verificación & Prueba con personal de operación en turno de noche de invierno, con los guantes efectivamente usados en faena. \\
Origen en las bases & Bases técnicas del caso Minería, cap. 9.4; Bases técnicas del caso Minería, criterio de aceptación 14; requisito técnico RT-13.08; requisito técnico RT-12.11 \\
\end{tablasintitulo}

\begin{tablasintitulo}{I{0.200}I{0.800}}{\cab{Campo} & \cab{RNF-USA-002}}
Problema & Problema 2 \\
Requerimiento & El sistema debe mantener legible la interfaz de terreno bajo radiación solar directa, de modo que el 100 \% de los registros se complete sin sombrear la pantalla, en una prueba con 10 personas operadoras. \\
Método de verificación & Prueba en cabina al mediodía en verano, con 10 personas operadoras, midiendo la proporción de registros completados sin sombrear la pantalla con la mano o el cuerpo. \\
Origen en las bases & Bases técnicas del caso Minería, cap. 4.3; Bases técnicas del caso Minería, cap. 9.4; Bases técnicas del caso Minería, cap. 13.3; requisito técnico RT-13.08; Derivado de RT-13.08 (legibilidad): la prueba con 10 operadores concreta cómo verificarla \\
\end{tablasintitulo}

\begin{tablasintitulo}{I{0.200}I{0.800}}{\cab{Campo} & \cab{RNF-USA-003}}
Problema & Problema 2 \\
Requerimiento & El Proponente debe especificar (Formulario T-11) equipamiento de terreno que funcione entre -6 °C y 34 °C, con resistencia a polvo en suspensión permanente y a vibración de equipo minero, sin degradar el tiempo de respuesta comprometido, aunque su compra sea de cargo del Cliente. \\
Método de verificación & Prueba en cabina de pala y de camión, con el equipamiento especificado, durante un turno completo de invierno y uno de verano. \\
Origen en las bases & Bases técnicas del caso Minería, cap. 6; requisito técnico RT-13.08; requisito técnico RT-08.10; requisito técnico RT-08.11; supuesto SUP-33; supuesto SUP-18 \\
\end{tablasintitulo}

\begin{tablasintitulo}{I{0.200}I{0.800}}{\cab{Campo} & \cab{RNF-USA-004}}
Problema & Problema 2 \\
Requerimiento & El sistema debe permitir al supervisor de turno clasificar el residuo de detenciones no resueltas automáticamente en no más de 5 minutos por turno, desde lista cerrada y sin texto libre. \\
Método de verificación & Medición del tiempo efectivo de clasificación durante 30 turnos, con el volumen real de detenciones. \\
Origen en las bases & Bases técnicas del caso Minería, cap. 7.2; Bases técnicas del caso Minería, restricción no negociable 8; Bases técnicas del caso Minería, criterio de aceptación 5; supuesto SUP-05 \\
\end{tablasintitulo}

\begin{tablasintitulo}{I{0.200}I{0.800}}{\cab{Campo} & \cab{RNF-USA-005}}
Problema & Problema 2 \\
Requerimiento & El sistema debe lograr que la persona operadora de pala registre por su cuenta, sin que el supervisor se lo pida, al menos el 90 \% de los eventos que debían registrarse (los que una fuente independiente, como el despacho, acredita que ocurrieron; en la botonera, los cambios de material). \\
Método de verificación & Medición en las cuatro últimas semanas de la marcha blanca de la Etapa 1, contra la línea base de dos iniciativas previas abandonadas. \\
Origen en las bases & Bases técnicas del caso Minería, cap. 9.4; Bases técnicas del caso Minería, criterio de aceptación 14; Derivado de la sección 9.4 y del criterio de aceptación 14 (que el registro nazca donde ocurre el hecho): el 90 \% concreta la meta de adopción \\
\end{tablasintitulo}

\begin{tablasintitulo}{I{0.200}I{0.800}}{\cab{Campo} & \cab{RNF-USA-006}}
Problema & Problema 2 \\
Requerimiento & El sistema debe presentar en español el 100 \% de sus interfaces, mensajes, alertas y documentación dirigida a personas usuarias. \\
Método de verificación & Revisión de la totalidad de las pantallas y mensajes en cada entrega. \\
Origen en las bases & requisito técnico RT-22.03; Bases administrativas, artículo 8 \\
\end{tablasintitulo}

\begin{tablasintitulo}{I{0.200}I{0.800}}{\cab{Campo} & \cab{RNF-USA-007}}
Problema & Problema 3 \\
Requerimiento & El sistema debe permitir completar un registro de seguridad en terreno, con firma, sin conexión y en no más de 3 minutos por registro de inspección. \\
Método de verificación & Medición con inspectores de seguridad en terreno, en zona sin cobertura, sobre 50 registros. \\
Origen en las bases & Bases técnicas del caso Minería, cap. 4.8; Bases técnicas del caso Minería, criterio de aceptación 9; requisito técnico RT-17.01; Derivado de la sección 4.8 y del criterio de aceptación 9: el máximo de 3 minutos concreta el registro ágil \\
\end{tablasintitulo}

\begin{tablasintitulo}{I{0.200}I{0.800}}{\cab{Campo} & \cab{RNF-USA-008}}
Problema & Problema 2 \\
Requerimiento & El sistema no debe requerir de la persona operadora de pala ninguna acción de registro distinta del aviso de cambio de material y de la identificación por credencial en cada relevo, ni en cabina ni en ningún otro dispositivo: 1 sola acción de registro en cabina y 1 acercamiento de credencial por relevo, de no más de 2 segundos en el percentil 95. \\
Método de verificación & Revisión funcional de la interfaz de cabina y medición del tiempo de ciclo durante la marcha blanca. \\
Origen en las bases & Bases técnicas del caso Minería, cap. 4.3; Bases técnicas del caso Minería, cap. 9.4; Bases técnicas del caso Minería, restricción no negociable 8; Bases técnicas del caso Minería, entrevista a Painemal \\
\end{tablasintitulo}

\begin{tablasintitulo}{I{0.200}I{0.800}}{\cab{Campo} & \cab{RNF-USA-009}}
Problema & Problema 2 \\
Requerimiento & El sistema no debe requerir de la persona operadora de equipo ninguna acción de clasificación de detenciones, ni en cabina ni en ningún otro dispositivo: 0 acciones de clasificación solicitadas a la persona operadora. \\
Método de verificación & Revisión funcional de las interfaces de cabina y de terreno durante la marcha blanca. \\
Origen en las bases & Bases técnicas del caso Minería, cap. 7.2; Bases técnicas del caso Minería, restricción no negociable 8; supuesto SUP-05 \\
\end{tablasintitulo}

\begin{tablasintitulo}{I{0.200}I{0.800}}{\cab{Campo} & \cab{RNF-USA-010}}
Problema & Problema 2 \\
Requerimiento & El sistema debe cumplir WCAG 2.2 nivel AA en todas sus interfaces destinadas a personas usuarias, con 0 incumplimientos de criterios de nivel A y AA. \\
Método de verificación & Verificación con herramientas automatizadas y pruebas manuales, con informe de conformidad. \\
Origen en las bases & requisito técnico RT-13.01 \\
\end{tablasintitulo}

\begin{tablasintitulo}{I{0.200}I{0.800}}{\cab{Campo} & \cab{RNF-USA-011}}
Problema & Problema 2 \\
Requerimiento & El sistema debe funcionar con diseño responsivo en escritorio, tableta y teléfono, adaptando el contenido y no sólo reduciéndolo, en las resoluciones 1920×1080, 1366×768, 768×1024 (vertical y horizontal) y 375×812. \\
Método de verificación & Revisión de cada pantalla en las cuatro resoluciones. \\
Origen en las bases & requisito técnico RT-13.02; requisito técnico RT-25.04 \\
\end{tablasintitulo}

\begin{tablasintitulo}{I{0.200}I{0.800}}{\cab{Campo} & \cab{RNF-USA-012}}
Problema & Problema 2 \\
Requerimiento & El sistema debe permitir alcanzar cada funcionalidad principal desde la pantalla de inicio del perfil en un máximo de 3 interacciones. \\
Método de verificación & Recorrido de cada funcionalidad principal por perfil, contando interacciones. \\
Origen en las bases & requisito técnico RT-13.05 \\
\end{tablasintitulo}

\begin{tablasintitulo}{I{0.200}I{0.800}}{\cab{Campo} & \cab{RNF-USA-013}}
Problema & Problema 2 \\
Requerimiento & El sistema debe entregar retroalimentación visual ante cada acción y presentar los errores indicando qué ocurrió y qué hacer, con causa y acción sugerida en el 100 \% de los mensajes de error y 0 mensajes técnicos dirigidos a la persona usuaria final. \\
Método de verificación & Revisión del catálogo de mensajes y prueba de usabilidad con personas usuarias del Cliente. \\
Origen en las bases & requisito técnico RT-13.06 \\
\end{tablasintitulo}

\begin{tablasintitulo}{I{0.200}I{0.800}}{\cab{Campo} & \cab{RNF-USA-014}}
Problema & Problema 2 \\
Requerimiento & El sistema debe soportar a personas usuarias con baja alfabetización digital mediante alto contraste, objetivos táctiles de al menos 44 × 44 píxeles, iconografía acompañada de texto en el 100 \% de los íconos y flujos guiados paso a paso. \\
Método de verificación & Revisión de interfaces de terreno y prueba con personas operadoras. \\
Origen en las bases & requisito técnico RT-13.07 \\
\end{tablasintitulo}

\begin{tablasintitulo}{I{0.200}I{0.800}}{\cab{Campo} & \cab{RNF-USA-015}}
Problema & Problema 2 \\
Requerimiento & El sistema debe ser navegable íntegramente por teclado, con el 100 \% de las funciones de escritorio operables sólo con teclado, orden de foco lógico y atajos para las operaciones frecuentes. \\
Método de verificación & Prueba manual de navegación por teclado en cada pantalla de escritorio. \\
Origen en las bases & requisito técnico RT-13.11 \\
\end{tablasintitulo}

\begin{tablasintitulo}{I{0.200}I{0.800}}{\cab{Campo} & \cab{RNF-OPE-001}}
Problema & Transversal \\
Requerimiento & El sistema debe mostrar en una sola consola el estado de sus componentes, de sus integraciones con terceros y de sus colas de sincronización de borde, de modo que el equipo de TI del Cliente (11 personas, 4 en faena) diagnostique una falla de integración en no más de 30 minutos, sin entrar a los sistemas subyacentes ni usar otras herramientas. \\
Método de verificación & Prueba de diagnóstico de una falla de integración simulada, cronometrada, ejecutada por personal de TI del Cliente desde la consola única, con un tiempo de diagnóstico no superior a 30 minutos. \\
Origen en las bases & Bases técnicas del caso Minería, cap. 8 (entrevista de Alcaíno); Bases técnicas del caso Minería, cap. 17.1; Bases técnicas del caso Minería, restricción no negociable 10; requisito técnico RT-14.01; requisito técnico RT-03.16; Derivado de la entrevista de Alcaíno y de RT-14.01: el diagnóstico en 30 minutos concreta ese requisito \\
\end{tablasintitulo}

\begin{tablasintitulo}{I{0.200}I{0.800}}{\cab{Campo} & \cab{RNF-OPE-002}}
Problema & Transversal \\
Requerimiento & El sistema debe emitir alertas anticipadas de degradación al equipo de operación: colas de sincronización cerca de su límite, integraciones de fabricante caídas e indicadores atrasados sobre sus umbrales (15 minutos los de turno y 4 horas los de gestión). En las colas, la alerta debe emitirse al llegar al 80 \% del límite de 150.000 eventos (120.000), en no más de 5 minutos desde que se cruza ese umbral, lo que deja un margen de al menos 4 horas al ritmo de diseño de unos 3.000 eventos por hora. \\
Método de verificación & Prueba de inyección controlada de degradación (crecimiento de cola, caída de integración) verificando el tiempo de emisión de la alerta y el margen respecto del umbral comprometido. Se verifica que la alerta se emite al cruzar 120.000 eventos, en no más de 5 minutos, y que el margen restante equivale a al menos 4 horas de acumulación al ritmo medio. \\
Origen en las bases & requisito técnico RT-14.04; Derivado de RNF-CAP-006 (cola local de 150.000 eventos): los umbrales de alerta se calculan a partir de esa capacidad \\
\end{tablasintitulo}

\begin{tablasintitulo}{I{0.200}I{0.800}}{\cab{Campo} & \cab{RNF-OPE-003}}
Problema & Transversal \\
Requerimiento & El sistema debe contar con una política de respaldo de 3 copias en 2 medios distintos, 1 fuera del sitio y 1 inmutable o fuera de línea, con 0 errores en las verificaciones de restauración (regla 3-2-1-1-0). Los respaldos deben ir cifrados en reposo y en tránsito con una clave gestionada independiente de la infraestructura respaldada, y las copias inmutables deben protegerse contra borrado y modificación durante su retención, incluso con credenciales administrativas comprometidas. El Proponente debe declarar, por tipo de dato (producción y trazabilidad, mantenimiento, seguridad y salud ocupacional, ambiental y laboratorio), la frecuencia de respaldo, el período de retención y el tiempo estimado de restauración completa. \\
Método de verificación & Revisión de la política y de la configuración de respaldo contra el esquema 3-2-1-1-0 y de los resultados de verificación de restauración; revisión del cifrado y de la gestión de claves; intento de borrado o modificación de una copia inmutable con credenciales administrativas; revisión de la declaración por dominio de datos. \\
Origen en las bases & requisito técnico RT-07.09; requisito técnico RT-07.10; requisito técnico RT-07.11; requisito técnico RT-07.13 \\
\end{tablasintitulo}

\begin{tablasintitulo}{I{0.200}I{0.800}}{\cab{Campo} & \cab{RNF-OPE-004}}
Problema & Transversal \\
Requerimiento & El sistema debe permitir probar la restauración de los respaldos al menos una vez al mes sobre una muestra representativa, midiendo y documentando el tiempo efectivo de restauración. \\
Método de verificación & Ejecución y acta de la prueba mensual de restauración, con el tiempo medido contra el declarado en el tiempo de restauración que declara el Proponente en RNF-OPE-003. \\
Origen en las bases & requisito técnico RT-07.12 \\
\end{tablasintitulo}

\begin{tablasintitulo}{I{0.200}I{0.800}}{\cab{Campo} & \cab{RNF-OPE-005}}
Problema & Transversal \\
Requerimiento & El sistema debe recuperarse ante una falla mayor del sitio central con un tiempo máximo de recuperación (RTO) de 4 horas y una pérdida máxima de datos (RPO) de 15 minutos para los servicios críticos centrales, con conmutación al sitio secundario probada en real al menos dos veces al año y antes del paso a producción. Los componentes de terreno tienen sus propios objetivos (RTO 1 hora, RPO 0). \\
Método de verificación & Prueba de conmutación real con informe del RTO y RPO efectivamente alcanzados y plan de corrección de brechas. \\
Origen en las bases & Bases técnicas del caso Minería, cap. 6; requisito técnico RT-07.04; requisito técnico RT-07.05; requisito técnico RT-07.07 \\
\end{tablasintitulo}

\begin{tablasintitulo}{I{0.200}I{0.800}}{\cab{Campo} & \cab{RNF-OPE-006}}
Problema & Transversal \\
Requerimiento & El sistema debe permitir seguir cualquier operación de punta a punta por todos los sistemas, en la nube y en faena, mediante un identificador único de transacción (estándar OpenTelemetry) presente en el 100 \% de las transacciones e integraciones, que registre la entrada y la salida de cada integración. Debe integrar el monitoreo del componente en faena en la misma plataforma que el de la nube y medir los indicadores de nivel de servicio sobre la experiencia real de las personas usuarias y no solo sobre pruebas simuladas. \\
Método de verificación & Seguimiento de una operación de negocio de punta a punta a través de todos los sistemas involucrados, en la nube y en faena, y verificación, sobre una muestra de transacciones e integraciones, de que todas portan el identificador único y registran su entrada y su salida, y de que el monitoreo del componente on-premise aparece en la misma plataforma de observabilidad que el de la nube. Revisión de la definición y del origen de los datos de cada indicador de nivel de servicio. \\
Origen en las bases & requisito técnico RT-14.01; requisito técnico RT-14.03; requisito técnico RT-05.19; requisito técnico RT-03.16 \\
\end{tablasintitulo}

\begin{tablasintitulo}{I{0.200}I{0.800}}{\cab{Campo} & \cab{RNF-TRV-02.06}}
Problema & Transversal \\
Requerimiento & El sistema debe garantizar que el 100 \% de las operaciones de escritura expuestas a reintentos sean idempotentes mediante una clave de idempotencia declarada por el cliente, con 0 efectos duplicados ante el reenvio de una misma solicitud dentro de la ventana de deduplicacion documentada de 72 horas, que cubre el corte de enlace de 48 horas más la sincronización posterior. \\
Método de verificación & Prueba: reenviar cada operacion de escritura expuesta a reintentos con la misma clave y verificar 0 efectos duplicados; inspeccion de la documentacion de la ventana de deduplicacion, incluido un reenvio a las 72 horas. \\
Origen en las bases & requisito técnico RT-02.06; Derivado de RT-02.06 (ventana de deduplicación documentada) y de las 48 horas sin enlace: las 72 horas cubren el corte más la sincronización \\
\end{tablasintitulo}

\begin{tablasintitulo}{I{0.200}I{0.800}}{\cab{Campo} & \cab{RNF-TRV-02.07}}
Problema & Transversal \\
Requerimiento & El sistema debe entregar cada evento de negocio al menos una vez, con 0 duplicados con efecto en el consumidor (deduplicacion en el consumidor) y con el orden preservado dentro de la particion o del agregado en los procesos que lo exijan. \\
Método de verificación & Prueba: inyectar fallas y reenvios en el bus de eventos y verificar 0 eventos perdidos, 0 efectos duplicados en el consumidor y orden por particion/agregado. \\
Origen en las bases & requisito técnico RT-02.07 \\
\end{tablasintitulo}

\begin{tablasintitulo}{I{0.200}I{0.800}}{\cab{Campo} & \cab{RNF-TRV-02.08}}
Problema & Transversal \\
Requerimiento & El sistema debe aplicar en el 100 \% de las llamadas remotas un tiempo de espera explicito, con reintento con retroceso exponencial y variacion aleatoria, cortacircuitos, mamparos de aislamiento y limites de tasa, con 0 llamadas remotas sin tiempo de espera (tiempo de espera de 1 segundo en consultas y de 2 segundos en escrituras; hasta 3 reintentos con retroceso exponencial de base 200 milisegundos y variación aleatoria de ±20 \%; límite de tasa de 50 llamadas por segundo por dependencia, con ráfaga de 100). \\
Método de verificación & Inspeccion de configuracion (1 s en consultas, 2 s en escrituras, 3 reintentos, 50 llamadas por segundo) de todas las llamadas remotas y prueba de falla inyectada de cada dependencia verificando el disparo del cortacircuitos, el aislamiento y el retroceso. \\
Origen en las bases & requisito técnico RT-02.08; Derivado de RT-02.08 (patrones de resiliencia): los tiempos, reintentos y límites concretan esos patrones \\
\end{tablasintitulo}

\begin{tablasintitulo}{I{0.200}I{0.800}}{\cab{Campo} & \cab{RNF-TRV-04.11}}
Problema & Transversal \\
Requerimiento & El Adjudicatario debe mantener una cobertura de pruebas automatizadas de al menos 70 \% sobre el código de lógica de negocio de la solución y bloquear la integración de todo cambio que deje la cobertura bajo ese umbral. \\
Método de verificación & Reporte de cobertura emitido por la herramienta de pruebas sobre el código de lógica de negocio en cada corrida de integración, con resultado igual o superior a 70 \%, y prueba de que un cambio bajo el umbral es rechazado. \\
Origen en las bases & requisito técnico RT-04.11 \\
\end{tablasintitulo}

\begin{tablasintitulo}{I{0.200}I{0.800}}{\cab{Campo} & \cab{RNF-TRV-07.03}}
Problema & Transversal \\
Requerimiento & El sistema debe replicar de forma continua los datos del sitio principal al sitio secundario, medir el retraso de replicación y alertar al equipo de operación cuando ese retraso supere 5 minutos (un tercio del RPO de 15 minutos de RNF-OPE-005), con 0 retrasos por sobre el umbral sin alerta emitida. \\
Método de verificación & Prueba en Preproducción y en el sitio secundario: inducir retraso de replicación por sobre el umbral de 5 minutos y verificar que la medición se registra y la alerta llega al equipo de operación; revisión de la métrica de retraso en la consola operativa durante la marcha blanca. \\
Origen en las bases & requisito técnico RT-07.03; Derivado de RT-07.03 (medir y alertar el retraso de replicación): el umbral de 5 minutos se fija a partir del RPO de RT-07.04 \\
\end{tablasintitulo}

\begin{tablasintitulo}{I{0.200}I{0.800}}{\cab{Campo} & \cab{RNF-TRV-11.04}}
Problema & Transversal \\
Requerimiento & El Adjudicatario debe mantener un programa de gestión de vulnerabilidades con escaneo continuo del sistema y remediar las vulnerabilidades críticas en un máximo de 7 días corridos, las altas en 15 y las medias en 30, contados desde su publicación o detección, con 0 vulnerabilidades fuera de plazo. \\
Método de verificación & Informe mensual del escáner con fecha de detección y de cierre por vulnerabilidad, contrastado por muestreo con la evidencia de remediación; auditoría del Cliente sobre el registro de vulnerabilidades. \\
Origen en las bases & requisito técnico RT-11.04 \\
\end{tablasintitulo}

\begin{tablasintitulo}{I{0.200}I{0.800}}{\cab{Campo} & \cab{RNF-TRV-11.07}}
Problema & Transversal \\
Requerimiento & El sistema debe publicar el 100 \% de sus servicios alcanzables desde fuera de la red del Cliente exclusivamente a través de una capa de borde que incluya red de distribución de contenidos, cortafuegos de aplicaciones web con reglas gestionadas y personalizadas, y protección contra denegación de servicio distribuida en capas 3, 4 y 7, con 0 servicios expuestos fuera de esa capa. \\
Método de verificación & Escaneo externo de dominios, puertos y servicios contrastado con la superficie de exposición declarada (RT-11.13); pruebas de ataque controlado en capas 3, 4 y 7 y de evasión de la capa de borde, con 0 accesos directos al origen. \\
Origen en las bases & requisito técnico RT-11.07 \\
\end{tablasintitulo}

\begin{tablasintitulo}{I{0.200}I{0.800}}{\cab{Campo} & \cab{RNF-TRV-11.11}}
Problema & Transversal \\
Requerimiento & El sistema debe aplicar, en la puerta de enlace de servicios y sobre el 100 \% de las llamadas a sus servicios, autenticación, autorización, cuotas, límites de tasa, validación de esquema e inspección de carga útil, con 0 llamadas que alcancen un servicio sin haber pasado por esos controles. La cuota es de 200.000 llamadas diarias y el límite de tasa de 50 llamadas por segundo, con ráfaga de 100, por cliente autenticado. \\
Método de verificación & Pruebas de casos negativos por servicio (sin credencial, sin permiso, sobre la cuota, sobre el límite de 50 llamadas por segundo, sobre la cuota diaria de 200.000, esquema inválido, carga útil maliciosa) y revisión de la configuración de la puerta de enlace, con 0 llamadas no controladas. \\
Origen en las bases & requisito técnico RT-11.11; Derivado de RT-11.11 (cuotas y límites de tasa): las cifras concretan esos límites \\
\end{tablasintitulo}

\begin{tablasintitulo}{I{0.200}I{0.800}}{\cab{Campo} & \cab{RNF-TRV-21.11}}
Problema & Transversal \\
Requerimiento & El sistema debe proveer capacitación en línea en modalidad de autoformación para las personas usuarias, con registro de avance por persona disponible para monitoreo, con una disponibilidad mensual del 99,5 \% y una cobertura de al menos un curso para cada uno de los cinco perfiles del plan de capacitación de TRA-009, con registro de avance del 100 \% de las personas inscritas. \\
Método de verificación & Inspección de demostración: completar un curso de autoformación con un usuario de prueba y verificar que el avance queda registrado y es consultable por el Cliente, y medición de la disponibilidad mensual contra 99,5 \% y de la cobertura de cursos por perfil. \\
Origen en las bases & requisito técnico RT-21.11; Derivado de RT-21.11 (capacitación en línea): la disponibilidad de 99,5 \% concreta el requisito \\
\end{tablasintitulo}

\begin{tablasintitulo}{I{0.200}I{0.800}}{\cab{Campo} & \cab{RNF-TRV-21.13}}
Problema & Transversal \\
Requerimiento & El sistema de mesa de ayuda del Adjudicatario debe exponer al Cliente los indicadores clave del proceso de soporte, con acceso a los indicadores específicos según nivel de autorización, y emitir mensualmente el reporte de cumplimiento de los niveles de servicio de atención con detalle por severidad y análisis de los incumplimientos. \\
Método de verificación & Inspección de demostración: verificar visibilidad de indicadores clave con perfil Cliente y restricción de los específicos por perfil; verificar la emisión del reporte mensual por severidad con análisis de incumplimientos. \\
Origen en las bases & requisito técnico RT-21.13; requisito técnico RT-21.18 \\
\end{tablasintitulo}

\begin{tablasintitulo}{I{0.200}I{0.800}}{\cab{Campo} & \cab{RNF-TRV-21.15}}
Problema & Transversal \\
Requerimiento & El sistema de mesa de ayuda del Adjudicatario debe proveer un canal único de registro de incidentes y solicitudes que asigne a cada caso un número de ticket único y una clasificación por severidad, que permita seguir su ciclo de vida completo hasta el cierre conforme, condicionando el cierre a la confirmación de la persona usuaria o al vencimiento del plazo de confirmación automática declarado en la propuesta (7 días corridos desde el aviso de solución, que cubre un ciclo completo de descanso del régimen de turnos 7×7), y conservar un registro histórico de las actividades de soporte con análisis de tendencia mensual. \\
Método de verificación & Prueba de aceptación: crear casos de cada severidad por el canal y verificar número único, clasificación y estados del ciclo hasta el cierre; verificar que un ticket no pasa a cerrado sin confirmación o sin vencer el plazo de 7 días corridos declarado; consultar el registro histórico y el análisis de tendencia mensual. \\
Origen en las bases & requisito técnico RT-21.15; requisito técnico RT-21.17; requisito técnico RT-21.09 \\
\end{tablasintitulo}

\seccion{Requerimientos de transición y preparación}
Estos requerimientos cubren lo que debe ocurrir antes y durante la puesta en producción: migración de datos, capacitación, marcha blanca y preparación de la operación. Cada ficha indica el método con que se verifica su cumplimiento.

\begin{tablasintitulo}{I{0.200}I{0.800}}{\cab{Campo} & \cab{TRA-001}}
Problema & Problema 1 \\
Requerimiento & El sistema debe registrar el tiempo que cada punto de muestreo obligatorio agrega al ciclo productivo y mostrarlo por punto y por turno durante la marcha blanca (período de prueba en paralelo con el método actual), para verificar con datos que no se agregan pasos. El máximo permitido por punto es de 2 segundos en el 95 \% de las muestras, el mismo umbral de las transacciones críticas de terreno. \\
Método de verificación & Medición cronometrada en la marcha blanca del tiempo agregado por cada punto de muestreo, con y sin el paso, informada por punto de muestreo y por turno. Se verifica que el tiempo agregado por cada punto no supera 2 segundos en el percentil 95. \\
Origen en las bases & Bases técnicas del caso Minería, cap. 9.4; Bases técnicas del caso Minería, restricción no negociable 8; Bases técnicas del caso Minería, decisión pendiente 12; supuesto SUP-12 \\
\end{tablasintitulo}

\begin{tablasintitulo}{I{0.200}I{0.800}}{\cab{Campo} & \cab{TRA-002}}
Problema & Problema 1 \\
Requerimiento & El sistema debe contener documentados, antes del cierre de la marcha blanca de la Etapa 1, todos los factores de ajuste de la conciliación metalúrgica vigentes al mes 1 del contrato (acumulados desde 2013), reconstruidos a partir de las planillas existentes y de los cierres históricos, cada uno con su fórmula, su vigencia y su fundamento, de modo que ninguno dependa del conocimiento de una persona en particular. \\
Método de verificación & Recálculo de doce cierres mensuales históricos con los factores reconstruidos, contrastando el resultado contra el cierre original de cada mes. \\
Origen en las bases & Bases técnicas del caso Minería, cap. 4.5; Bases técnicas del caso Minería, criterio de aceptación 2; Bases técnicas del caso Minería, criterio de aceptación 3; Bases administrativas, artículo 17; Bases técnicas del caso Minería, entrevistas; supuesto SUP-02 \\
\end{tablasintitulo}

\begin{tablasintitulo}{I{0.200}I{0.800}}{\cab{Campo} & \cab{TRA-003}}
Problema & Problema 2 \\
Requerimiento & El Adjudicatario debe capacitar a la totalidad del personal de cada área, programando la capacitación en al menos 2 ciclos completos del régimen de turnos 7×7 (dos períodos de 14 días) y capacitando a la dotación de turno de noche en su propio turno. \\
Método de verificación & Revisión del plan de capacitación (programación en al menos 2 ciclos por área y capacitación del turno de noche en su turno) y del registro de cobertura efectiva por área y turno. \\
Origen en las bases & Bases técnicas del caso Minería, cap. 9.4; Bases técnicas del caso Minería, cap. 13.3; requisito técnico RT-22.04; Bases administrativas, artículo 17; Bases administrativas, artículo 90 \\
\end{tablasintitulo}

\begin{tablasintitulo}{I{0.200}I{0.800}}{\cab{Campo} & \cab{TRA-004}}
Problema & Problema 2 \\
Requerimiento & El sistema debe registrar el tiempo que el aviso de cambio de material agrega al ciclo de carguío (la palada) y mostrarlo por equipo y por turno durante la marcha blanca, para acreditar que no se agregan pasos. El umbral exigido es 0 segundos de pala detenida atribuibles al aviso. \\
Método de verificación & Registro del tiempo del aviso en el ciclo de carguío por equipo y por turno durante la marcha blanca, contrastado con cronometraje en terreno de una muestra de ciclos con y sin aviso. \\
Origen en las bases & Bases técnicas del caso Minería, cap. 4.3; Bases técnicas del caso Minería, cap. 9.4; Bases técnicas del caso Minería, restricción no negociable 8 \\
\end{tablasintitulo}

\begin{tablasintitulo}{I{0.200}I{0.800}}{\cab{Campo} & \cab{TRA-005}}
Problema & Problema 2 \\
Requerimiento & El sistema debe medir y publicar la adopción de cada función de registro en terreno por área, por turno y por persona usuaria, como la proporción de eventos registrados sobre los eventos esperados, y compararla con la meta de 95 \% o más en SITCO y de 90 \% en las declaraciones con la botonera de mina (valores fijados por el proponente, no cifras de las bases). Evento esperado es todo hecho operacional que una fuente independiente del registro manual acredita que ocurrió (un ciclo de carguío detectado por la telemetría o el despacho, un viaje de camión, un pesaje en báscula, un ingreso por portería), contado por función, área y turno; el denominador lo genera esa fuente y no la persona usuaria. \\
Método de verificación & Contraste mensual entre eventos registrados y eventos esperados por función, área y turno, con muestra auditada, contra la meta de RNF-USA-005. \\
Origen en las bases & Bases técnicas del caso Minería, cap. 13.3; Bases técnicas del caso Minería, cap. 17.6; Bases técnicas del caso Minería, criterio de aceptación 14; Bases administrativas, artículo 89 \\
\end{tablasintitulo}

\begin{tablasintitulo}{I{0.200}I{0.800}}{\cab{Campo} & \cab{TRA-006}}
Problema & Problema 2 \\
Requerimiento & El sistema debe permitir revertir la puesta en producción de una función de registro en terreno (volver a la versión anterior) por área o por sitio, sin detener la operación ni afectar a las demás áreas, en no más de 30 minutos por función y área desde la decisión de revertir hasta la operación restablecida. Así cabe con holgura en la ventana de baja actividad de 03:00 a 05:00. \\
Método de verificación & Ensayo de reversión de una función en un área durante la marcha blanca, verificando que las demás áreas siguen operando y midiendo el tiempo de reversión contra el máximo de 30 minutos. \\
Origen en las bases & Bases técnicas del caso Minería, cap. 13.3; Bases técnicas del caso Minería, cap. 17.6; requisito técnico RT-20.02; requisito técnico RT-04.06; Derivado de RT-20.02 y de la sección 17.6 (declarar el tiempo de reversión): los 30 minutos son ese tiempo \\
\end{tablasintitulo}

\begin{tablasintitulo}{I{0.200}I{0.800}}{\cab{Campo} & \cab{TRA-007}}
Problema & Problema 2 \\
Requerimiento & El sistema debe registrar la capacitación efectivamente recibida por cada persona usuaria, con área, turno y fecha, y exponer la cobertura alcanzada (personas capacitadas sobre la dotación del área) por área y turno. \\
Método de verificación & Contraste del registro de capacitación contra la nómina por área y turno, verificando la cobertura calculada. \\
Origen en las bases & Bases técnicas del caso Minería, cap. 13.3; requisito técnico RT-22.04; Bases administrativas, artículo 17; Bases administrativas, artículo 90; Derivado de RT-22.04 y de los artículos 17 y 90 de las Bases administrativas (capacitación por turnos): el registro por persona permite demostrar la cobertura \\
\end{tablasintitulo}

\begin{tablasintitulo}{I{0.200}I{0.800}}{\cab{Campo} & \cab{TRA-008}}
Problema & Problema 2 \\
Requerimiento & El sistema debe permitir operar en paralelo con el método actual durante la marcha blanca, registrar ambos flujos y conciliarlos, hasta el cierre conforme según el Artículo 17.3 de las Bases administrativas. \\
Método de verificación & Informe de conciliación diaria de la marcha blanca entre la solución y el registro vigente, con las diferencias explicadas, y acta de apagado del método anterior. \\
Origen en las bases & Bases técnicas del caso Minería, cap. 13.3; Bases técnicas del caso Minería, cap. 17.6; Bases administrativas, artículo 17 \\
\end{tablasintitulo}

\begin{tablasintitulo}{I{0.200}I{0.800}}{\cab{Campo} & \cab{TRA-MIG-001}}
Problema & Transversal \\
Requerimiento & El Proponente debe presentar en su propuesta un plan de migración de datos históricos con alcance, origen, volumen, reglas de transformación, criterios de calidad, estrategia de ejecución y plan de reversión. \\
Método de verificación & Revisión documental del plan de migración contra los siete elementos exigidos (alcance, origen, volumen, transformación, calidad, ejecución, reversión). \\
Origen en las bases & requisito técnico RT-05.11; supuesto SUP-30; supuesto SUP-31 \\
\end{tablasintitulo}

\begin{tablasintitulo}{I{0.200}I{0.800}}{\cab{Campo} & \cab{TRA-MIG-002}}
Problema & Problema 1 \\
Requerimiento & El Adjudicatario debe incluir en la migración una etapa de perfilado y saneamiento previo de los datos de origen, con informe de los defectos detectados y la decisión adoptada sobre cada uno, dejando cada defecto detectado documentado con su decisión antes del primer ensayo de migración. \\
Método de verificación & Revisión del informe de perfilado y saneamiento y verificación de que cada defecto detectado tiene una decisión registrada. \\
Origen en las bases & requisito técnico RT-05.12; supuesto SUP-31 \\
\end{tablasintitulo}

\begin{tablasintitulo}{I{0.200}I{0.800}}{\cab{Campo} & \cab{TRA-MIG-003}}
Problema & Transversal \\
Requerimiento & El Adjudicatario debe ejecutar al menos 2 ensayos completos de migración sobre el ambiente de Preproducción antes de la migración definitiva, midiendo y documentando el tiempo total y el resultado de la conciliación de cada ensayo. \\
Método de verificación & Revisión de las bitácoras de los dos ensayos y contraste del resultado de conciliación de cada uno. \\
Origen en las bases & requisito técnico RT-05.13; supuesto SUP-31 \\
\end{tablasintitulo}

\begin{tablasintitulo}{I{0.200}I{0.800}}{\cab{Campo} & \cab{TRA-MIG-004}}
Problema & Problema 1 \\
Requerimiento & El Adjudicatario debe conciliar la migración definitiva con recuentos y sumas de control sobre el 100 \% de los registros migrados y muestreo dirigido, y dejar explicada toda diferencia antes del cierre de la marcha blanca de la etapa en que migra cada origen. Los requerimientos de migración por dominio (TRA-MIG-005) remiten a esta conciliación y no la repiten. \\
Método de verificación & Recuento y suma de control sobre la totalidad de los registros migrados, más muestreo dirigido sobre los casos de mayor riesgo, contrastando origen contra destino. \\
Origen en las bases & requisito técnico RT-05.14; Bases administrativas, artículo 17; supuesto SUP-31 \\
\end{tablasintitulo}

\begin{tablasintitulo}{I{0.200}I{0.800}}{\cab{Campo} & \cab{TRA-MIG-005}}
Problema & Problema 1 \\
Requerimiento & El Adjudicatario debe migrar los datos históricos de cada dominio, por los años que fija el Caso para cada dominio, conforme al plan de migración de TRA-MIG-001 y con la conciliación de TRA-MIG-004, con la totalidad de los registros de ese período migrada y conciliada, y dejar los registros anteriores a esa ventana consultables en un repositorio durante el período de retención fijado por el Caso (entre 5 y 10 años según el tipo de registro), como establece RF-ADM-001. Particularidades por dominio (ventana de migración fijada por el Caso): - Producción y trazabilidad de mineral: 5 años. Conservar el vínculo entre polígono, viaje, stock, lote de chancado y lote comercial, según la unidad mínima de trazabilidad adoptada (el viaje de camión). - Mantenimiento: 3 años. Incluye órdenes de trabajo y pautas históricas de SAP (SAP no se reemplaza ni se modifica, restricción 4). - Seguridad y salud ocupacional: 10 años. Se tratan como dato sensible (cifrado a nivel de campo, Caso RT-11.10; RT-16.09). Los registros históricos que hoy están en papel deben quedar digitalizados y disponibles para fiscalización (Compañía Minera Altos de Aranda [CMAA], 2026b, cap. 4.8): el alcance es la digitalización como imagen, indexada por tipo de registro, fecha, área, empresa, persona involucrada y folio y sin transcripción de campos, de hasta 200.000 formularios (SUP-32), volumen que equivale a unos 10 años de incidentes, cuasi accidentes y observaciones (≈11.500 al año) e inspecciones planificadas (≈7.800 al año). - Monitoreo ambiental: 10 años. Incluye informes PDF de laboratorio externo (agua, aire y relaves) con trazabilidad a la medición de origen (criterio 11). - Laboratorio: 5 años. Incluye resultados de análisis hoy traspasados a mano a las planillas de balance; el sistema de laboratorio se mantiene e integra en ambos sentidos (restricción 6). \\
Método de verificación & Por cada dominio: recuento y suma de control origen contra destino sobre el 100 \% de los registros de la ventana, más muestreo dirigido (conciliación según TRA-MIG-004); verificación de que los registros fuera de la ventana son consultables en el repositorio, por dominio, durante el período de retención. Para los registros de seguridad digitalizados, verificación de legibilidad e indexación de las imágenes sobre una muestra. \\
Origen en las bases & Bases técnicas del caso Minería, cap. 4.5; Bases técnicas del caso Minería, restricción no negociable 4; Bases técnicas del caso Minería, restricción no negociable 6; Bases técnicas del caso Minería, criterio de aceptación 1; Bases técnicas del caso Minería, criterio de aceptación 11; Bases técnicas del caso Minería, criterio de aceptación 9; requisito técnico RT-05.15; requisito técnico RT-05.10; requisito técnico RT-16.10; requisito técnico RT-05.07; supuesto SUP-30; supuesto SUP-31; supuesto SUP-32 \\
\end{tablasintitulo}

\begin{tablasintitulo}{I{0.200}I{0.800}}{\cab{Campo} & \cab{TRA-009}}
Problema & Transversal \\
Requerimiento & El Adjudicatario debe presentar un plan de capacitación estructurado por perfil (personas usuarias finales, usuarias avanzadas, administradoras, equipo técnico y soporte de niveles 1 y 2); entregar todo el material de capacitación en español, en formato editable y de propiedad del Cliente (manuales por perfil, guías rápidas, preguntas frecuentes, videos tutoriales y base de conocimiento consultable); y evaluar y certificar a las personas usuarias administradoras y al equipo técnico del Cliente antes del cierre de cada marcha blanca, como condición de ese cierre. Se exige la certificación de al menos el 95 \% de las personas usuarias administradoras y del equipo técnico del Cliente (redondeando hacia arriba el número de personas), incluido el personal del turno de noche. \\
Método de verificación & Revisión documental del plan de capacitación contra los perfiles exigidos; revisión del material entregado en cuanto a idioma, formato editable y cesión de propiedad; registro de asistencia y evaluación por persona y acta de certificación. Cálculo del porcentaje certificado por perfil y por turno contra el mínimo de 95 \%. \\
Origen en las bases & requisito técnico RT-22.01; requisito técnico RT-22.03; requisito técnico RT-22.05; Bases administrativas, artículo 90; Bases administrativas, artículo 17; Derivado del artículo 90 de las Bases administrativas y de RT-22.05 (certificación de administradores y equipo técnico): el 95 \% concreta la meta \\
\end{tablasintitulo}

\begin{tablasintitulo}{I{0.200}I{0.800}}{\cab{Campo} & \cab{TRA-010}}
Problema & Problema 2 \\
Requerimiento & El sistema debe permitir declarar la meta de adopción por función y por área y alertar cuando la adopción medida (TRA-005) quede por debajo de la meta durante el número de turnos parametrizado, para activar el plan de acción de adopción de la propuesta. \\
Método de verificación & Prueba de alerta con adopción simulada bajo la meta durante el número parametrizado de turnos, y revisión de que el plan de acción de adopción se activa. \\
Origen en las bases & Bases técnicas del caso Minería, cap. 13.3; Bases técnicas del caso Minería, cap. 17.6 punto 6 (cómo se mide la adopción, con qué meta y qué se hace si no la alcanza); Bases técnicas del caso Minería, criterio de aceptación 14 \\
\end{tablasintitulo}

\seccion{Requerimientos del proyecto}
Los requerimientos del proyecto son obligaciones del Adjudicatario sobre la forma de ejecutar el contrato, como certificaciones, declaraciones y entregables. Cada ficha indica el método con que se verifica su cumplimiento.

\begin{tablasintitulo}{I{0.200}I{0.800}}{\cab{Campo} & \cab{RP-001}}
Problema & Problema 3 \\
Requerimiento & El Adjudicatario debe mantener vigentes las certificaciones ISO 45001 e ISO 14001 y acreditar conocimiento del Reglamento de Seguridad Minera y experiencia comprobable en faena minera en operación. \\
Método de verificación & Revisión de los certificados ISO 45001 e ISO 14001 y de su fecha de vigencia, de la acreditación de conocimiento del Reglamento de Seguridad Minera y de la evidencia de experiencia en faena minera en operación (contratos o certificados de clientes). \\
Origen en las bases & Bases técnicas del caso Minería, restricción no negociable 12; requisito técnico RT-15.02 \\
\end{tablasintitulo}

\begin{tablasintitulo}{I{0.200}I{0.800}}{\cab{Campo} & \cab{RP-002}}
Problema & Problema 2 \\
Requerimiento & Si la propuesta incorpora analítica predictiva, el Adjudicatario debe declarar, en la propuesta, con qué datos se entrena el modelo, la tasa de falsos positivos esperada y el mecanismo de mejora del modelo. \\
Método de verificación & Revisión documental, con el área de mantenimiento del Cliente, de los datos de entrenamiento, de la tasa de falsos positivos declarada y del mecanismo de mejora. \\
Origen en las bases & Bases técnicas del caso Minería, cap. 9.3; requisito técnico RT-18.08; requisito técnico RT-05.30 \\
\end{tablasintitulo}

\begin{tablasintitulo}{I{0.200}I{0.800}}{\cab{Campo} & \cab{RP-003}}
Problema & Transversal \\
Requerimiento & El Adjudicatario debe declarar, componente por componente, qué se ejecuta en la faena y qué en la nube, con el 100 \% de los componentes clasificados, y mantener esa declaración actualizada con cada entrega. \\
Método de verificación & Revisión de la documentación de arquitectura en cada hito de entrega. \\
Origen en las bases & Bases técnicas del caso Minería, cap. 8 (entrevista de Alcaíno); Bases técnicas del caso Minería, cap. 6; Bases técnicas del caso Minería, cap. 9.8; Bases técnicas del caso Minería, cap. 17.4; Bases técnicas del caso Minería, restricción no negociable 2 \\
\end{tablasintitulo}

\begin{tablasintitulo}{I{0.200}I{0.800}}{\cab{Campo} & \cab{RP-004}}
Problema & Transversal \\
Requerimiento & El Adjudicatario debe entregar el estudio de cobertura del rajo, identificando las zonas de sombra (sectores sin señal) y proyectando su evolución a 3 años con el avance del rajo, y mantenerlo actualizado durante el contrato cada seis meses, antes de cada detención programada de planta de marzo y septiembre, y además cuando el plan minero modifique en forma relevante el avance del rajo. \\
Método de verificación & Revisión del estudio y contraste con la medición de cobertura en terreno, y verificación de una actualización cada seis meses. \\
Origen en las bases & Bases técnicas del caso Minería, cap. 6; Bases técnicas del caso Minería, cap. 13.2; requisito técnico RT-03.24; requisito técnico RT-03.23; supuesto SUP-26 \\
\end{tablasintitulo}

\begin{tablasintitulo}{I{0.200}I{0.800}}{\cab{Campo} & \cab{RP-005}}
Problema & Transversal \\
Requerimiento & El Adjudicatario debe levantar y mantener un inventario técnico versionado de los componentes de la solución y de los sistemas con que se integra, con su versión, su responsable y su dependencia, vigente en cada hito de entrega. Debe además entregar, en cada hito, la documentación de arquitectura e integraciones de la solución, versionada y en español. \\
Método de verificación & Revisión documental en cada hito de entrega. \\
Origen en las bases & Bases técnicas del caso Minería, cap. 5; Bases técnicas del caso Minería, cap. 17.5; Bases técnicas del caso Minería, cap. 6; Bases técnicas del caso Minería, restricción no negociable 10; requisito técnico RT-21.04; Bases administrativas, artículo 91; Bases administrativas, artículo 8 \\
\end{tablasintitulo}

\begin{tablasintitulo}{I{0.200}I{0.800}}{\cab{Campo} & \cab{RP-006}}
Problema & Transversal \\
Requerimiento & El Adjudicatario debe ofrecer como servicio costeado toda función de operación de la solución que requiera un especialista dedicado que la compañía no tiene, sin dejar ninguna función crítica dependiente de un perfil inexistente en las 11 personas del equipo de TI. \\
Método de verificación & Revisión del modelo de operación propuesto contra la dotación declarada del Cliente. \\
Origen en las bases & Bases técnicas del caso Minería, cap. 8 (entrevista de Alcaíno); Bases técnicas del caso Minería, restricción no negociable 10 \\
\end{tablasintitulo}

\begin{tablasintitulo}{I{0.200}I{0.800}}{\cab{Campo} & \cab{RP-007}}
Problema & Transversal \\
Requerimiento & El Adjudicatario debe operar un centro de atención en español, 24×7×365, con capacidad de respuesta presencial en faena para los incidentes de severidad crítica, y dimensionar su dotación con fundamento cuantitativo (teoría de colas o Erlang C) a partir del volumen mensual de contactos que estime y declare como supuesto en la propuesta. \\
Método de verificación & Revisión del horario y el idioma de atención; prueba de escalamiento de un incidente crítico simulado con respuesta presencial en faena; revisión de la memoria de cálculo del dimensionamiento (método, volumen de contactos estimado y supuestos) contra los 1.250 contactos mensuales declarados \\
Origen en las bases & requisito técnico RT-21.06; requisito técnico RT-21.07; requisito técnico RT-21.14; supuesto SUP-27 \\
\end{tablasintitulo}

\begin{tablasintitulo}{I{0.200}I{0.800}}{\cab{Campo} & \cab{RP-008}}
Problema & Transversal \\
Requerimiento & El personal del Adjudicatario que asista en faena debe estar acreditado como contratista y contar con examen de altura geográfica vigente (faena a 3.180 metros), y el Adjudicatario debe incluir en la oferta el costo y el tiempo de traslado de su personal desde Antofagasta (145 km por Ruta 25 y camino privado). \\
Método de verificación & Revisión de las acreditaciones y de los exámenes de altura vigentes del personal en faena; revisión del plan de dotación, de los tiempos de traslado declarados y del costo de traslado en la oferta. \\
Origen en las bases & Bases técnicas del caso Minería, cap. 2.1; Bases técnicas del caso Minería, cap. 17.5; Bases técnicas del caso Minería, restricción no negociable 12; requisito técnico RT-21.16; supuesto SUP-35 \\
\end{tablasintitulo}

\begin{tablasintitulo}{I{0.200}I{0.800}}{\cab{Campo} & \cab{RP-009}}
Problema & Transversal \\
Requerimiento & El Adjudicatario debe documentar los procedimientos de operación críticos de la solución (arranque, monitoreo, respaldo, restauración, conmutación, escalamiento y atención de fallas de integración) y lograr que el personal del Cliente los ejecute sin asistencia antes del cierre de la marcha blanca, de modo que el equipo de 11 personas de TI pueda ejecutarlos sin apoyo del Adjudicatario. \\
Método de verificación & Prueba de ejecución de cada procedimiento crítico por personal del Cliente sin asistencia, observada y registrada. \\
Origen en las bases & Bases técnicas del caso Minería, cap. 8 (entrevista de Alcaíno); requisito técnico RT-21.04; requisito técnico RT-07.05; requisito técnico RT-22.05; Bases administrativas, artículo 91; Bases administrativas, artículo 17 \\
\end{tablasintitulo}

\begin{tablasintitulo}{I{0.200}I{0.800}}{\cab{Campo} & \cab{RP-010}}
Problema & Transversal \\
Requerimiento & El Adjudicatario debe proveer acompañamiento y mentoría al equipo técnico del Cliente durante al menos 6 meses posteriores al paso a producción de la Etapa 2; ejecutar al menos 2 jornadas anuales de actualización para el personal del Cliente durante la Operación; capacitar sin costo adicional al personal nuevo del Cliente durante la Operación; y mantener y actualizar la base de conocimiento durante todo el Contrato, con métricas de uso y de utilidad percibida. \\
Método de verificación & Registro de las sesiones de acompañamiento y mentoría (fecha, participantes y temas); actas de las jornadas anuales de actualización; registro de capacitación del personal nuevo y verificación de que no se cobró cargo adicional; revisión de la base de conocimiento y de sus métricas. \\
Origen en las bases & Bases técnicas del caso Minería, cap. 8 (entrevista de Alcaíno); requisito técnico RT-22.06; requisito técnico RT-22.07; requisito técnico RT-22.08; requisito técnico RT-21.10 \\
\end{tablasintitulo}

\begin{tablasintitulo}{I{0.200}I{0.800}}{\cab{Campo} & \cab{RP-011}}
Problema & Transversal \\
Requerimiento & El Adjudicatario debe planificar las validaciones que requieran al personal clave de operaciones (planta y mina) fuera de las detenciones programadas de marzo y septiembre, e indicar en su cronograma la fecha de cada validación. \\
Método de verificación & Revisión del cronograma de validaciones contra las fechas de las detenciones programadas de marzo y septiembre. \\
Origen en las bases & Bases técnicas del caso Minería, restricción no negociable 11 (cap. 10) \\
\end{tablasintitulo}

\anexo{Listado de Supuestos, Exclusiones y Restricciones}
Este anexo detalla los 35 supuestos declarados en los puntos que el caso no define, con su fundamento, su impacto si resultan equivocados y la instancia en que se validan. Incluye además las exclusiones y las restricciones no negociables que fijan los capítulos 11 y 10 de las Bases técnicas del caso Minería. Complementa la sección 2.5 del Subdocumento 2.

\seccion{Supuestos}
Los supuestos se numeran de SUP-01 a SUP-35. Los SUP-01 a SUP-14 responden a las 14 decisiones pendientes del numeral 16.1 del caso, páginas 26 y 27. Los SUP-20 a SUP-25 son vacíos que identificó el proponente y los restantes son supuestos del alcance. Cada ficha indica el tipo y el tema del supuesto, su número de decisión cuando corresponde, la decisión o supuesto adoptado, su fundamento, su impacto si resulta equivocado, la instancia de validación, el módulo responsable y si además se trata como riesgo.

\begin{tablasintitulo}{I{0.200}I{0.800}}{\cab{Campo} & \cab{SUP-01}}
Tipo & Decisión del numeral 16.1 \\
Tema & Unidad mínima de trazabilidad \\
N.º de decisión & 1 \\
Decisión o supuesto adoptado & La unidad es el viaje de camión en la mina, el lote de turno en la planta, el camión a puerto y el lote comercial. \\
Fundamento & Son hechos que hoy ya registra algún sistema: el despacho registra cada viaje, la planta trabaja por turno y la báscula pesa cada camión a puerto. Una unidad más fina exigiría nuevos registros en el ciclo productivo, y con estas unidades se reconstruye el origen de un lote en menos de 2 horas. \\
Impacto si resulta equivocado & Si el comprador o el esquema de cadena de custodia exige una unidad más fina, habría que agregar registros en terreno y rehacer el modelo de la cadena, y el criterio de aceptación del origen de un lote podría no cumplirse en el mes 16. \\
Instancia de validación & Levantamiento de los meses 1 a 3, con planificación, metalurgia y el área comercial. Se confirma con el comprador en la revisión del lote piloto en 2028. \\
Módulo responsable & Trazabilidad y Embarque \\
¿Se trata también como riesgo? & No \\
\end{tablasintitulo}

\begin{tablasintitulo}{I{0.200}I{0.800}}{\cab{Campo} & \cab{SUP-02}}
Tipo & Decisión del numeral 16.1 \\
Tema & Fuente de verdad del tonelaje \\
N.º de decisión & 2 \\
Decisión o supuesto adoptado & La balanza de alimentación es la fuente oficial del mineral que entra a la planta, y el pesaje de a bordo es la fuente de la mina, corregido con el sesgo estimado por camión a partir de la báscula, con intervalo de confianza de 95 \%. \\
Fundamento & Hoy existen tres mediciones que no coinciden y ninguna tiene un rol definido. Dar a cada medición un tramo convierte la diferencia entre ellas en un dato por tramo, y la báscula ya pesa una muestra de los viajes, lo que permite calcular el factor sin equipos nuevos. \\
Impacto si resulta equivocado & Si finanzas o metalurgia no la ratifican, cambia la base del balance, se rehacen las reglas de Conciliación Metalúrgica y peligran el cierre en 3 días hábiles y las metas de diferencia que compromete el proponente (bajo 3 \% y bajo 2 \%; la referencia de la industria es bajo 2 \%). \\
Instancia de validación & Ratificación de la gerencia de finanzas en el mes 3, junto con la superintendencia de planta y metalurgia. \\
Módulo responsable & Conciliación Metalúrgica \\
¿Se trata también como riesgo? & Sí \\
\end{tablasintitulo}

\begin{tablasintitulo}{I{0.200}I{0.800}}{\cab{Campo} & \cab{SUP-03}}
Tipo & Decisión del numeral 16.1 \\
Tema & Mezcla en stocks \\
N.º de decisión & 3 \\
Decisión o supuesto adoptado & Cada día del stock es una capa con su tonelaje, su ley y sus orígenes, y el material retomado toma el promedio ponderado de las capas vigentes. El resultado se compara cada mes con la ley de alimentación medida por el laboratorio y se verifica con una campaña de muestreo de retoma en la marcha blanca. \\
Fundamento & En los stocks intermedios es donde hoy se rompe la trazabilidad y no existe una medición del material retomado. El modelo de capas no requiere sensores nuevos, y la comparación mensual lo vuelve verificable. \\
Impacto si resulta equivocado & Si el stock no se comporta como capas, el origen asignado a los lotes sería impreciso y la diferencia de conciliación no quedaría bien explicada por tramo. Habría que ajustar el modelo, por ejemplo con capas por frente de descarga. \\
Instancia de validación & Levantamiento de los meses 1 a 3 con planificación y metalurgia, campaña de muestreo de retoma y comparación mensual desde la marcha blanca de la Etapa 1. El modelo se ajusta con un año de datos antes del mes 28. \\
Módulo responsable & Trazabilidad y Embarque \\
¿Se trata también como riesgo? & No \\
\end{tablasintitulo}

\begin{tablasintitulo}{I{0.200}I{0.800}}{\cab{Campo} & \cab{SUP-04}}
Tipo & Decisión del numeral 16.1 \\
Tema & Identificación del operador \\
N.º de decisión & 4 \\
Decisión o supuesto adoptado & El operador de pala se identifica acercando su credencial personal al lector de la botonera en el relevo, sin login manual, sin contraseña y sin pantallas adicionales, en no más de 2 segundos en el percentil 95. En camiones, se usa el ingreso del operador al despacho. Si nadie se identifica, el registro queda como operador no identificado y se avisa al supervisor de turno. \\
Fundamento & La credencial ya existe y el despacho ya identifica al operador de camión. Un toque por relevo no agrega pasos al ciclo productivo (restricción 8 del Caso) y deja identificado quién declaró cada cambio. \\
Impacto si resulta equivocado & Si los operadores no tocan la credencial o el lector no la reconoce, los registros quedan sin autor. Habría que agregar un aviso en la botonera o cambiar el lector. \\
Instancia de validación & Piloto con una pala en la marcha blanca de la Etapa 1, aprobado por la gerencia de operaciones mina. \\
Módulo responsable & Operación Mina \\
¿Se trata también como riesgo? & No \\
\end{tablasintitulo}

\begin{tablasintitulo}{I{0.200}I{0.800}}{\cab{Campo} & \cab{SUP-05}}
Tipo & Decisión del numeral 16.1 \\
Tema & Definición de detención \\
N.º de decisión & 5 \\
Decisión o supuesto adoptado & Se usa una taxonomía única de demoras y detenciones, acordada con mina y planta. El sistema clasifica automáticamente a partir de la telemetría, el despacho y el historiador de proceso, y el supervisor del turno clasifica solo el residuo que queda sin resolver. \\
Fundamento & Hoy el 23 \% del tiempo registrado queda sin clasificar y el 61 \% de las detenciones no tiene causa raíz. El supervisor conoce la causa de lo que el sistema no resuelve y puede clasificarlo dentro del turno sin cargar al operador. \\
Impacto si resulta equivocado & Si la taxonomía no refleja la operación o el supervisor no alcanza a clasificar, quedan detenciones pendientes y no se cumple el 100 \% clasificado en el turno siguiente. Habría que rediseñar la taxonomía o reasignar la tarea. \\
Instancia de validación & Diseño con supervisores y despachadores durante el levantamiento. Se mide en la marcha blanca de la Etapa 1 para la mina y de la Etapa 2 para la planta. \\
Módulo responsable & Operación Mina y Equipos y Mantenimiento \\
¿Se trata también como riesgo? & No \\
\end{tablasintitulo}

\begin{tablasintitulo}{I{0.200}I{0.800}}{\cab{Campo} & \cab{SUP-06}}
Tipo & Decisión del numeral 16.1 \\
Tema & Muestras rechazadas \\
N.º de decisión & 6 \\
Decisión o supuesto adoptado & Un lote cuyo análisis es rechazado queda retenido hasta tener un análisis válido, y se pide una nueva toma de muestra con un plazo de 24 horas desde el rechazo. \\
Fundamento & Un lote sin análisis válido no tiene una ley confiable y no puede embarcarse (CMAA, 2026b, cap. 16.1). Conformarlo así expone a penalizaciones como las que hoy recibe la compañía: 4 de 31 embarques del año anterior, por arsénico. \\
Impacto si resulta equivocado & Si retener el lote atrasa un embarque, o el cliente prefiere liberar con un análisis provisorio, habría que cambiar la regla de retención y aceptar el riesgo de penalización. \\
Instancia de validación & Levantamiento de los meses 1 a 3 con metalurgia, el laboratorio y el área comercial. Se prueba en la marcha blanca de la Etapa 1. \\
Módulo responsable & Trazabilidad y Embarque \\
¿Se trata también como riesgo? & No \\
\end{tablasintitulo}

\begin{tablasintitulo}{I{0.200}I{0.800}}{\cab{Campo} & \cab{SUP-07}}
Tipo & Decisión del numeral 16.1 \\
Tema & Concentrado devuelto \\
N.º de decisión & 7 \\
Decisión o supuesto adoptado & El concentrado devuelto se reincorpora como lote de reproceso ligado a su lote original. \\
Fundamento & Así la cadena no pierde el origen y el tonelaje devuelto no se cuenta dos veces ni desaparece del balance. \\
Impacto si resulta equivocado & Si el cliente registra la devolución como un lote nuevo sin vínculo, el origen no se puede recorrer hacia atrás y la conciliación queda con diferencias sin explicar. \\
Instancia de validación & Levantamiento de los meses 1 a 3 con el área comercial y metalurgia. Se prueba en la marcha blanca de la Etapa 1. \\
Módulo responsable & Trazabilidad y Embarque \\
¿Se trata también como riesgo? & No \\
\end{tablasintitulo}

\begin{tablasintitulo}{I{0.200}I{0.800}}{\cab{Campo} & \cab{SUP-08}}
Tipo & Decisión del numeral 16.1 \\
Tema & Autenticación de contratistas \\
N.º de decisión & 8 \\
Decisión o supuesto adoptado & En terreno, el trabajador contratista se identifica con su credencial y un PIN numérico. En el portal, la empresa contratista entra con doble factor. \\
Fundamento & El cliente exige doble factor y prohíbe las credenciales compartidas. Una contraseña alfanumérica no es práctica en terreno con guantes, y la credencial ya existe. \\
Impacto si resulta equivocado & Si el cliente no acepta la credencial más el PIN como doble factor, habría que agregar otro factor en terreno, lo que suma tiempo y dispositivos. \\
Instancia de validación & Revisión con el área de tecnologías de información y el área de seguridad del cliente durante el levantamiento. Se confirma en la capacitación por empresa antes de la marcha blanca. \\
Módulo responsable & Acreditación y Acceso \\
¿Se trata también como riesgo? & No \\
\end{tablasintitulo}

\begin{tablasintitulo}{I{0.200}I{0.800}}{\cab{Campo} & \cab{SUP-09}}
Tipo & Decisión del numeral 16.1 \\
Tema & Zona de sombra con cambio de turno \\
N.º de decisión & 9 \\
Decisión o supuesto adoptado & La botonera y la tablet guardan cifrado hasta 48 horas continuas sin señal, con la identidad de cada operador, y lo envían solas al recuperar señal, con sincronización en no más de 4 horas y sin intervención manual. \\
Fundamento & La red LTE no cubre todo el rajo y un operador puede terminar su turno en una zona sin señal y ser reemplazado por otro. Así cada registro queda con su autor. \\
Impacto si resulta equivocado & Si una zona de sombra dura más de un turno o el dispositivo pierde registros, se pierde el origen y el autor de esos cambios. Habría que aumentar la capacidad del dispositivo o ampliar la cobertura. \\
Instancia de validación & Estudio de cobertura LTE en la Etapa 1, prueba de corte controlado de 48 horas en la detención de marzo de 2028 y piloto con una pala. \\
Módulo responsable & Operación Mina \\
¿Se trata también como riesgo? & No \\
\end{tablasintitulo}

\begin{tablasintitulo}{I{0.200}I{0.800}}{\cab{Campo} & \cab{SUP-10}}
Tipo & Decisión del numeral 16.1 \\
Tema & Dueño de la telemetría \\
N.º de decisión & 10 \\
Decisión o supuesto adoptado & Se asume que el cliente puede extraer la telemetría de sus equipos, y la especificación de acceso con cada fabricante parte en el mes 1 y la negociación de los acuerdos, en el mes 3, usando el estándar abierto de intercambio de datos de equipos mineros (ISO 15143-3) cuando el fabricante lo soporte, conforme a RT-05.23. \\
Fundamento & La propiedad de esos datos hoy no está clara, y Equipos y Mantenimiento se deja en la Etapa 2 para dar tiempo a la negociación. \\
Impacto si resulta equivocado & Si un fabricante no entrega los datos o los cobra, la vista única queda incompleta para sus equipos y habría que reordenar la integración o su costo. \\
Instancia de validación & Negociación con cada fabricante desde el mes 3, cerrada antes del inicio de la Etapa 2 en el mes 13. \\
Módulo responsable & Equipos y Mantenimiento \\
¿Se trata también como riesgo? & Sí \\
\end{tablasintitulo}

\begin{tablasintitulo}{I{0.200}I{0.800}}{\cab{Campo} & \cab{SUP-11}}
Tipo & Decisión del numeral 16.1 \\
Tema & Carga en borde de polígono \\
N.º de decisión & 11 \\
Decisión o supuesto adoptado & La carga en el contacto entre dos polígonos se asigna íntegra al polígono donde está posicionada la pala, y el viaje queda marcado como borde de forma automática, a partir de la posición de la pala respecto del contacto que informa el despacho, sin pedir ninguna acción al operador. \\
Fundamento & Hoy todo se asigna al polígono de la pala y nadie ha decidido mantenerlo o corregirlo. El operador ve el contacto pero no tiene que decidir a qué polígono pertenece la carga, y la marca permite tratar esos viajes aparte en la conciliación. \\
Impacto si resulta equivocado & Si los viajes de borde son muchos, la ley asignada a esos viajes es imprecisa y afecta la conciliación. Habría que definir una regla de reparto entre los dos polígonos. \\
Instancia de validación & Levantamiento con planificación. Se revisa la proporción de viajes marcados como borde en la marcha blanca de la Etapa 1. \\
Módulo responsable & Operación Mina \\
¿Se trata también como riesgo? & No \\
\end{tablasintitulo}

\begin{tablasintitulo}{I{0.200}I{0.800}}{\cab{Campo} & \cab{SUP-12}}
Tipo & Decisión del numeral 16.1 \\
Tema & Muestreo versus ciclo de mina \\
N.º de decisión & 12 \\
Decisión o supuesto adoptado & El muestreo se hace en los puntos de transferencia de la planta, con los muestreadores y el procedimiento del laboratorio, y la mina solo aporta un accionador. El tiempo que cada punto de muestreo agrega al ciclo no supera 2 segundos en el percentil 95. \\
Fundamento & La mina no puede perder rendimiento de pala y ningún registro nuevo puede agregar pasos al ciclo productivo. La planta ya cuenta con muestreadores y procedimiento. \\
Impacto si resulta equivocado & Si la planta necesita muestras tomadas en la mina, habría que agregar pasos al ciclo productivo, lo que va contra una restricción del cliente, y rediseñar el muestreo. \\
Instancia de validación & Acuerdo entre la gerencia de operaciones mina y la superintendencia de planta durante el levantamiento. \\
Módulo responsable & Conciliación Metalúrgica \\
¿Se trata también como riesgo? & No \\
\end{tablasintitulo}

\begin{tablasintitulo}{I{0.200}I{0.800}}{\cab{Campo} & \cab{SUP-13}}
Tipo & Decisión del numeral 16.1 \\
Tema & Biometría y sindicato \\
N.º de decisión & 13 \\
Decisión o supuesto adoptado & La verificación facial es uno a uno contra la persona asociada a la credencial presentada, y entra solo con acuerdo previo del sindicato. Si la persona no tiene plantilla enrolada o el lector facial está fuera de servicio, el ingreso se controla con la credencial y la validación de habilitación. La plantilla se guarda cifrada y se elimina cuando la persona pierde el acceso. Este supuesto rige solo para la portería: el acceso a la sala técnica usa biometría facial con respaldo por huella dactilar (AFIS), como exige RT-06.20. \\
Fundamento & La verificación facial resuelve el préstamo de credenciales, pero el sindicato objeta la biometría y los datos biométricos requieren el resguardo de datos sensibles de la Ley N° 21.719 (las bases citan esa ley para los datos de salud ocupacional, fatiga, alcohol y drogas). Por eso se deja en la Etapa 2. \\
Impacto si resulta equivocado & Si no hay acuerdo, la portería funciona con credencial y validación de habilitación, y el préstamo de credenciales no se resuelve. \\
Instancia de validación & Mesa de trabajo con el sindicato, seguridad y recursos humanos desde el mes 2, con un protocolo escrito antes de instalar los lectores faciales en la Etapa 2. \\
Módulo responsable & Acreditación y Acceso \\
¿Se trata también como riesgo? & Sí \\
\end{tablasintitulo}

\begin{tablasintitulo}{I{0.200}I{0.800}}{\cab{Campo} & \cab{SUP-14}}
Tipo & Decisión del numeral 16.1 \\
Tema & Versión del plan \\
N.º de decisión & 14 \\
Decisión o supuesto adoptado & Cada cambio del plan crea una versión con fecha y hora. El cumplimiento se conoce al inicio del turno siguiente, comparado con la versión del plan efectivamente vigente. \\
Fundamento & Hoy el plan cambia sin registro, por ejemplo cuando la tronadura desplaza el contacto, y el cumplimiento se mide con diez a quince días de desfase. \\
Impacto si resulta equivocado & Si los planificadores no registran los cambios, el cumplimiento se mide contra un plan que ya no regía. Habría que tomar las versiones directo del software de planificación. \\
Instancia de validación & La superintendencia de planificación valida el versionado en la marcha blanca de la Etapa 1. \\
Módulo responsable & Operación Mina \\
¿Se trata también como riesgo? & No \\
\end{tablasintitulo}

\begin{tablasintitulo}{I{0.200}I{0.800}}{\cab{Campo} & \cab{SUP-15}}
Tipo & Supuesto del alcance \\
Tema & Fecha de inicio del contrato \\
N.º de decisión & no aplica \\
Decisión o supuesto adoptado & El mes 1 del contrato es enero de 2027 y no es posterior a abril de 2027. \\
Fundamento & Es la fecha estimada de inicio: las bases cuentan el mes 1 como el primer mes completo desde la fecha de inicio del contrato (CMAA, 2026a, Art. 10, numeral 2). Con el mes 1 en enero de 2027, la Etapa 1 entra en producción en abril de 2028. La Etapa 2 entra en producción en septiembre de 2028. La Etapa 1 acumula nueve meses de operación antes de la exigencia del comprador de enero de 2029 (CMAA, 2026b, hitos externos). \\
Impacto si resulta equivocado & Con un inicio posterior a abril de 2027, la Etapa 2 y el Portal para Compradores entran en producción a partir de enero de 2029 y habría que adelantar a la Etapa 1 la huella de carbono y el portal del comprador, con más dotación y más riesgo sobre la marcha blanca. \\
Instancia de validación & Confirmación de la fecha de inicio con el Cliente en la presentación preparatoria 3 y firma del contrato. \\
Módulo responsable & Todos los módulos \\
¿Se trata también como riesgo? & No \\
\end{tablasintitulo}

\begin{tablasintitulo}{I{0.200}I{0.800}}{\cab{Campo} & \cab{SUP-16}}
Tipo & Supuesto del alcance \\
Tema & Lectura de sistemas existentes \\
N.º de decisión & no aplica \\
Decisión o supuesto adoptado & El sistema de despacho de flota, SAP S/4HANA, el sistema de gestión de laboratorio y el historiador de proceso exponen interfaces documentadas (servicios, OPC UA o archivos estructurados) que permiten leer y, en SAP y en el laboratorio, también recibir los datos que envía SITCO, sin modificar esos sistemas. SAP sigue emitiendo los documentos contables. El Cliente entrega su documentación y habilita el acceso a un ambiente de prueba dentro de los 20 días hábiles siguientes al inicio del Contrato. Si un sistema no expone interfaz, se usa una exportación automática programada; si ni eso es posible, la capa intermedia o el servicio adicional de su proveedor se trata como cambio conforme al Artículo 72.º de las Bases administrativas. \\
Fundamento & Los cinco sistemas se mantienen y deben integrarse (CMAA, 2026b, cap. 5 y restricciones 3 a 7), pero el Cliente declara no tener documentadas sus interfaces, versiones ni condiciones de acceso. Veinte días hábiles es el plazo que permite cerrar el levantamiento antes del hito H2 (mes 4). \\
Impacto si resulta equivocado & Con una exportación programada los datos llegan con menos frecuencia, lo que puede atrasar el balance del turno. Si alguna interfaz no existe o no está documentada, hay que construir una capa intermedia no prevista y las integraciones críticas de la Etapa 1 se atrasan. \\
Instancia de validación & Levantamiento del inventario técnico (RP-005), entregado en el mes 3, y levantamiento técnico de interfaces en los meses 1 y 3. \\
Módulo responsable & Integrador de Sistemas \\
¿Se trata también como riesgo? & Sí \\
\end{tablasintitulo}

\begin{tablasintitulo}{I{0.200}I{0.800}}{\cab{Campo} & \cab{SUP-17}}
Tipo & Supuesto del alcance \\
Tema & Credencial actual \\
N.º de decisión & no aplica \\
Decisión o supuesto adoptado & La credencial actual es de proximidad y se puede leer con el lector de la botonera y de la tablet. Si no es así, se cambia el lector, no la credencial. \\
Fundamento & La credencial ya está en manos de todo el personal y cambiarla afectaría a toda la faena. \\
Impacto si resulta equivocado & Habría que cambiar el lector de las botoneras y las tablets, con efecto en su especificación y su costo. \\
Instancia de validación & Levantamiento de los meses 1 a 3. \\
Módulo responsable & Operación Mina y Acreditación y Acceso \\
¿Se trata también como riesgo? & No \\
\end{tablasintitulo}

\begin{tablasintitulo}{I{0.200}I{0.800}}{\cab{Campo} & \cab{SUP-18}}
Tipo & Supuesto del alcance \\
Tema & Alimentación en cabinas de pala \\
N.º de decisión & no aplica \\
Decisión o supuesto adoptado & Las cabinas de las palas tienen alimentación de 24 volts disponible para la botonera. \\
Fundamento & La botonera se instala junto a la pantalla del despacho, que ya funciona con la alimentación de la cabina. \\
Impacto si resulta equivocado & Habría que agregar un conversor o una alimentación nueva, instalada con mantenimiento en una ventana de intervención, lo que suma costo y plazo. \\
Instancia de validación & Inspección de las cinco palas en faena durante el levantamiento. \\
Módulo responsable & Operación Mina \\
¿Se trata también como riesgo? & No \\
\end{tablasintitulo}

\begin{tablasintitulo}{I{0.200}I{0.800}}{\cab{Campo} & \cab{SUP-19}}
Tipo & Supuesto del alcance \\
Tema & Réplica del historiador de proceso \\
N.º de decisión & no aplica \\
Decisión o supuesto adoptado & El historiador de proceso permite replicar las señales seleccionadas hacia un servidor espejo en la zona desmilitarizada. \\
Fundamento & Así se leen las mediciones de la planta sin escribir nunca en la red de control. \\
Impacto si resulta equivocado & Si no lo permite, el tonelaje de la balanza de alimentación y las detenciones de planta tendrían que leerse de otra forma, con efecto en el balance del turno. \\
Instancia de validación & Levantamiento con la superintendencia de planta, que define las señales que se leen. \\
Módulo responsable & Conciliación Metalúrgica y Equipos y Mantenimiento \\
¿Se trata también como riesgo? & No \\
\end{tablasintitulo}

\begin{tablasintitulo}{I{0.200}I{0.800}}{\cab{Campo} & \cab{SUP-20}}
Tipo & Vacío identificado por el proponente \\
Tema & Fuente de verdad de la ley medida \\
N.º de decisión & no aplica \\
Decisión o supuesto adoptado & El resultado del sistema de gestión de laboratorio es la fuente de verdad de la ley medida. El modelo de bloques y el plan son estimaciones que se reconcilian contra ella, y la diferencia entre ley estimada y medida se publica por polígono como indicador de calidad del modelo. \\
Fundamento & El numeral 16.1 n.º 2 pregunta por el tonelaje, pero la conciliación metalúrgica es de metal contenido: tonelaje por ley (glosario, Cap. C). La ley también tiene tres fuentes que discrepan (Cap. 8, Sanhueza: 0,58 del plan contra 0,51 de cabeza). De la ley depende el contenido de arsénico, que es el origen comercial de esta licitación. \\
Impacto si resulta equivocado & La conciliación cubre el tonelaje pero no el metal contenido, y la diferencia de ley queda sin responsable ni explicación, con lo que el criterio n.º 3 se cumple solo a medias. \\
Instancia de validación & Recálculo retroactivo de doce meses contrastando ley estimada contra ley medida por polígono, en una instancia de validación distinta de la del SUP-02. Se realiza en el mes 15. \\
Módulo responsable & Conciliación Metalúrgica \\
¿Se trata también como riesgo? & No \\
\end{tablasintitulo}

\begin{tablasintitulo}{I{0.200}I{0.800}}{\cab{Campo} & \cab{SUP-21}}
Tipo & Vacío identificado por el proponente \\
Tema & Material anterior al sistema \\
N.º de decisión & no aplica \\
Decisión o supuesto adoptado & El material acopiado antes del arranque del sistema se marca como “origen anterior al sistema”, con la mejor estimación disponible y esa condición visible en toda consulta. Se declara el plazo durante el cual la trazabilidad será parcial y se compromete su extinción antes de enero de 2029. \\
Fundamento & El criterio n.º 1 exige reconstruir el origen de cualquier lote, pero los lotes de los primeros meses contendrán material sin origen registrado; el acopio de puerto es de 25.000 toneladas (Cap. 3). Desde enero de 2029 el comprador exige trazabilidad verificable (Cap. 13.2). \\
Impacto si resulta equivocado & Si la rotación real del acopio es más lenta que la estimada, queda material sin origen después de enero de 2029 y el compromiso con el comprador no se cumple para esos lotes. \\
Instancia de validación & Medir la rotación efectiva del acopio durante la marcha blanca y proyectar la fecha de extinción del material sin origen; si excede el hito, se acuerda con el Cliente un tratamiento específico para el remanente. \\
Módulo responsable & Trazabilidad y Embarque \\
¿Se trata también como riesgo? & Sí \\
\end{tablasintitulo}

\begin{tablasintitulo}{I{0.200}I{0.800}}{\cab{Campo} & \cab{SUP-22}}
Tipo & Vacío identificado por el proponente \\
Tema & Certificado para el comprador \\
N.º de decisión & no aplica \\
Decisión o supuesto adoptado & El portal de trazabilidad entrega al comprador un certificado derivado —origen por sector, fechas, stocks recorridos, análisis y cadena de custodia— y no la cadena completa con ley por polígono. La publicación requiere aprobación del área comercial y toda consulta queda registrada. \\
Fundamento & RT-16.30 exige el portal pero no define qué se expone. La ley por polígono revela la geología del yacimiento, que es información competitivamente sensible. El Cap. 9.1 distingue tres destinatarios sin declarar que deben ver cosas distintas; esta decisión lo declara. \\
Impacto si resulta equivocado & Si el comprador exige el detalle completo como condición del contrato de 2029, hay que renegociar qué se expone o aceptar publicar información sensible del yacimiento. \\
Instancia de validación & Validación con el Cliente en la presentación preparatoria 3 y, por su intermedio, del contenido del certificado con el comprador antes del cierre del diseño del portal. \\
Módulo responsable & Trazabilidad y Embarque \\
¿Se trata también como riesgo? & Sí \\
\end{tablasintitulo}

\begin{tablasintitulo}{I{0.200}I{0.800}}{\cab{Campo} & \cab{SUP-23}}
Tipo & Vacío identificado por el proponente \\
Tema & Vector de elementos del lote \\
N.º de decisión & no aplica \\
Decisión o supuesto adoptado & El modelo de lote lleva un vector de elementos —cobre, oro, plata y arsénico como elemento deletéreo— y no un único valor de ley de cobre. La conciliación y la liquidación se calculan sobre ese vector. \\
Fundamento & El producto de Aranda es concentrado de cobre con oro y plata como subproductos (Cap. 2.1), que se liquidan, se pagan y tienen sus propias deducciones (Cap. 12), pero el resto del caso habla sólo de cobre. Un modelo que solo represente la ley de cobre no soporta la liquidación real. \\
Impacto si resulta equivocado & Si el vector se define incompleto, la liquidación del embarque no se puede reproducir desde el sistema y queda dependiendo de la planilla comercial, que es el problema que la propuesta viene a resolver. \\
Instancia de validación & Contrastar el vector modelado contra tres liquidaciones históricas reales del comprador durante la marcha blanca. \\
Módulo responsable & Trazabilidad y Embarque \\
¿Se trata también como riesgo? & No \\
\end{tablasintitulo}

\begin{tablasintitulo}{I{0.200}I{0.800}}{\cab{Campo} & \cab{SUP-24}}
Tipo & Vacío identificado por el proponente \\
Tema & Operador portuario \\
N.º de decisión & no aplica \\
Decisión o supuesto adoptado & Se diseña para no depender del operador portuario: gabinete de borde con operación desconectada en el acopio de puerto, con el mismo esquema de registro local y reconciliación del SUP-09. Se propone al Cliente negociar un acuerdo de nivel de servicio como actividad temprana. \\
Fundamento & El acopio y el embarque ocurren en recinto de terceros, con conectividad provista por el operador portuario y sin acuerdo de nivel de servicio (Cap. 3 y 6). Ahí se conforma el lote comercial, que es el objeto último de la trazabilidad. \\
Impacto si resulta equivocado & Si el operador portuario además restringe la instalación física del gabinete, no hay registro autónomo en el punto donde se conforma el lote y la cadena se corta justo al final. \\
Instancia de validación & Validación con el Cliente, en la presentación preparatoria 3, de las condiciones del contrato con el operador portuario y visita técnica al recinto antes del diseño de detalle. \\
Módulo responsable & Trazabilidad y Embarque \\
¿Se trata también como riesgo? & Sí \\
\end{tablasintitulo}

\begin{tablasintitulo}{I{0.200}I{0.800}}{\cab{Campo} & \cab{SUP-25}}
Tipo & Vacío identificado por el proponente \\
Tema & Fuente horaria común \\
N.º de decisión & no aplica \\
Decisión o supuesto adoptado & Se define una fuente horaria común y trazable para faena, Centro Integrado de Operaciones y puerto, con sincronización periódica de los dispositivos y registro de la deriva del reloj local en cada evento generado en modo desconectado. \\
Fundamento & Sin fuente horaria común los sellos de tiempo de los tres sitios no son comparables y la evidencia pierde valor probatorio. Es crítico porque el SUP-09 obliga a sellar en el dispositivo y porque se exige firma electrónica avanzada para las actas de aceptación de embarque y los registros de seguridad. \\
Impacto si resulta equivocado & Los eventos de los tres sitios no son ordenables entre sí y la firma electrónica avanzada pierde sustento, con lo que el expediente de trazabilidad deja de ser prueba ante la autoridad y ante el comprador. \\
Instancia de validación & Prueba de deriva de reloj sobre el corte controlado de 48 horas del SUP-09: la deriva registrada debe permitir ordenar los eventos de los tres sitios sin ambigüedad. Se realiza en la marcha blanca de la Etapa 1. \\
Módulo responsable & Integrador de Sistemas y Trazabilidad y Embarque \\
¿Se trata también como riesgo? & No \\
\end{tablasintitulo}

\begin{tablasintitulo}{I{0.200}I{0.800}}{\cab{Campo} & \cab{SUP-26}}
Tipo & Supuesto del alcance \\
Tema & Cobertura inalámbrica del rajo \\
N.º de decisión & no aplica \\
Decisión o supuesto adoptado & La cobertura inalámbrica actual del rajo, más los gabinetes de borde propuestos, sostiene el registro en terreno durante el horizonte del contrato. \\
Fundamento & La construcción de la red de comunicaciones del rajo está excluida del alcance (Cap. 11), pero el registro en terreno depende de ella. El proponente entrega el estudio de cobertura con proyección a tres años y la especificación de lo que faltaría (RT-03.24 del caso). \\
Impacto si resulta equivocado & Zonas de sombra crecientes con el avance del rajo degradan el criterio n.º 13 y obligan a inversión adicional en red, de cargo del Cliente, no prevista en el alcance. \\
Instancia de validación & Estudio de cobertura del rajo con proyección a tres años, incorporado a la EDT, y medición de eventos encolados por sector durante la marcha blanca. Se realiza al comienzo de la Etapa 1. \\
Módulo responsable & Integrador de Sistemas \\
¿Se trata también como riesgo? & Sí \\
\end{tablasintitulo}

\begin{tablasintitulo}{I{0.200}I{0.800}}{\cab{Campo} & \cab{SUP-27}}
Tipo & Supuesto del alcance \\
Tema & Volumen de contactos del centro de atención \\
N.º de decisión & no aplica \\
Decisión o supuesto adoptado & El centro de atención se dimensiona para 1.250 contactos mensuales, por encima del mínimo de 900, en español y 24x7x365, con respuesta presencial en faena para incidentes críticos. \\
Fundamento & El caso no entrega volumen histórico de contactos de soporte. La cifra se estima desde la dotación usuaria proyectada y el régimen de turnos, y condiciona la dotación del servicio de operación. \\
Impacto si resulta equivocado & Si el volumen real es mayor, la dotación comprometida no sostiene los niveles de servicio y el Adjudicatario queda expuesto a multas por incumplimiento de acuerdo de nivel de servicio. \\
Instancia de validación & Medición del volumen real de contactos durante los tres primeros meses de operación, con revisión de dotación acordada con el Cliente. Se realiza en los primeros meses de la operación. \\
Módulo responsable & Administración y Servicios Comunes \\
¿Se trata también como riesgo? & No \\
\end{tablasintitulo}

\begin{tablasintitulo}{I{0.200}I{0.800}}{\cab{Campo} & \cab{SUP-28}}
Tipo & Supuesto del alcance \\
Tema & Reparto del consumo entre áreas \\
N.º de decisión & no aplica \\
Decisión o supuesto adoptado & El consumo de agua, energía y combustible se reparte entre áreas —planta, mina, campamento y supresión de polvo— aplicando una regla de reparto declarada, ante la ausencia de medición independiente por área. \\
Fundamento & El caso no declara que exista medición segregada por área. La regla de reparto es una definición del proponente y queda visible junto al resultado, conforme al requerimiento de exhibir el supuesto aplicable junto al dato calculado. \\
Impacto si resulta equivocado & Si la autoridad o el verificador de la huella de carbono exigen medición independiente por área, el dato calculado no es aceptable y hay que instrumentar medición adicional, de cargo del Cliente. \\
Instancia de validación & Validación con el Cliente, en la presentación preparatoria 3, de la instrumentación existente; verificación por tercero de la huella de carbono, que se pronuncia sobre la regla de reparto declarada. Se realiza en la presentación preparatoria 3 y en el mes 19. \\
Módulo responsable & Cumplimiento Ambiental \\
¿Se trata también como riesgo? & No \\
\end{tablasintitulo}

\begin{tablasintitulo}{I{0.200}I{0.800}}{\cab{Campo} & \cab{SUP-29}}
Tipo & Supuesto del alcance \\
Tema & Etiquetas de proceso que se ingestan \\
N.º de decisión & no aplica \\
Decisión o supuesto adoptado & De las \textasciitilde{}14.500 etiquetas de proceso que el historiador de planta declara hoy —16.000 en la proyección a tres años— se ingestan de forma continua \textasciitilde{}5.000: las que alimentan la conciliación metalúrgica, la balanza de alimentación, el monitoreo ambiental y el cálculo de la huella de carbono. El resto permanece en el historiador, que no se reemplaza, y se lee por el conector del historiador de proceso bajo demanda; se asume que el resto no es necesario para la conciliación, el monitoreo ambiental ni la huella de carbono. \\
Fundamento & El numeral 14.1 declara el universo de etiquetas, no cuáles debe capturar la solución. La ingesta continua de las 14.500 llevaría la capa de series de tiempo de 2 a \textasciitilde{}5,7 TB al año sin que ningún criterio del Capítulo 18 lo exija, y el historiador sigue siendo la fuente para la señal que no se ingesta. La selección es, por tanto, definición del Proponente. \\
Impacto si resulta equivocado & Si el Cliente exige la ingesta completa, la capa de series de tiempo y el ancho de banda del enlace de planta se redimensionan al triple, con impacto en la Oferta Económica y en el Formulario T-11; no cambian los componentes ni los contratos de interfaz. \\
Instancia de validación & Levantamiento de las señales del historiador con la Superintendencia de Planta y Metalurgia, dentro de la sección 4.2 del Subdocumento 4, y validación con el Cliente en la presentación preparatoria 3. Se realiza en la presentación preparatoria 3 y al comienzo de la Etapa 1. \\
Módulo responsable & Conciliación Metalúrgica y Cumplimiento Ambiental \\
¿Se trata también como riesgo? & No \\
\end{tablasintitulo}

\begin{tablasintitulo}{I{0.200}I{0.800}}{\cab{Campo} & \cab{SUP-30}}
Tipo & Supuesto del alcance \\
Tema & Históricos anteriores a la ventana de migración \\
N.º de decisión & no aplica \\
Decisión o supuesto adoptado & Los registros históricos anteriores a la ventana de migración existen, son exportables y quedan accesibles en un repositorio de consulta de solo lectura que el Proponente habilita en la Etapa 1 y conserva durante el plazo de retención del Capítulo 15. \\
Fundamento & El RT-05.15 del documento transversal exige que los datos que no se migran queden accesibles durante el período de retención que fije el caso. El Capítulo 15 declara qué se migra y, por separado, plazos de retención mayores —diez años para trazabilidad y embarques contra cinco migrados—, pero no declara si esos registros anteriores existen, en qué soporte residen ni en qué estado de conservación se encuentran. A falta de esa precisión, el Proponente asume su existencia y compromete el repositorio. \\
Impacto si resulta equivocado & Si el volumen o el formato de esos históricos difieren de lo supuesto —soporte en papel, sistemas discontinuados o exportación no estructurada—, el dimensionamiento y el costo del repositorio cambian, y la carga puede exceder la ventana de la Etapa 1. Si el Cliente, en cambio, mantiene esos registros en sus propios sistemas, el repositorio sale del alcance. \\
Instancia de validación & Consulta N.º 14, formulada al Cliente el 30 de agosto de 2026. El dimensionamiento definitivo se ajusta con la respuesta, antes de la firma. Se realiza durante el período de consultas. \\
Módulo responsable & Administración y Servicios Comunes \\
¿Se trata también como riesgo? & Sí \\
\end{tablasintitulo}

\begin{tablasintitulo}{I{0.200}I{0.800}}{\cab{Campo} & \cab{SUP-31}}
Tipo & Supuesto del alcance \\
Tema & Volumen y calidad de los datos a migrar \\
N.º de decisión & no aplica \\
Decisión o supuesto adoptado & Los datos históricos a migrar no superan en más de un 20 \% la volumetría actual multiplicada por su ventana de migración: 1,71 millones de viajes y 6,85 millones de ciclos de carguío (5 años), 355.000 muestras de laboratorio (5 años), 102.000 órdenes de trabajo de mantenimiento (3 años) y el monitoreo de 186 instrumentos ambientales y geotécnicos (10 años). La tasa de registros con campos obligatorios vacíos o inconsistentes no supera el 5 \% por dominio. Sobre cualquiera de esos umbrales, el saneamiento y la carga adicionales se tratan como cambio conforme al Artículo 72.º de las Bases administrativas. \\
Fundamento & El Caso fija las ventanas de migración (Cap. 15, RT-05.15) y la volumetría anual (Cap. 14.1), pero no el volumen ni la calidad de los históricos, y reconoce no tener inventario técnico de sus sistemas (Cap. 5). Las cifras resultan de multiplicar la volumetría anual actual por cada ventana. La calidad no se conoce: parte de los datos vive en planillas mantenidas a mano desde 2013 (Cap. 4.5). \\
Impacto si resulta equivocado & Si el volumen o la tasa de defectos es mayor, el perfilado y los dos ensayos de migración (TRA-MIG-002 y TRA-MIG-003) se alargan y la migración definitiva puede no quedar lista para la marcha blanca de la Etapa 1 (mes 13). \\
Instancia de validación & Informe de perfilado y saneamiento de los datos de origen (TRA-MIG-002), con recuento por dominio y tasa de defectos. Se realiza entre los meses 3 y 8. \\
Módulo responsable & Administración y Servicios Comunes \\
¿Se trata también como riesgo? & Sí \\
\end{tablasintitulo}

\begin{tablasintitulo}{I{0.200}I{0.800}}{\cab{Campo} & \cab{SUP-32}}
Tipo & Supuesto del alcance \\
Tema & Registros de seguridad en papel \\
N.º de decisión & no aplica \\
Decisión o supuesto adoptado & Los registros de seguridad y salud ocupacional en papel de los últimos 10 años no superan 200.000 formularios y se digitalizan como imagen, indexados por fecha, tipo de registro, área, empresa y persona involucrada, quedando consultables y exportables para fiscalización. No se incluye la transcripción campo a campo. Sobre 200.000 formularios, o si el Cliente exige transcripción estructurada, el trabajo adicional se trata como cambio conforme al Artículo 72.º de las Bases administrativas. \\
Fundamento & El Caso exige migrar 10 años de seguridad y salud ocupacional (Cap. 15) y declara que esos registros nacen en papel y se digitan días después (Cap. 4.8 y 7.3). La cifra resulta de la volumetría actual (Cap. 14.1): 11.500 incidentes, cuasi accidentes y observaciones más 7.800 inspecciones al año, por 10 años, unos 193.000 formularios. El criterio n.º 9 exige que estén disponibles para fiscalización sin buscar carpetas, lo que la imagen indexada cumple. \\
Impacto si resulta equivocado & Si el archivo en papel es mayor o la autoridad exige el dato estructurado, el esfuerzo de digitación se multiplica y compite con la migración de la Etapa 1. \\
Instancia de validación & Inventario físico de los archivos en papel con el área de seguridad y salud ocupacional. Se realiza en el mes 2. \\
Módulo responsable & Seguridad y Salud Ocupacional \\
¿Se trata también como riesgo? & Sí \\
\end{tablasintitulo}

\begin{tablasintitulo}{I{0.200}I{0.800}}{\cab{Campo} & \cab{SUP-33}}
Tipo & Supuesto del alcance \\
Tema & Entrega del hardware de terreno \\
N.º de decisión & no aplica \\
Decisión o supuesto adoptado & El Proponente entrega la especificación aprobada del hardware de terreno (Formulario T-11) en el mes 6 para la Etapa 1 y en el mes 12 para la Etapa 2, y el Cliente entrega ese hardware en faena a más tardar al término del mes 7. Cada semana de atraso en la entrega se registra como causa no imputable al Adjudicatario para efectos del Artículo 80.º y se trata en el Comité Ejecutivo conforme al Artículo 72.º de las Bases administrativas. \\
Fundamento & El hardware de terreno lo adquiere el Cliente y el Proponente especifica qué comprar (Cap. 11; RT-08.10). La instalación, la configuración y las pruebas de aceptación de la certificación de la Etapa 1 (hito H5, mes 12) requieren el hardware definitivo en faena con al menos tres meses de holgura. \\
Impacto si resulta equivocado & Sin el hardware en faena, las pruebas de aceptación y la marcha blanca de la Etapa 1 se hacen sin los dispositivos definitivos, y los umbrales de terreno (RNF-DES-001, RNF-USA-001) no se pueden verificar con el equipo real. \\
Instancia de validación & Aprobación de la especificación en el mes 6 y seguimiento de la orden de compra en el Comité de Proyecto. Se realiza entre los meses 4 y 9. \\
Módulo responsable & Integrador de Sistemas \\
¿Se trata también como riesgo? & Sí \\
\end{tablasintitulo}

\begin{tablasintitulo}{I{0.200}I{0.800}}{\cab{Campo} & \cab{SUP-34}}
Tipo & Supuesto del alcance \\
Tema & Entrega de la sala técnica \\
N.º de decisión & no aplica \\
Decisión o supuesto adoptado & El Proponente entrega la especificación técnica de la sala técnica principal de faena en el mes 1, y el Cliente entrega la obra civil de base terminada (losa, acometida eléctrica y canalizaciones) conforme a esa especificación al término del mes 4, y el Proponente instala sobre ella el módulo prefabricado (T11-29). Cada semana de atraso se registra como causa no imputable al Adjudicatario para efectos del Artículo 80.º y se trata en el Comité Ejecutivo conforme al Artículo 72.º de las Bases administrativas. \\
Fundamento & La obra civil es de cargo del Cliente y su especificación del Proponente (RT-06.06). La sala actual de 40 m² no cumple el Capítulo 6 del documento transversal y debe habilitarse o reemplazarse (Cap. 15, RT-06.01). La infraestructura híbrida debe estar habilitada en el hito H3 (mes 6), lo que exige la sala disponible dos meses antes para instalar y probar. \\
Impacto si resulta equivocado & Sin la sala técnica, los ambientes de faena del hito H3 no se habilitan y se arrastra todo el desarrollo que depende del componente on-premise. \\
Instancia de validación & Acta de recepción de la sala en el Comité de Proyecto. Se realiza en el mes 4. \\
Módulo responsable & Integrador de Sistemas \\
¿Se trata también como riesgo? & Sí \\
\end{tablasintitulo}

\begin{tablasintitulo}{I{0.200}I{0.800}}{\cab{Campo} & \cab{SUP-35}}
Tipo & Supuesto del alcance \\
Tema & Acreditación del personal del Adjudicatario \\
N.º de decisión & no aplica \\
Decisión o supuesto adoptado & La acreditación de cada persona del Adjudicatario que deba ingresar a faena demora en promedio 9 días desde la entrega de su documentación completa, con examen de altura geográfica vigente. La planificación inicia la acreditación 15 días hábiles antes de cada actividad en faena. Si la acreditación se extiende por causa del Cliente, el atraso se registra como causa no imputable al Adjudicatario para efectos del Artículo 80.º de las Bases administrativas. \\
Fundamento & El personal del Adjudicatario se acredita como cualquier contratista (restricción n.º 12) y debe contar con examen de altura geográfica vigente (Cap. 15, RT-21.16). Nueve días es el promedio actual de acreditación de una persona contratista nueva (Cap. 4.8 y 7.3). El Caso exige que el plan refleje estos tiempos (Cap. 17.5). \\
Impacto si resulta equivocado & Si la acreditación demora más, se atrasan la instalación en faena, las pruebas y el acompañamiento de las marchas blancas, que el Caso exige presencial en día y en noche (Cap. 13.3). \\
Instancia de validación & Acreditación del primer grupo del Adjudicatario al inicio del Contrato. Se realiza en el mes 1. \\
Módulo responsable & Acreditación y Acceso \\
¿Se trata también como riesgo? & Sí \\
\end{tablasintitulo}

\seccion{Exclusiones}
Las exclusiones son lo que la compañía declara que no está pidiendo, según el capítulo 11 de las Bases técnicas del caso Minería, página 20. Que algo esté excluido no significa que pueda ignorarse, ya que genera dependencias que se gestionan como riesgos del proyecto. La columna “Dependencia que genera” es de elaboración propia.

\begin{tabla}{I{0.110}I{0.329}I{0.164}I{0.397}}{Exclusiones del caso y dependencias que generan.}{tab:exclusiones-caso-dependencias-generan}
  {\cab{ID} & \cab{Exclusión} & \cab{Supuestos relacionados} & \cab{Dependencia que genera}}
  {Fuente: elaboración propia a partir de las Bases técnicas del caso Minería (CMAA, 2026b, cap. 11, p. 20).}
EXC-01 & No se pide reemplazar el ERP corporativo ni ninguno de sus módulos. & SUP-16 & Integración con SAP S/4HANA por interfaz; depende de que existan interfaces documentadas (SUP-16). Restricción 4. \\
EXC-02 & No se pide reemplazar el sistema de despacho de flota, el de gestión de laboratorio ni el software de planificación minera. & SUP-14, SUP-16 & Los tres se integran; la calidad de los datos y la frecuencia de lectura dependen de sus interfaces (SUP-16) y del versionado del plan (SUP-14). Restricciones 5, 6 y 7. \\
EXC-03 & No se pide intervenir el control de proceso de la planta ni modificar la lógica de control. & SUP-19 & Solo lectura vía réplica del historiador (SUP-19); si no existe la réplica, se afecta la lectura de tonelaje y detenciones de planta. Restricción 3. \\
EXC-04 & No se pide automatizar equipos mineros (operación autónoma de camiones o perforadoras). & SUP-10 & La solución solo lee telemetría; su alcance en Equipos y Mantenimiento depende de que el Cliente pueda extraerla (SUP-10). \\
EXC-05 & No se pide el diseño del plan minero ni la estimación de recursos y reservas. & SUP-14 & La solución consume las salidas de planificación; el cumplimiento se mide contra la versión vigente del plan (SUP-14). Restricción 7. \\
EXC-06 & No se pide la gestión de remuneraciones ni la administración de personal (residen en el ERP). & No aplica & Sin dependencia identificada en el registro; la solución no toma datos de remuneraciones. \\
EXC-07 & No se pide la operación del terminal portuario (lo administra un tercero), sin perjuicio de la integración necesaria. & SUP-24 & El registro en puerto se diseña para no depender del operador (gabinete de borde, operación desconectada); condiciones del contrato con el terminal por validar (SUP-24). \\
EXC-08 & No se pide construir la red de comunicaciones del rajo, aunque sí evaluar si la cobertura actual es suficiente y especificar qué faltaría. & SUP-26, SUP-09 & La cobertura actual más los gabinetes de borde debe sostener el registro (SUP-26); las zonas de sombra se cubren con almacenamiento local de 48 horas (SUP-09). Dependencia tratada en el plan y en el registro de riesgos. \\
EXC-09 & El hardware de terreno lo adquiere el Cliente; el Proponente especifica qué comprar, cuánto y con qué características (CMAA, 2026c, cap. 8). & SUP-33, SUP-34, SUP-18 & Sin el hardware en faena no hay pruebas de aceptación ni marcha blanca de la Etapa 1 (SUP-33); la sala técnica (SUP-34) y la alimentación en cabinas (SUP-18) son de cargo del Cliente. \\
\end{tabla}

\seccion{Restricciones no negociables}
Las restricciones son condiciones que la compañía no pone en discusión, según el capítulo 10 de las Bases técnicas del caso Minería, páginas 19 y 20. Cada restricción se relaciona con los requerimientos del catálogo que la operacionalizan y con los supuestos que la mencionan.

\begin{tabla}{I{0.095}I{0.326}C{0.189}I{0.224}I{0.166}}{Restricciones no negociables y requerimientos que las operacionalizan.}{tab:restricciones-no-negociables-requerimientos-operacionalizan}
  {\cab{ID} & \cab{Restricción} & \cab{N.º de requerimientos} & \cab{Requerimientos que la operacionalizan} & \cab{Supuestos relacionados}}
  {Fuente: elaboración propia a partir de las Bases técnicas del caso Minería (CMAA, 2026b, cap. 10, pp. 19–20).}
RES-01 & La mina y la planta no se detienen (24/7, 365 días). Únicas ventanas mayores: dos detenciones de planta de seis días, en marzo y septiembre. & 1 & RNF-DIS-004 & No aplica \\
RES-02 & La operación no puede depender del enlace con Antofagasta; debe seguir funcionando durante un corte de fibra prolongado. & 7 & RF-CNX-002, RF-CNX-003, RF-CNX-005, RF-EQU-006, RNF-DIS-001, RNF-DIS-002, RP-003 & SUP-09, SUP-24 \\
RES-03 & La red de control de proceso de la planta es de solo lectura: se lee del historiador, no se escribe ni se instalan componentes sin la segregación aprobada. & 4 & RF-PLA-006, RF-PLA-009, RNF-INT-001, RNF-SEG-004 & SUP-16, SUP-19 \\
RES-04 & SAP no se reemplaza ni se modifica; la solución le entrega datos y SAP emite. No hay dos verdades contables. & 8 & RF-CON-002, RF-CON-013, RF-CON-015, RF-EQU-004, RF-EQU-005, RF-TRZ-019, RNF-INT-004, TRA-MIG-005 & SUP-16 \\
RES-05 & El sistema de despacho de flota se mantiene (contrato hasta 2031); se integra, no se sustituye. & 3 & RF-TER-004, RF-TRZ-001, RNF-INT-003 & SUP-16 \\
RES-06 & El sistema de gestión de laboratorio se mantiene y se integra en ambos sentidos. & 6 & RF-CON-009, RF-CON-012, RF-CON-016, RF-CON-017, RNF-INT-002, TRA-MIG-005 & SUP-16 \\
RES-07 & El software de planificación minera se mantiene; la solución consume sus salidas. & 1 & RF-PLA-001 & SUP-16 \\
RES-08 & No se agregan pasos al ciclo productivo; todo registro adicional al operador se justifica en segundos y se demuestra en la marcha blanca. & 14 & RF-ACR-011, RF-CON-011, RF-EQU-005, RF-PLA-010, RF-TER-001, RF-TER-002, RNF-DES-001, RNF-DES-004, RNF-DES-016, RNF-USA-004, RNF-USA-008, RNF-USA-009, TRA-001, TRA-004 & SUP-04, SUP-12 \\
RES-09 & La identidad se federa con el directorio corporativo del grupo; doble factor obligatorio, sin credenciales compartidas y con registro completo de accesos. & 6 & RF-ACR-013, RF-ACR-016, RF-ADM-005, RF-ADM-007, RF-CNX-003, RNF-SEG-003 & SUP-08 \\
RES-10 & El equipo de tecnologías de información del Cliente son once personas; toda función que exija un especialista dedicado se ofrece como servicio y costeada. & 4 & RF-ADM-002, RNF-OPE-001, RP-005, RP-006 & No aplica \\
RES-11 & El personal clave de operaciones tiene disponibilidad limitada en las detenciones de marzo y septiembre; las validaciones que lo requieran se planifican fuera de esas ventanas. & 1 & RP-011 & No aplica \\
RES-12 & Toda intervención en faena está sujeta al Reglamento de Seguridad Minera y a los procedimientos del Cliente; el personal del Adjudicatario se acredita como cualquier contratista. & 2 & RP-001, RP-008 & SUP-35 \\
\end{tabla}

\anexo{Otros listados}
Este anexo reúne dos listados de trazabilidad que complementan a los anteriores: los criterios de aceptación del caso frente a los requerimientos, y los supuestos frente a los requerimientos que dependen de ellos. A partir de C.3 reúne además el detalle de las secciones 2.2 a 2.4 del Subdocumento 2 que el cuerpo sintetiza.

\seccion{Trazabilidad de los criterios de aceptación a los requerimientos}
Los 14 criterios son los resultados de negocio con que la compañía juzgará el proyecto, según el capítulo 18 de las Bases técnicas del caso Minería, páginas 32 y 33. Cada criterio se relaciona con los requerimientos que lo citan en su origen. La columna “Resultado esperado” es una síntesis del texto del caso.

\begin{tabla}{C{0.105}I{0.409}C{0.190}I{0.295}}{Criterios de aceptación y requerimientos que los sustentan.}{tab:criterios-aceptacion-requerimientos-sustentan}
  {\cab{N.º de criterio} & \cab{Resultado esperado} & \cab{N.º de requerimientos} & \cab{Requerimientos que lo sustentan}}
  {Fuente: elaboración propia a partir de las Bases técnicas del caso Minería (CMAA, 2026b, cap. 18, pp. 32–33).}
1 & Origen del mineral de cualquier lote embarcado reconstruido con evidencia en menos de dos horas (hoy 4 a 6 semanas). & 10 & RF-TRZ-002, RF-TRZ-003, RF-TRZ-004, RF-TRZ-005, RF-TRZ-006, RF-TRZ-007, RF-TRZ-010, RF-TRZ-015, RNF-DES-005, TRA-MIG-005 \\
2 & Conciliación metalúrgica mensual cierra en no más de tres días hábiles (hoy 18). & 3 & RF-CON-003, RNF-DES-006, TRA-002 \\
3 & Diferencia de conciliación converge y se sostiene bajo el umbral comprometido; la diferencia persistente queda explicada por punto de la cadena (hoy 3,1 \% a 6,8 \% sin explicación). & 8 & RF-CON-002, RF-CON-004, RF-CON-005, RF-CON-006, RF-CON-007, RF-CON-010, RNF-DES-014, TRA-002 \\
4 & Cumplimiento del plan minero conocido al inicio del turno siguiente, contra la versión vigente del plan (hoy 10 a 15 días de desfase). & 3 & RF-PLA-002, RF-PLA-003, RNF-DES-008 \\
5 & Totalidad de detenciones de equipo y de planta clasificadas con causa dentro del turno siguiente (hoy 23 \% sin clasificar; 61 \% sin causa raíz). & 6 & RF-PLA-005, RF-PLA-006, RF-PLA-007, RF-PLA-008, RNF-DES-009, RNF-USA-004 \\
6 & El supervisor dispone de estado, historial, alarmas de telemetría y pautas pendientes de un equipo en una sola vista (hoy cuatro sistemas y una radio). & 5 & RF-EQU-001, RF-EQU-002, RF-EQU-004, RNF-DES-011, RNF-DIS-005 \\
7 & Ninguna persona rechazada en portería por documentación vencida que el sistema pudo advertir (hoy 14 \% de los rechazos). & 1 & RF-ACR-006 \\
8 & Persona contratista nueva acreditada en menos de 48 horas (hoy 9 días). & 5 & RF-ACR-001, RF-ACR-002, RF-ACR-003, RNF-DES-015, RNF-SEG-003 \\
9 & Registros de seguridad, inspecciones e incidentes nacen digitales en el lugar del hecho y disponibles para fiscalización. & 7 & RF-ACR-014, RF-CUM-001, RF-CUM-002, RF-CUM-003, RF-CUM-004, RNF-USA-007, TRA-MIG-005 \\
10 & Toda alerta de fatiga tiene acción documentada y trazable. & 2 & RF-CUM-005, RF-CUM-006 \\
11 & Reportes a la autoridad ambiental generados desde el sistema, con trazabilidad hasta la medición de origen. & 3 & RF-CUM-010, RF-CUM-012, TRA-MIG-005 \\
12 & Huella de carbono por tonelada de cobre fino calculada mensualmente y auditable por un tercero (hoy anual, con consultora). & 6 & RF-CUM-011, RF-CUM-015, RF-CUM-016, RF-CUM-017, RNF-DES-013, RNF-POR-002 \\
13 & La operación continúa y registra durante 48 horas sin enlace y sincroniza sin intervención manual. & 7 & RF-CNX-001, RF-CNX-003, RF-CNX-005, RF-EQU-006, RF-TER-005, RNF-DES-012, RNF-DIS-003 \\
14 & El operador de pala usa la solución en su turno, con guantes, y por decisión propia (hoy dos iniciativas abandonadas). & 6 & RF-TER-001, RF-TER-002, RNF-USA-001, RNF-USA-005, TRA-005, TRA-010 \\
\end{tabla}

\seccion{Requerimientos que dependen de cada supuesto}
Si un supuesto resulta equivocado, los requerimientos de este listado son los que habría que revisar. Los supuestos SUP-07 y SUP-15 no tienen requerimientos que los citen, porque afectan al alcance y al cronograma y no a un requerimiento específico.

\begin{tabla}{I{0.110}I{0.249}I{0.280}I{0.361}}{Requerimientos que dependen de cada supuesto.}{tab:requerimientos-dependen-cada-supuesto}
  {\cab{ID} & \cab{Tema} & \cab{Requerimientos que dependen del supuesto} & \cab{Origen en las bases}}
  {Fuente: elaboración propia.}
SUP-01 & Unidad mínima de trazabilidad & RF-TRZ-001, RF-TRZ-002, RF-TRZ-003, RF-TRZ-005 & Bases técnicas del caso Minería, cap. 16.1, decisión 1 \\
SUP-02 & Fuente de verdad del tonelaje & RF-CON-001, RF-CON-002, RF-CON-006, TRA-002 & Bases técnicas del caso Minería, cap. 16.1, decisión 2 \\
SUP-03 & Mezcla en stocks & RF-TRZ-003, RF-TRZ-006, RF-TRZ-007, RF-TRZ-008, RF-TRZ-009, RNF-DES-005 & Bases técnicas del caso Minería, cap. 16.1, decisión 3 \\
SUP-04 & Identificación del operador & RF-ACR-011, RF-ACR-014, RNF-DES-004 & Bases técnicas del caso Minería, cap. 16.1, decisión 4 \\
SUP-05 & Definición de detención & RF-PLA-005, RF-PLA-006, RF-PLA-007, RF-PLA-008, RNF-DES-009, RNF-USA-004, RNF-USA-009 & Bases técnicas del caso Minería, cap. 16.1, decisión 5 \\
SUP-06 & Muestras rechazadas & RF-TRZ-018 & Bases técnicas del caso Minería, cap. 16.1, decisión 6 \\
SUP-07 & Concentrado devuelto & No aplica & Bases técnicas del caso Minería, cap. 16.1, decisión 7 \\
SUP-08 & Autenticación de contratistas & RF-ACR-005, RF-ACR-012, RF-ACR-013, RNF-SEG-003 & Bases técnicas del caso Minería, cap. 16.1, decisión 8 \\
SUP-09 & Zona de sombra con cambio de turno & RF-CNX-003, RF-EQU-006, RF-TER-005, RNF-DES-012, RNF-DIS-001, RNF-DIS-003 & Bases técnicas del caso Minería, cap. 16.1, decisión 9 \\
SUP-10 & Dueño de la telemetría & RF-EQU-003, RNF-DIS-005 & Bases técnicas del caso Minería, cap. 16.1, decisión 10 \\
SUP-11 & Carga en borde de polígono & RF-TER-001, RF-TRZ-003, RF-TRZ-010, RF-TRZ-011, RNF-DES-005 & Bases técnicas del caso Minería, cap. 16.1, decisión 11 \\
SUP-12 & Muestreo versus ciclo de mina & RF-CON-011, RF-CON-017, TRA-001 & Bases técnicas del caso Minería, cap. 16.1, decisión 12 \\
SUP-13 & Biometría y sindicato & RF-ACR-008, RF-ACR-009, RF-ACR-010 & Bases técnicas del caso Minería, cap. 16.1, decisión 13 \\
SUP-14 & Versión del plan & RF-PLA-002, RF-PLA-003, RNF-DES-008 & Bases técnicas del caso Minería, cap. 16.1, decisión 14 \\
SUP-15 & Fecha de inicio del contrato & No aplica & Bases administrativas, art. 17 · Bases técnicas del caso Minería, hitos externos (enero de 2029) \\
SUP-16 & Lectura de sistemas existentes & RF-CON-009, RF-CON-015, RF-CON-017, RF-PLA-001, RF-TRZ-001, RNF-INT-001, RNF-INT-002, RNF-INT-003, RNF-INT-004 & Bases técnicas del caso Minería, cap. 5 y cap. 10 \\
SUP-17 & Credencial actual & RF-ACR-011 & Bases técnicas del caso Minería, cap. 4.8 (control de acceso con lectura de credencial) \\
SUP-18 & Alimentación en cabinas de pala & RNF-USA-003 & Bases técnicas del caso Minería, cap. 6 (equipamiento de terreno) \\
SUP-19 & Réplica del historiador de proceso & RNF-INT-001 & Bases técnicas del caso Minería, restricción 3 (red de control de solo lectura) · Bases técnicas transversales, RT-05.23 \\
SUP-20 & Fuente de verdad de la ley medida & RF-CON-004, RF-CON-009, RF-CON-012 & Vacío adicional identificado por el proponente (CMAA, 2026b, cap. 16.1, nota final) \\
SUP-21 & Material anterior al sistema & RF-TRZ-003, RF-TRZ-005, RF-TRZ-007, RNF-DES-005 & Vacío adicional identificado por el proponente (CMAA, 2026b, cap. 16.1, nota final) \\
SUP-22 & Certificado para el comprador & RF-TRZ-016 & Vacío adicional identificado por el proponente (CMAA, 2026b, cap. 16.1, nota final) \\
SUP-23 & Vector de elementos del lote & RF-TRZ-006, RF-TRZ-012 & Vacío adicional identificado por el proponente (CMAA, 2026b, cap. 16.1, nota final) \\
SUP-24 & Operador portuario & RF-TRZ-013, RF-TRZ-014, RNF-DIS-001, RNF-INF-001 & Vacío adicional identificado por el proponente (CMAA, 2026b, cap. 16.1, nota final) \\
SUP-25 & Fuente horaria común & RF-CNX-003, RF-TRZ-017, RNF-CUM-002 & Vacío adicional identificado por el proponente (CMAA, 2026b, cap. 16.1, nota final) \\
SUP-26 & Cobertura inalámbrica del rajo & RP-004 & Bases técnicas del caso Minería, cap. 6 y cap. 11 \\
SUP-27 & Volumen de contactos del centro de atención & RP-007 & Bases técnicas del caso Minería, cap. 14 (contactos mensuales “a estimar y declarar como supuesto”) · Bases técnicas transversales, RT-21.14 \\
SUP-28 & Reparto del consumo entre áreas & RF-CUM-010, RF-CUM-011, RF-CUM-015 & Bases técnicas del caso Minería, entrevista de Guerrero (un solo medidor y una regla de reparto de 2019) \\
SUP-29 & Etiquetas de proceso que se ingestan & RF-CON-001, RF-CUM-008, RNF-CAP-004 & Bases técnicas del caso Minería, cap. 14.1 (≈14.500 etiquetas, ≈16.000 a tres años) \\
SUP-30 & Históricos anteriores a la ventana de migración & RF-ADM-001, TRA-MIG-001, TRA-MIG-005 & Bases técnicas transversales, RT-05.15 \\
SUP-31 & Volumen y calidad de los datos a migrar & TRA-MIG-001, TRA-MIG-002, TRA-MIG-003, TRA-MIG-004, TRA-MIG-005 & Bases técnicas del caso Minería, cap. 14.1 y cap. 15 · Bases técnicas transversales, RT-05.11 a RT-05.15 \\
SUP-32 & Registros de seguridad en papel & TRA-MIG-005 & Bases técnicas del caso Minería, cap. 4.8, cap. 7.3, cap. 14.1 y cap. 15 \\
SUP-33 & Entrega del hardware de terreno & RF-CNX-004, RNF-USA-003 & Bases técnicas del caso Minería, cap. 11 · Bases técnicas transversales, RT-08.10 \\
SUP-34 & Entrega de la sala técnica & RNF-INF-001 & Bases técnicas del caso Minería, cap. 6 y cap. 15 · Bases técnicas transversales, RT-06.06 \\
SUP-35 & Acreditación del personal del Adjudicatario & RP-008 & Bases técnicas del caso Minería, cap. 4.8, cap. 10 restricción 12 y cap. 15, RT-21.16 \\
\end{tabla}

\seccion{Lugares de operación de la compañía}
La operación de la compañía se reparte entre la faena, Antofagasta y Santiago, en lugares con condiciones muy distintas entre sí, lo que influye en cómo se genera y cómo viaja la información. Los lugares donde opera son los siguientes:

\begin{itemize}
\item Rajo Aranda: es donde se perfora, se trona y se carga y transporta el mineral y el estéril hacia los botaderos y los stocks intermedios. Se trabaja a la intemperie, entre −6 °C y 34 °C, con viento, polvo en suspensión y radiación solar extrema, y la cobertura de red es desigual.
\item Planta concentradora: es donde el mineral pasa por chancado, molienda, flotación, espesamiento y filtrado. Opera de forma continua y cuenta con una sala de control de proceso en una red segregada.
\item Depósito de relaves: se ubica a 6 km de la planta, recibe los relaves espesados y concentra el monitoreo geotécnico y de aguas que se reporta a la autoridad, con una instrumentación que transmite de forma intermitente.
\item Área de servicios y portería: es el acceso a la faena, donde se acredita y controla a los contratistas y donde están la policlínica, las bodegas y los talleres de mantenimiento. Concentra el flujo de personas y es el punto donde hoy se producen los rechazos de ingreso.
\item Acopio y embarque en Antofagasta: recibe el concentrado desde la faena, lo acopia hasta 25.000 toneladas y lo embarca en un terminal portuario administrado por un tercero, junto a la nave, el agente de aduana y el organismo fiscalizador.
\item Centro Integrado de Operaciones: es la sala de control integrada de mina y planta, inaugurada en 2022 en Antofagasta, a 145 km de la faena. Depende por completo del enlace de comunicaciones y hoy se usa cerca del 40 \% de su capacidad.
\item Oficina corporativa: se ubica en Las Condes, Santiago, y reúne a las áreas de finanzas, comercial, sustentabilidad, abastecimiento y tecnologías de información.
\end{itemize}

\seccion{Sistemas de la compañía}
La compañía ya cuenta con sistemas en funcionamiento y definió qué ocurre con cada uno de ellos: cuáles se mantienen, cuáles deben integrarse y cuáles pueden reemplazarse. Los sistemas con los que cuenta son los siguientes:

\begin{itemize}
\item ERP corporativo (SAP S/4HANA): cubre finanzas, abastecimiento, inventarios, recursos humanos y planificación de mantenimiento.
\item Sistema de despacho de flota: asigna los camiones y registra los tiempos de ciclo, el pesaje a bordo y el posicionamiento.
\item Sistema de control distribuido e historiador de proceso: controla y registra el proceso de la planta.
\item Sistema de gestión de laboratorio: registra las muestras, los análisis y sus resultados. Se mantiene y debe integrarse en ambos sentidos.
\item Software de planificación minera: contiene el modelo de bloques, el diseño, el secuenciamiento y el plan minero.
\item Sistema de pesaje de báscula de camiones: registra el pesaje de control a la salida del área de descarga y en el despacho a puerto.
\item Portales de telemetría de fabricantes: son tres, uno por marca de equipo, y entregan alarmas, horas, consumos y códigos de falla.
\item Plataforma de control de fatiga y somnolencia: es un sistema de un tercero con cámaras en cabina y alertas.
\item Control de acceso de portería: opera con torniquetes y lectura de credencial.
\item Planillas de cálculo y carpetas compartidas: sostienen la acreditación, el balance metalúrgico, el control de stocks, el reporte ambiental y el cumplimiento del plan.
\end{itemize}

\seccion{Particularidades operacionales y estacionales de la faena}
La operación de Aranda tiene características que condicionan cómo se genera la información. Se identifican las siguientes:

\begin{itemize}
\item Opera de forma continua, las 24 horas de los 365 días del año, en dos turnos de doce horas (08:00 a 20:00 y 20:00 a 08:00) y con régimen 7x7 en faena. La operación funciona sin pausas y la mitad de las personas trabaja de noche.
\item El mineral pasa por muchas manos antes de convertirse en un embarque, y cada etapa tiene su propio responsable y su propia forma de registrar, tal como se muestra en la Figura \refext{S2}{fig:ciclo-mineral-aranda-flujo-actual-informacion}.
\item Las condiciones físicas del rajo son extremas: 3.180 metros de altura, temperaturas entre −6 °C y 34 °C, ráfagas sobre 70 km/h, polvo en suspensión permanente y radiación solar extrema. El personal trabaja con guantes gruesos y lentes de seguridad.
\item La conectividad no es pareja. La red LTE privada instalada en 2021 tiene zonas de sombra permanentes en el fondo del rajo, en el botadero norte y en dos sectores de rampa, y esas zonas cambian a medida que el rajo se profundiza. El enlace de fibra con Antofagasta depende de un único proveedor y registra entre tres y seis cortes al año, de cuatro a veinte horas. Durante un corte, el respaldo satelital solo alcanza para telefonía y correo.
\item La planta tiene una red de control de proceso segregada que la compañía no permite intervenir. Se puede leer de su historiador, pero no se escribe en ella.
\item Una parte importante de la operación se supervisa lejos de la faena, desde el CIO en Antofagasta y desde la oficina de Santiago, y ambos dependen de que la información llegue desde terreno.
\item Cerca del 57 \% de las personas con acceso a la faena son contratistas (2.400 de 4.180), y su número sube hasta 6.700 durante las detenciones mayores de planta.
\end{itemize}
El flujo de la información que resume la Figura \refext{S2}{fig:ciclo-mineral-aranda-flujo-actual-informacion} del cuerpo del Subdocumento 2 se describe en tres tramos.

En el primer tramo, de la planificación al carguío que corresponde a las etapas 1,2 y 3, el plan semanal se exporta cada domingo y se carga a mano al sistema de despacho, y los cambios que ocurren durante la semana se comunican por radio y por el grupo de mensajería del turno. Solo dos de las nueve perforadoras registran de forma automática, y el avance del resto se anota en una planilla al final del turno. Cuando el operador de la pala detecta que el material cambió, lo avisa por radio y el despachador reasigna el destino del camión, pero la razón del cambio no queda registrada.

En el segundo tramo, de los stocks intermedios al procesamiento que corresponde a las etapas 4 y 5, el control de cada stock se lleva en una planilla que actualiza el supervisor de turno. El laboratorio entrega sus resultados entre cuatro y treinta y seis horas después de tomada la muestra, y estos se transcriben a mano a las planillas de metalurgia. El balance del turno y el del mes se calculan en planillas distintas.

En el tercer tramo que corresponde a las etapas 6, 7 y 8, del transporte a la venta, el concentrado viaja a Antofagasta con una guía de despacho y una planilla de pesaje en papel. En el acopio del terminal portuario se forma el lote comercial mezclando material, sin vínculo con su origen, y la liquidación con el comprador se realiza meses después, sobre una muestra que pueden resolver hasta tres laboratorios.

El calendario de Aranda tiene momentos en que la operación cambia de ritmo, y algunos coinciden justamente con los momentos en que el problema se hace más visible. Se identifican los siguientes:

\begin{itemize}
\item Detenciones programadas de planta en marzo y septiembre, de seis días cada una. Son las únicas ventanas de intervención mayor, ingresan hasta 4.300 contratistas adicionales y el personal clave de planta no está disponible para otra cosa. Es también cuando más presión recibe la acreditación de personas.
\item Invierno altiplánico en enero y febrero, con episodios de lluvia y tormenta eléctrica de uno a tres días que pueden cortar la energía o suspender la faena.
\item Cierre contable durante los primeros cinco días hábiles de cada mes. La conciliación metalúrgica, que hoy tarda 18 días hábiles, termina mucho después de que el cierre contable ya ocurrió.
\item Embarques de concentrado, 31 al año (aproximadamente uno cada doce días), que duran entre 36 y 48 horas.
\item Cambios de turno a las 08:00 y a las 20:00, con 45 a 60 minutos de alta carga en registros, accesos y comunicaciones. Es el momento en que se pierden, en promedio, 47 minutos efectivos por turno.
\item Tronaduras casi diarias, con ventanas de 30 a 45 minutos de evacuación del sector y pérdida transitoria de cobertura.
\item Ciclo 7x7: cada grupo pasa siete días en faena y siete fuera, por lo que llegar a todas las personas de un área toma al menos catorce días.
\item Reportes a la autoridad ambiental con frecuencias mensuales, trimestrales y anuales, y al menos una auditoría anual de la autoridad minera, de tres a cinco días.
\end{itemize}

\seccion{Marco normativo y compromisos con terceros}
La minería del cobre es una de las industrias más fiscalizadas del país, y Aranda responde a varias autoridades al mismo tiempo, cada una con sus propias exigencias de registro y de reporte. En la Tabla \ref{tab:marco-normativo-compromisos-terceros-relevantes-problema} se presenta el marco que se identificó como relevante para el problema, junto con la situación actual de la compañía en cada ámbito.

\begin{tabla}{I{0.155}I{0.283}I{0.284}I{0.278}}{Marco normativo y compromisos con terceros relevantes para el problema.}{tab:marco-normativo-compromisos-terceros-relevantes-problema}
  {\cab{Ámbito} & \cab{Norma o compromiso} & \cab{Obligación que impone} & \cab{Situación en Aranda}}
  {Fuente: elaboración propia a partir de las Bases técnicas del caso Minería (CMAA, 2026b, caps. 7, 8 y 12, pp. 12, 15 y 20–21).}
Seguridad minera & Reglamento de Seguridad Minera (Ministerio de Minería, 2004), fiscalizado por Sernageomin & Condiciones de seguridad de la faena, registro y reporte de incidentes e inspecciones & Seis de las once observaciones de la última auditoría fueron por trazabilidad de registros, que hoy nacen en papel \\
Cierre de faenas & Ley N° 20.551 (2011) & Información periódica y trazable del avance de la operación y de las obras comprometidas & La información de avance se consolida a mano \\
Medio ambiente & Resolución de Calificación Ambiental de la faena y Superintendencia del Medio Ambiente (Ley N° 20.417, 2010) & Compromisos de monitoreo y reporte con frecuencia, punto de medición y destinatario definidos & Reportes armados a mano con datos de cinco áreas y una observación en 2025 por inconsistencia \\
Salud ocupacional & Decreto 594 (1999) y protocolos de vigilancia por altura geográfica y exposición a sílice & Exámenes y periodicidades que condicionan la habilitación de cada persona & Vencimientos no advertidos a tiempo, que explican el 14 \% de los rechazos en portería \\
Datos personales & Ley N° 21.719 (2024) & Protección reforzada de datos sensibles, como los de salud & Datos de salud ocupacional, fatiga y control de alcohol y drogas circulan hoy en planillas y correos \\
Redes industriales & Serie IEC 62443 (International Electrotechnical Commission [IEC], s.f.) & Segregación entre tecnologías de información y de operación & Condición para cualquier lectura desde la red de control de la planta \\
Cobre responsable & The Copper Mark (2026) y exigencias de los compradores & Evidencia verificable del origen y de la custodia del material en cada etapa & Origen de esta licitación, ya que hoy el origen de un lote es una estimación \\
Huella de carbono & Protocolo de Gases de Efecto Invernadero, alcances 1, 2 y 3 (WRI y WBCSD, 2004) & Cálculo por alcance con factores de emisión y verificación por un tercero & Hoy se calcula una vez al año con una consultora y toma dos meses \\
Comercialización & Contratos de venta y prácticas de liquidación de concentrados & Ley pagable, deducciones, cargos de tratamiento y refinación y penalidades por deletéreos & Cuatro de 31 embarques penalizados por arsénico en 2025 \\
\end{tabla}

\seccion{Puntos de quiebre de la trazabilidad del mineral}
Las nueve brechas de trazabilidad del mineral se detallan a continuación:

\begin{itemize}
\item No existe un registro de la versión del plan que estaba vigente en cada momento. El plan semanal se exporta cada domingo, pero cambia tres o cuatro veces por semana y esos cambios se comunican por radio y por el grupo de mensajería del turno.
\item Cuando el operador de pala detecta que el material cambió, lo avisa por radio y el despachador reasigna el destino del camión. El cambio queda en el sistema de despacho, pero la razón del cambio no.
\item Cuando la pala carga en el borde entre dos polígonos, todo el material se asigna al polígono donde está la pala. Es una simplificación conocida por la operación, sin una decisión formal sobre mantenerla o corregirla.
\item Existen tres mediciones del tonelaje (el pesaje a bordo del camión, la báscula que pesa solo una muestra de los viajes y la balanza de alimentación de la planta) que no coinciden, y ninguna ha sido declarada como fuente de verdad. Cada área utiliza la medición que emplea en su propio informe.
\item En los stocks intermedios se asume que la ley del material retomado es el promedio ponderado de lo que se acopió durante semanas. Ese supuesto no ha sido verificado.
\item Los resultados del laboratorio llegan entre cuatro y treinta y seis horas después de tomada la muestra y luego se transcriben a mano a las planillas de metalurgia.
\item El balance metalúrgico del turno y el del mes se calculan en planillas distintas, y el del mes aplica factores de ajuste acumulados desde 2013 que solo dos personas saben explicar.
\item El concentrado viaja a puerto con una guía de despacho y una planilla de pesaje en papel, y en el acopio portuario el lote comercial se forma mezclando material sin vínculo con su origen.
\item Reconstruir el origen de un lote embarcado toma entre cuatro y seis semanas y el resultado es una estimación, mientras que la liquidación con el comprador se cierra meses después sobre muestras que dirimen hasta tres laboratorios.
\end{itemize}

\seccion{Debilidades de la productividad operacional}
En el caso de la productividad, pese a que la operación cuenta con un sistema de despacho que el personal usa todos los días, la información de cada equipo vive repartida en lugares distintos y se junta recién cuando alguien la necesita. A raíz de esto, se identifican las siguientes debilidades:

\begin{itemize}
\item Ante un camión detenido, el supervisor debe abrir el sistema de despacho para ver dónde está, el portal del fabricante para ver si hubo alarmas y SAP para ver si tenía mantenimiento pendiente, y además llamar por radio. Son cuatro sistemas y una radio, y entre quince y veinte minutos con un camión de 240 toneladas detenido.
\item Los fabricantes de camiones, palas y perforadoras entregan la telemetría en tres portales web distintos, sin integración con SAP ni con el despacho. No se revisan de forma sistemática y uno de ellos casi no se consulta.
\item No existe claridad contractual sobre quién es dueño de los datos de telemetría que generan los propios equipos de la compañía, ni sobre si pueden extraerse por vía programática.
\item El 23 \% del tiempo registrado corresponde a demoras sin clasificar y el 61 \% de las detenciones no tiene causa raíz documentada. No hay una definición compartida de qué es una detención, cómo se clasifica ni quién la clasifica.
\item Los avisos de falla se levantan por radio y un planificador los transcribe después a SAP, y el tiempo entre una falla y su diagnóstico es de 4,2 horas.
\item Solo dos de las nueve perforadoras registran automáticamente sus pozos. Los datos de perforación, que entregan información sobre la dureza de la roca, no se aprovechan, y el área de geología los pidió tres veces sin poder usarlos.
\item El cumplimiento del plan se calcula con diez a quince días de desfase y contra el plan exportado el domingo, que para ese momento ya no era el vigente. Por eso el indicador puede mostrar un incumplimiento aún cuando la operación siguió el plan vigente.
\item En cada cambio de turno se pierden 47 minutos efectivos, casi el doble de la meta interna de 25 minutos.
\item El CIO de Antofagasta, inaugurado en 2022, se usa cerca del 40 \% de su capacidad, porque varias funciones dependían de información que solo existe en la faena o en el conocimiento de quienes trabajan en ella.
\end{itemize}

\seccion{Brechas de cumplimiento regulatorio y gestión de personas}
Al ser una operación con más de cuatro mil personas en régimen, cerca de noventa empresas contratistas y compromisos ante varias autoridades, el control de las personas y el cumplimiento de las obligaciones de reporte requieren un seguimiento riguroso. Sin embargo, en este ámbito se identifican los siguientes puntos:

\begin{itemize}
\item Acreditar a un trabajador contratista nuevo demora en promedio nueve días, y la validación documental se apoya en una planilla que mantiene una sola persona del área de seguridad y salud ocupacional.
\item El 14 \% de los rechazos en portería corresponde a documentos vencidos que no se advirtieron con anticipación, y el trabajador rechazado ya viajó seiscientos kilómetros para llegar a la faena.
\item Las credenciales se prestan, situación que reconoce el jefe de seguridad, y el mecanismo que él propone (biometría) fue objetado formalmente por el sindicato.
\item Los incidentes, cuasi accidentes, observaciones de conducta e inspecciones se registran en formularios de papel que alguien digita días después. Ante una fiscalización, los registros deben buscarse en carpetas físicas.
\item Las alertas de fatiga y somnolencia llegan a un correo y a un teléfono de turno, pero lo que se hizo con cada una no queda registrado en ninguna parte.
\item Los reportes a la autoridad ambiental se arman a mano con datos que llegan por correo desde cinco áreas y desde laboratorios externos en PDF, lo que ya generó una observación por inconsistencia.
\item El consumo de agua se mide con un solo medidor general y se reparte entre áreas con una regla de 2019, y el combustible se conoce en litros comprados y no en litros consumidos por área.
\item La huella de carbono se calcula una vez al año con una consultora y el proceso toma dos meses. Los compradores la exigirán trimestral y auditada a partir de 2029.
\item Los indicadores del convenio con la comunidad, de empleo local y compras locales, se construyen cada semestre pidiendo planillas a recursos humanos y a abastecimiento.
\item Datos sensibles, como los resultados de exámenes, de control de fatiga y de alcohol y drogas, circulan en planillas y correos sin un control de quién los consulta.
\end{itemize}

\seccion{Indicadores del problema}
La Tabla \ref{tab:indicadores-problema-comparados-referencia-industria} reúne todos los indicadores del problema, agrupados según los tres problemas, junto con la referencia de la industria o la meta interna de la compañía con que se contrastan. El cuerpo del Subdocumento 2 presenta una selección de estas filas.

\begin{tabla}{I{0.374}D{0.200}D{0.206}I{0.220}}{Indicadores del problema comparados con la referencia de la industria o la meta interna de la compañía.}{tab:indicadores-problema-comparados-referencia-industria}
  {\cab{Indicador} & \cab{Valor Aranda 2025} & \cab{Referencia o exigencia} & \cab{Brecha respecto de la referencia}}
  {Fuente: elaboración propia a partir de las Bases técnicas del caso Minería (CMAA, 2026b, cap. 7, pp. 11–12).}
\multicolumn{4}{I{1.000}}{Problema 1: Trazabilidad comercial y conciliación} \\
Diferencia de conciliación mina-planta-embarque & 3,1 \% a 6,8 \% mensual & bajo 2 \% (referencia de industria) & 1,5 a 3,4 veces la referencia \\
Días hábiles para cerrar la conciliación del mes & 18 & 3 a 5 (referencia de industria) & 3,6 a 6 veces más lento \\
Tiempo para reconstruir el origen de un lote & 4 a 6 semanas, estimado & horas, verificable (referencia de industria) & El resultado no sirve como evidencia \\
Embarques con penalización por arsénico & 4 de 31 & no informada & 12,9 \% de los embarques \\
Verificación del supuesto de ley en retoma de stocks & nunca realizada & no informada & sin evidencia \\
\multicolumn{4}{I{1.000}}{Problema 2: Productividad de la operación} \\
Utilización efectiva de la flota & 71 \% & 80 \% (meta interna) & 9 puntos bajo la meta \\
Demoras operacionales sin clasificar & 23 \% del tiempo & bajo 5 \% (meta interna) & 4,6 veces la meta \\
Horas de detención no programada de planta & 214 h al año & bajo 120 h (meta interna) & 94 horas por sobre la meta \\
Detenciones sin causa raíz documentada & 61 \% & bajo 10 \% (meta interna) & 6 veces la meta \\
Tiempo entre falla y diagnóstico en equipo mina & 4,2 h & 1 h (meta interna) & 3,2 horas más por falla \\
Pérdida efectiva por cambio de turno & 47 min por turno & 25 min (meta interna) & 22 minutos de más en cada turno \\
Desfase en la medición del cumplimiento del plan & 10 a 15 días & turno siguiente (meta interna) & Se replanifica con datos antiguos \\
\multicolumn{4}{I{1.000}}{Problema 3: Cumplimiento regulatorio y gestión de personas} \\
Días para acreditar a un contratista nuevo & 9 & menos de 2, meta del mandante & 4,5 veces la meta del mandante \\
Rechazos en portería por documentos vencidos no advertidos & 14 \% de los rechazos & no informada & Rechazos evitables \\
Observaciones en la última auditoría de la autoridad minera & 11 (6 por trazabilidad) & no informada & 55 \% ligadas a registros \\
Registros de seguridad en papel digitados después & prácticamente todos & no informada & Sin valor inmediato ante fiscalización \\
Alertas de fatiga con acción documentada & sin registro & no informada & Exposición ante un accidente \\
Frecuencia de cálculo de la huella de carbono & anual, con consultora & trimestral y auditada (2029) & Cuatro veces la frecuencia actual \\
\end{tabla}

\seccion{Volumen de hechos operacionales y forma de registro}
Otra manera de dimensionar el problema es contar cuántas veces al año ocurre un hecho que debería quedar registrado y cómo se registra hoy. La Tabla \ref{tab:volumen-hechos-operacionales-forma-se-registran} detalla esos hechos a partir de la volumetría entregada por la compañía.

\begin{tabla}{I{0.288}D{0.235}I{0.476}}{Volumen de hechos operacionales y forma en que se registran hoy.}{tab:volumen-hechos-operacionales-forma-se-registran}
  {\cab{Hecho operacional} & \cab{Volumen anual actual (proyección a 3 años)} & \cab{Forma de registro actual}}
  {\emph{Fuente: elaboración propia a partir de las Bases técnicas del caso Minería.}}
Viajes de camión de extracción & 342.000 (367.000) & Sistema de despacho, sin registro de la razón de los cambios de destino \\
Ciclos de carguío de pala & 1,37 millones (1,47) & Sistema de despacho. El cambio de material se avisa por radio \\
Pozos de perforación & 46.000 (49.000) & Automático en 2 de 9 perforadoras y el resto en planilla \\
Muestras de laboratorio & 71.000 (78.000) & Sistema de laboratorio, luego se transcriben a planillas \\
Viajes de concentrado a puerto & 12.400 (13.000) & Guía de despacho y planilla de pesaje en papel \\
Órdenes de trabajo de mantenimiento & 34.000 (36.000) & SAP, aunque los avisos correctivos llegan por radio \\
Incidentes, cuasi accidentes y observaciones & 11.500 (12.000) & Formulario de papel digitado días después \\
Inspecciones planificadas de seguridad & 7.800 (8.200) & Formulario de papel digitado días después \\
Alertas de fatiga y somnolencia & 4.600 (5.000) & Correo y teléfono de turno, sin registro de la acción tomada \\
Ingresos y salidas en portería & 1,06 millones (1,13) & Torniquete y credencial, con la habilitación gestionada en una planilla \\
Personas con acceso a faena & 4.180 (4.400) y hasta 8.500 en detención & Planilla de acreditación mantenida a mano \\
Equipos con telemetría de fabricante & 52 (70) & Tres portales web separados \\
Instrumentos ambientales y geotécnicos & 186 (230) & Instrumentos propios, laboratorios externos y correo \\
Señales del historiador de planta & 14.500 etiquetas (16.000) & Consulta manual cuando se requieren \\
\end{tabla}

\seccion{Actores internos de la compañía}
Los actores internos que muestra la Figura \refext{S2}{fig:actores-internos-operacion-gobierno-corporativo} del cuerpo del Subdocumento 2 son los siguientes:

\begin{itemize}
\item Directorio y gerencia general: máxima instancia de decisión de la compañía, responsable de su dirección estratégica y de adjudicar la licitación.
\item Gerencia de finanzas: área responsable de la contabilidad, el cierre mensual y la conciliación de la compañía, ubicada en la oficina corporativa de Santiago.
\item Gerencia de operaciones mina: área responsable de la extracción en el rajo, es decir, de la perforación, la tronadura, el carguío y el transporte del mineral y el estéril, y de la productividad de la flota.
\item Superintendencia de planta y metalurgia: área a cargo del procesamiento del mineral en la planta concentradora y del balance metalúrgico.
\item Superintendencia de planificación minera: área que elabora el plan minero a partir del modelo de bloques, el diseño y el secuenciamiento de la mina.
\item Jefatura de mantenimiento: área responsable del mantenimiento de los camiones, las palas, las perforadoras y los equipos de apoyo de la faena.
\item Metalurgistas y laboratorio químico: profesionales que generan y procesan los datos de ley y recuperación del mineral: el laboratorio analiza las muestras y los metalurgistas calculan el balance.
\item Personal del Centro Integrado de Operaciones (CIO): equipo de 115 personas que opera en turnos rotativos la sala de control integrada de mina y planta, ubicada en Antofagasta.
\item Supervisores y despachadores: personal que articula la operación en tiempo real: los supervisores coordinan el trabajo en terreno y los despachadores asignan y reasignan los destinos de los camiones.
\item Operadores de equipos (palas, camiones de extracción, perforadoras): trabajadores que manejan los equipos en el rajo y que generan el dato de origen de la operación.
\item Sindicato: organización que representa a los trabajadores propios de la compañía.
\end{itemize}

\seccionsinnumero{Referencias}
\begin{referencias}
\item Compañía Minera Altos de Aranda. (2026a). \emph{Bases administrativas}. Licitación N° TFEP-01/2026.
\item Compañía Minera Altos de Aranda. (2026b). \emph{Bases técnicas del caso Minería}. Licitación N° TFEP-01/2026.
\item Compañía Minera Altos de Aranda. (2026c). \emph{Bases técnicas transversales}. Licitación N° TFEP-01/2026.
\item International Electrotechnical Commission. (s.f.). \emph{IEC 62443 background}. https://syc-se.iec.ch/?p=4375
\item Ley N° 20.417. (2010). \emph{Crea el Ministerio, el Servicio de Evaluación Ambiental y la Superintendencia del Medio Ambiente}. Diario Oficial de la República de Chile.
\item Ley N° 20.551. (2011). \emph{Regula el cierre de faenas e instalaciones mineras}. Diario Oficial de la República de Chile.
\item Ley N° 21.719. (2024). \emph{Regula la protección y el tratamiento de los datos personales y crea la Agencia de Protección de Datos Personales}. Diario Oficial de la República de Chile.
\item Ministerio de Minería. (2004). \emph{Decreto Supremo N.º 132, Reglamento de Seguridad Minera}. Diario Oficial de la República de Chile, 7 de febrero de 2004.
\end{referencias}
\end{document}
